%%%%%%%%%%%%%%%%%%%%%%%%%%%%%%%%%%%%%%%%%%%%%%%%%%%%%%%%%%%%%%%%%%%%%%%%%%%%%%%%%%%%%%%%%%%%%%%%%%%%%%%%%%%%%%%%%%%%%%%%%%%%%%%%%%%%%%%%%%%%%%%%%%
%%%%%%%%%%%%%%%%%%%%%%%%%%%%%%%%%%%%%%%%%%%%%%%%%%%%%%%%%%%%%%%%%%%%%%%%%%%%%%%%%%%%%%%%%%%%%%%%%%%%%%%%%%%%%%%%%%%%%%%%%%%%%%%%%%%%%%%%%%%%%%%%%%

\documentclass[a4paper, 10pt, oneside]{elegantbook}
\usepackage{geometry}
\usepackage{longtable}
\usepackage{booktabs}
\usepackage{calc}
\usepackage{amsmath}
\usepackage{imakeidx}
\usepackage{hyperref}
% Recommended for modern LaTeX
\makeindex[intoc, title=Index, options={-s \jobname.ist}]         % Mandatory command to initialize indexing
\usepackage[T1]{fontenc}
\usepackage{graphicx}
\usepackage{pdfpages} % For inserting your cover photo as a PDF
\usepackage{makecell}
\usepackage{tabularx}
\usepackage{fancyhdr}
\usepackage[utf8]{inputenc}
%\inputpackage{csquotes}
%\allowdisplaybreaks
\usepackage{amssymb}
 \usepackage{afterpage}
%\usepackage[pagestyles]{titlesec}
\hypersetup{colorlinks} % Comment this line if you don't wish to have colored links
\pagestyle{fancy}
\fancypagestyle{plain}{
  \fancyhf{}
  \fancyfoot[C]{\hyperlink{toc}{\thepage}}
}
\usepackage{mathtools}
\usepackage{listings}
\usepackage{ifthen}
\usepackage{perpage}
\usepackage{wrapfig}
\DeclarePairedDelimiter\abs{\lvert}{\rvert}
 \usepackage{tikz}
\usepackage{pgfplots}
\pgfplotsset{compat=1.18}

%\IfFileExists{indentfirst.sty}{%
%  \usepackage{indentfirst}
%}{}

%\usepackage{calc]

% for running heads
\usepackage{fancyhdr}
%\usepackage{xcolor}
\usepackage[dvipsnames]{xcolor} % Accesses a wide range of colors
\usepackage{booktabs} % For professional horizontal lines (\toprule, \midrule, \bottomrule)
\usepackage{enumitem} % For tight list spacing using "nosep"
\usepackage{tabularx} % For auto-wrapping text columns to text width
%\usepackage{array} % Required for enhanced table controls
%\usepackage{pdfstartpage=1}
\usepackage{ragged2e,array,booktabs}

%%%%%%%%%%%%%%%%%%%%%%%%%%%%%%%%%%%%%%%%%%%%%%%%%%%%%%%%%%%%%%%%%
%%%% Interlude:  Set up PRINTING (default) or SCREEN VIEWING %%%%
%%%%%%%%%%%%%%%%%%%%%%%%%%%%%%%%%%%%%%%%%%%%%%%%%%%%%%%%%%%%%%%%%

% ForPrinting=true (default)           false
% Asymmetric margins                   Symmetric margins
% Black hyperlinks                     Blue hyperlinks
% Start Preface, ToC, etc. recto       No blank verso pages
%
% Chapter-like ``Sections'' start both recto and verso in the scanned
% book. This behavior has been retained.
%\cleardoublepage
\newboolean{ForPrinting}

%% UNCOMMENT the next line for a PRINT-OPTIMIZED VERSION of the text %%
%\setboolean{ForPrinting}{true}

%% Initialize values to ForPrinting=false
\newcommand{\Margins}{hmarginratio=1:1}     % Symmetric margins
\newcommand{\HLinkColor}{blue}              % Hyperlink color
\newcommand{\PDFPageLayout}{SinglePage}
\bigskip

% Re-crop screen-formatted version, accommodating wide displays
%\ifthenelse{\boolean{ForPrinting}}
%  {}
%  {\hypersetup{pdfpagescrop= 20 0 412 580}}


%%newcounter{mytnote}
%newcommand{\TNote}[1]{%
%  \refstepcounter{mytnote}%
%\title{Quantum Mechanics(reference guide)}


%\subtitle{With magnetic deflection, notes(Theory, barely descriptive)}
%\author{Owino Stephen}

%\institute{\copyright \the\year{} This book was typeset using \LaTeX{} Program}

%\date{May, 18, 2026}

%\version{1.0}

%\bioinfo{Bio}{https://github.com/StevenOwino}

%\extrainfo{The principles of reality are atoms, and emptiness, while other things are mere opinion.
%(Democritus)}

%\afterpage{\clearpage}


%\logo{infinity.png}

%cover{infinity.png}

\definecolor{customcolor}{RGB}{32,178,170}
\colorlet{coverlinecolor}{customcolor}

% Patch hyperref to ensure index page numbers remain clickable links
%\newcommand*{\hyperpage}[1]{\hyper
\begin{filecontents}[force]{\jobname.ist}
headings_flag 1
heading_prefix "\\par\\penalty-50\\textbf{"
heading_suffix "}\\par\\nobreak"
\end{filecontents}

%\newboolean{ForPrinting}

%\newcommand{\DPPageSep}[2]{\Pagelabel{#2}}

%\def\ebook{43}

%% Macro for reducing space between adjacent display environments
\newcommand{\CancelMathSkip}{\vspace{-\belowdisplayskip}}
\newcommand{\BindMath}[1]{%
{\vspace{\abovedisplayskip}%
{\setlength{\abovedisplayskip}{0pt}%
\setlength{\belowdisplayskip}{0pt}%
#1}\vspace{\belowdisplayskip}}%
}

%\newcommand{\ToCBox}[1]{%
%  \PadTxt[r]{XVIII.}{#1}\quad%
%}

%\newcommand{\ToCAnchor}{}

%\ToCLine[font style]{Chapter}{Title}{label}
\newcommand{\ToCLine}[4][\scshape]{%
  \par\noindent\label{toc:#4}%
  \ifthenelse{\not\equal{\pageref{toc:#4}}{\ToCAnchor}}{%
    \noindent\textsc{\scriptsize Chapter\hfill Page}\medskip%
    \renewcommand{\ToCAnchor}{\pageref{toc:#4}}\par%
  }{}%
  \settowidth{\TmpLen}{9999}%
  \noindent\parbox[b]{\linewidth-\TmpLen}{\ToCBox{#2}%
    \rule{0pt}{16pt}#1\hangindent6em #3\dotfill}% [** TN: Hard-coded spaces]
  \PadTxt[r]{9999}{\pageref{#4}}\smallskip%
}

\newcommand{\ChapterSkip}{\vspace*{0.1\linewidth}}
%\newcommand{\flushpage}{\clearpage\fancyhf{}\cleardoublepage}

% Numbered: \Chapter[Running Head]{number}{Title}
\newcommand{\Chapter}[3][]{%
  \flushpage
  \phantomsection\label{chap:#2}\pdfbookmark[0]{Chapter #2}{Chapter #2}%
  \ifthenelse{\not\equal{#1}{}}{%
    \SetOddHead{#1}%
  }{%
    \SetOddHead{#3}%
  }%

  \addtocontents{toc}{\protect\ToCLine{#2.}{#3}{chap:#2}}
  \thispagestyle{empty}
  \ChapterSkip%
  \section*{\centering\normalfont\Large CHAPTER #2.}
  \subsection*{\centering\normalfont\large\MakeUppercase{#3}.}
}

% Unnumbered, but has a ToC entry: \AltChapter[Running Head]{Title}
%[** TN: Uses the book's chapter structure to select ToC entry font]
\newcommand{\AltChapter}[2][]{%
  \ifthenelse{\equal{#2}{Prologue}}{%
    \pagestyle{fancy}
  }{}
  \flushpage
  \phantomsection\label{#2}\pdfbookmark[0]{#2}{#2}%
  \ifthenelse{\not\equal{#1}{}}{%
    \SetOddHead{#1}%
    \addtocontents{toc}{\protect\ToCLine{}{#1}{#2}}%
  }{%
    \SetOddHead{#2}%
    \ifthenelse{\equal{#2}{Table of Standard Forms}}{%
      \addtocontents{toc}{\protect\ToCLine[\bfseries]{}{#2}{#2}}%
    }{%
      \addtocontents{toc}{\protect\ToCLine{}{#2}{#2}}%
      \ChapterSkip%
    }%
  }%

  %\thispagestyle{empty}
  \section*{\centering\normalfont\Large\MakeUppercase{#2.}}
}

% For the Preface and Epilogue; no ToC entry
\newcommand\ChapterStar[2][]{%
  \flushpage
  \phantomsection

  \ifthenelse{\not\equal{#1}{}}{%
    \SetOddHead{#1}%
    \pdfbookmark[0]{#1}{#1}%
  }{%
    \SetOddHead{#2}%
    \pdfbookmark[0]{#2}{#2}%
  }%
%  \thispagestyle{empty}
  \ChapterSkip%
  \section*{\centering\normalfont\Large\MakeUppercase{#2.}}
}

%\newcounter{SectionNo}
%\setcounter{SectionNo}{0}

\newcommand\Section[2][]{%
  \section*{\centering\normalfont\normalsize\bfseries#2}%
  \refstepcounter{SectionNo}%
  \phantomsection\pdfbookmark[1]{#2}{#2}\label{section:\theSectionNo}%
  \ifthenelse{\not\equal{#1}{}}{\SetOddHead{#1}}{}%
}

%tweaking \subsection - added italics & indent
\newcommand\Subsection[1]{%
  \medskip\pagebreak[1]\par\textit{#1}\pagebreak[0]\par%
}

%tweaking \paragraph - added italics & indent
\newcommand{\Paragraph}[1]{\medskip\pagebreak[1]\par\textit{#1}}

\newcommand{\Note}[1]{%
  %\clearpage
  \section*{\centering\normalfont\normalsize #1}
}

% Dedicated structural macros
\newcommand{\Case}[1]{\Subsection{Case~\upshape{#1}.}}

% List item formatting
\newcommand{\ListInit}{%
  \setlength{\leftmargin}{0pt}%
  \setlength{\labelwidth}{\parindent}%
  \setlength{\labelsep}{0.5em}%
  \setlength{\itemsep}{0pt}%
  \setlength{\listparindent}{\parindent}
  \setlength{\itemindent}{2\parindent}%
}

\newcommand{\License}{%
%  \flushpage
%  \thispagestyle{empty}
%   \phantomsection\pdfbookmark[0]{License}{License}
  \SetBothHeads{License}
}

%\newcommand{\Z}{\phantom{0}}
\newcommand{\DStrut}{\rule[-12pt]{0pt}{32pt}}
\newcommand{\Strut}{\rule{0pt}{16pt}}

\newcommand{\First}[1]{\noindent\textsc{#1}}

% "exponent fraction" factored out for special handling, if desired
\newcommand{\efrac}[2]{\frac{#1}{#2}}

\DeclareMathOperator{\Arccos}{arc\,cos}
\renewcommand{\arccos}{\Arccos}

\DeclareMathOperator{\Arcsin}{arc\,sin}
\renewcommand{\arcsin}{\Arcsin}

\DeclareMathOperator{\Arctan}{arc\,tan}
\renewcommand{\arctan}{\Arctan}

\DeclareMathOperator{\arcsec}{arc\,sec}
\DeclareMathOperator{\cosec}{cosec}
\DeclareMathOperator{\cotan}{cotan}
\DeclareMathOperator{\sech}{sech}

% DPalign
\makeatletter
\providecommand\shortintertext\intertext
\newcount\DP@lign@no
\newtoks\DP@lignb@dy
\newif\ifDP@cr
\newif\ifbr@ce
\def\f@@zl@bar{\null}
\def\addto@DPbody#1{\global\DP@lignb@dy\@xp{\the\DP@lignb@dy#1}}
\def\parseb@dy#1{\ifx\f@@zl@bar#1\f@@zl@bar
    \addto@DPbody{{}}\let\@next\parseb@dy
  \else\ifx\end#1
    \let\@next\process@DPb@dy
    \ifDP@cr\else\addto@DPbody{\DPh@@kr&\DP@rint}\@xp\addto@DPbody\@xp{\@xp{\the\DP@lign@no}&}\fi
    \addto@DPbody{\end}
  \else\ifx\intertext#1
    \def\@next{\eat@command0}%
  \else\ifx\shortintertext#1
    \def\@next{\eat@command1}%
  \else\ifDP@cr\addto@DPbody{&\DP@lint}\@xp\addto@DPbody\@xp{\@xp{\the\DP@lign@no}&\DPh@@kl}
          \DP@crfalse\fi
    \ifx\begin#1\def\begin@stack{b}
      \let\@next\eat@environment
  \else\ifx\lintertext#1
    \let\@next\linter@text
  \else\ifx\rintertext#1
    \let\@next\rinter@text
  \else\ifx\\#1
    \addto@DPbody{\DPh@@kr&\DP@rint}\@xp\addto@DPbody\@xp{\@xp{\the\DP@lign@no}&\\}\DP@crtrue
    \global\advance\DP@lign@no\@ne
    \let\@next\parse@cr
  \else\check@braces#1!Q!Q!Q!\ifbr@ce\addto@DPbody{{#1}}\else
    \addto@DPbody{#1}\fi
    \let\@next\parseb@dy
  \fi\fi\fi\fi\fi\fi\fi\fi\@next}
\def\process@DPb@dy{\let\lintertext\@gobble\let\rintertext\@gobble
  \@xp\start@align\@xp\tw@\@xp\st@rredtrue\@xp\m@ne\the\DP@lignb@dy}
\def\linter@text#1{\@xp\DPlint\@xp{\the\DP@lign@no}{#1}\parseb@dy}
\def\rinter@text#1{\@xp\DPrint\@xp{\the\DP@lign@no}{#1}\parseb@dy}
\def\DPlint#1#2{\@xp\def\csname DP@lint:#1\endcsname{\text{#2}}}
\def\DPrint#1#2{\@xp\def\csname DP@rint:#1\endcsname{\text{#2}}}
\def\DP@lint#1{\ifbalancedlrint\@xp\ifx\csname
DP@lint:#1\endcsname\relax\phantom
  {\csname DP@rint:#1\endcsname}\else\csname DP@lint:#1\endcsname\fi
  \else\csname DP@lint:#1\endcsname\fi}
\def\DP@rint#1{\ifbalancedlrint\@xp\ifx\csname
DP@rint:#1\endcsname\relax\phantom
  {\csname DP@lint:#1\endcsname}\else\csname DP@rint:#1\endcsname\fi
  \else\csname DP@rint:#1\endcsname\fi}
\def\eat@command#1#2{\ifcase#1\addto@DPbody{\intertext{#2}}\or
  \addto@DPbody{\shortintertext{#2}}\fi\DP@crtrue
  \global\advance\DP@lign@no\@ne\parseb@dy}
\def\parse@cr{\new@ifnextchar*{\parse@crst}{\parse@crst{}}}
\def\parse@crst#1{\addto@DPbody{#1}\new@ifnextchar[{\parse@crb}{\parseb@dy}}
\def\parse@crb[#1]{\addto@DPbody{[#1]}\parseb@dy}
\def\check@braces#1#2!Q!Q!Q!{\def\dp@lignt@stm@cro{#2}\ifx
  \empty\dp@lignt@stm@cro\br@cefalse\else\br@cetrue\fi}
\def\eat@environment#1{\addto@DPbody{\begin{#1}}\begingroup
  \def\@currenvir{#1}\let\@next\digest@env\@next}
\def\digest@env#1\end#2{%
  \edef\begin@stack{\push@begins#1\begin\end \@xp\@gobble\begin@stack}%
  \ifx\@empty\begin@stack
    \@checkend{#2}
    \endgroup\let\@next\parseb@dy\fi
    \addto@DPbody{#1\end{#2}}
    \@next}
\def\lintertext{lint}\def\rintertext{rint}
\newif\ifbalancedlrint
\let\DPh@@kl\empty\let\DPh@@kr\empty
\def\DPg@therl{&\omit\hfil$\displaystyle}
\def\DPg@therr{$\hfil}

\newenvironment{DPalign*}[1][a]{%
  \if m#1\balancedlrintfalse\else\balancedlrinttrue\fi
  \global\DP@lign@no\z@\DP@crfalse
  \DP@lignb@dy{&\DP@lint0&}\parseb@dy
}{%
  \endalign
}
\newenvironment{DPgather*}[1][a]{%
  \if m#1\balancedlrintfalse\else\balancedlrinttrue\fi
  \global\DP@lign@no\z@\DP@crfalse
  \let\DPh@@kl\DPg@therl
  \let\DPh@@kr\DPg@therr
  \DP@lignb@dy{&\DP@lint0&\DPh@@kl}\parseb@dy
}{%
  \endalign
}
%\makeatother

\pagenumbering{arabic}

\begin{document}

\begin{titlepage}
 % Set the cover background color (e.g., MidnightBlue)
    %\pagecolor{MidnightBlue}
% Change text color to white for readability
    %\color{white}
    \centering
    \noindent
    \includegraphics[width=4.2cm]{infinity.pdf}
    %\vspace*{\fill}
    %\vspace*{1cm} % Adjust to push text down from the top
    {\Huge \textbf{Pre-QFT Notes}}\\[2cm]
    {\Large OWINO STEVE}\\[0.5cm]
    %\vfill % Pushes the remaining text/publisher info to the bottom
    %{\large S\\ 2026}  %Publisher Name\\ Year
    %\clearpage
    This book was typeset using \LaTeX{}; The Photon(Massless Particle in QFT): U($\Lambda$)$\Psi_{p, \sigma}$ = $\sqrt{\frac{(\Lambda)P^{0}}{P^{0}}}$ $e^{i \sigma \theta(\Lambda, P)}$ $\Psi_{\Lambda P,\sigma}$
    \date{May, 18, 2026}
    %\version{1.0}
    \bioinfo{Bio}{https://github.com/StevenOwino}
    \extrainfo{Someone asks if it's measurable, and you say, yes it is. You’re going to feel pretty silly if they spend a million and a half dollars to do the measurement, and can’t find it, because you’ve put a $(2 \pi)^2$ in a numerator, when it should have been in the denominator... - Sidney Coleman}
\end{titlepage}
%\makecoverpage % This prints the cover with the equation


%\newpage

%\thispagestyle{empty} % Removes page number from cover
%\hypertarget{toc}{}

\setcounter{page}{1}
\maketitle
%%\afterpage{\clearpage}

%\thispagestyle{empty}
%\clearpage
\nopagecolor

\setcounter{tocdepth}{4}
\mbox{}
%\pagenumbering{arabic}

%\frontmatter
%\tableofcontents

% Table of Contents
% Redefine headings for table of contents & index
\renewcommand{\contentsname}{%
  \protect\Large\protect\normalfont\protect\centering%
  CONTENTS.\\
}

%\mainmatter
\section{Units}

C = 1, $\hbar$ = $1$  \\

C (Speed of Light in Vacuum) = $299792458 \quad m s^{-1}$  \\

$\hbar$ (Reduced Plancks' Constant) = $1.054571817$ \ldots $\times$ $10^{-34}$ J s  \\

$G$ (Newtonian Constant of Gravitation) = $6.674 30(15)$ $\times$ $10^{-11} m^{3} kg^{-1} s^{-2}$  \\

C, $\hbar$, and G, are the three fundamental Constants.





















\newpage

\section{Preface}


The \emph{cover} is a figure of the infinity symbol. These are my notes, NOT pedagogical. I shall try to make QFT penetrable. A lot of intellectual buggage isn't exhaustive in this material at all. \\

 I assume the \emph{"Principle of Locality"} in general,   where the action (S), in the equation of motion is:  \\

 $\int$ ${d^4}$(x) $\left[ - \frac{1}{\sum}\partial_\mu \phi \quad \partial^{\mu} \phi - V(\phi) \right]$; $\forall$ parameters in V($\phi$) constant.  \\

The framework of attempting to comprehend nature is on Energy scales \ldots. I do not give References/Citations of original papers, and the Bibliography given, makes no pretence of being exhaustive.

Abstract Mathematics(models of inductive reasoning), is availed to avoid non-sensical answers, and bogus non-differentiable functions. This is the task for the human mind; to see this truth, and preserve it.  \\
It's accepted that the Angular Momentum of a closed system, is conserved(hence production of particles by conversion of Gravitational Energy, into matter Energy, conditioned by Biological factors).  \\
The origin of the Universe(a Physical picture of a polarizable space, and characterized as young, by $\frac{\epsilon} {\epsilon_0}$ $\approx$ 1, with an Energy density when expressed in \emph{Newtonian Units}), is not a reproducible event, and cannot be duplicated in the laboratory. In trying to eliminate the Riemann curvature tensor, there's also the Ricci curvature tensor, assuming the Photon travels without a shift in Energy or Frequency(red shift).  \\
String Theory is a highly sophisticated Theoretical laboratory(rather than experimental), it's defined at the Planck Energy($10^{19}$ GeV), as any Physical theory that you can possibly come up with, it cannot be tested with the present technology, and it's tightly constrained, as much as the Graviton emerges as a massless state, with Spin-2.  \\
Space-time is beyond being verified experimentally, and of the vacua states, only one is needed to be isolated, and the  Cosmological constant is very close to 0. However, in this theory, when there's spontaneous symmetry breaking, one does not know how to keep the Cosmological constant 0, perhaps a solution will be found to isolate the one vacuum.

There's no Space/Time for boredom! Have fun, and ask a Physics lab fundamental questions.

The letters in bold transform like vectors!, and the entire notes parses on TeX Live 2023/Debian kpathsea, version 6.3.5.

These are notes/references, contains a few functional code!, and compiling.


\First{An} electron plays a central role, and can be assumed to move in a spherically symmetric field, where all binding energies are negative. These are advanced core notes(reference sheet) in Quantum Mechanics(describing how particles behave like waves -
the state of a given system[internal structure], and its' probabilities), 
and some equations of Mathematical Physics[some theorems are mentioned to emphasize the concept of unitarity].
The key important formulas are the Schr$\ddot{o}$dinger equation, 
Uncertainty Principle, energy levels of the hydrogen atom, harmonic oscillator energy,
and the Hamiltonian energy(of which describes how a system evolves in time). 
Included, are physical quantities, long expressions, Gauge Theory, 
important systems(like the Harmonic Oscillator, the energy levels of the Hydrogen atom), Mathematical tools, 
important inequalities, and some advanced notes on Quantum Field Theory. The multiplication by $i$ makes it possible
to carry over to QM(Quantum Mechanics)...this is one of the most striking features of Classical Mechanics(CM)-namely,
the correspondence between observables, and one-parameter groups of symmetries. We hereby have a Mathematical model
capable in principle, of supplying the answer to any Physical question we might ask. The fact that the eigenvalues check with
the observed spectrum of the hydrogen atom, when one takes $\mathbf K$ = $\frac{\hbar}{2 \pi}$, which provides a confirmation
of the identity of $\hbar$, with $\frac{\hbar}{2 \pi}$. Your objective, dear reader, is to show in theorems of analysis, providing
asymptotic formulas for the eigenvalue distribution of differential operators.                                                                      
Don't try to read these notes linearly, instead, come back here as a formula reference. \\

\section{Why Quantum Mechanics?}

A complete description of the state of a Physical System, in Classical Mechanics, is effected by stating all its Coordinates, and Velocities at a given instant; With these initial data, the \emph{Equations Of Motion} completely determine the behaviour of the system at all subsequent instants. \\
In Quantum Mechanics, such a description is in principle impossible, since the Coordinates, and the corresponding Velocities, cannot exist simultaneously.
Thus a description of the state of a Quantum System is effected by means of a smaller number of quantities, than in Classical Mechanics, i.e. it is less detailed than a Classical description.  \\
A very important consequence follows from this, regarding the nature of the \textbf{predictions}, made in Quantum Mechanics. Whereas a Classical description suffices to predict the future motion of a mechanical system with complete accuracy, the less detailed description given in Quantum Mechanics evidently cannot be enough to do this. This means that, even if an \emph{Electron} is in a state described in the most complete manner possible in Quantum Mechanics, its' behaviour at subsequent instants, is still in principle \textbf{uncertain}. Hence Quantum Mechanics, cannot make completely definite predictions concerning the future behaviour of the Electron. For a given initial state of the Electron, a subsequent \textbf{measurement} can give various results. The problem in Quantum Mechanics, consists in determining the \textbf{probability} of obtaining various results on performing this \emph{measurement}. It is understood, of course, that in some cases, the \emph{probability} of a given result of \emph{measurement} may be equal to  \textbf{unity}, i.e. \emph{certainty}, so that the result of that \emph{measurement} is \emph{unique}.  \\
Completely described states, occur as a result of the \emph{simultaneous measurement} of a complete set of Physical quantities. From the results of such a \emph{measurement}, it's possible to determine the probability of various results of any subsequent \emph{measurement}, regardless of the history of the Electron, prior to the first \emph{measurement}.




















\newpage

\section{List of Figures in Chapters}

\begin{enumerate}

\item Binary Entropy Diagram
\item Vector pointing towards the Sun
\item Quantum Harmonic Oscillator
\item Anharmonic vs simple Harmonic Oscillator
\item Relativistic Mass increase
\item Ground state search
\item Sampled wave-function distribution
\item Breit-Wigner distribution
\item Stern-Gerlach projection

\end{enumerate}

\newpage

\section{Diagnostic}

%%%% PG BOILERPLATE %%%%
%$\phantomsection
%\pdfbookmark[0]{Boilerplate}{Boilerplate}
%\begin{center}
%\begin{minipage}{\textwidth}
%\small
%\section{Diagnostic}
%\begin{PGtext}

(a) Show that $y = 2x^2$ is a solution to the differential equation: $\frac{1}{2}x^2y''-xy'+y$ = 0, \\

(b) If $y_1(t)$ = $e^{3t}$, and $y_2(t)$ = $e^{-3t}$ are solutions to $y'' - 9y$ = 0, what is the general solution? Note that $y_1$ and  \\
$y_2$ are NOT constant multiples of one another, so they are linearly independent. Thus the general solution in the differential equation is  \\
y(t) = $c_1e^{3t} + c_2e^{-3t}$. \\

(c) If $y^{''} - 2y^{'} + 5y $= 0, has the associated characteristic equation $\lambda^2 - 2\lambda + 5$ = 0,  \\
by the quadratic formula, the roots of the characteristic equation are 1$\pm 2{i}$.
Therefore, the general solution to this differential equation is y(x) = $e^x(c_1cos2{x}+c_2sin2{x})$.  \\

(d) Knowing when to apply the $\mathbf{\textbf{Product and Chain Rules, Total and Partial Derivatives}}$.  \\

NOTE: $\frac{\hbar}{i}$ = -i$\hbar$ (Multiplying by -i is the same as dividing by i.)  \\

(e) $\int (a_0 x^4 + 4a_1 x^3 + 6a_2 x^2 + 4a_3 x + a_4)^{- \frac{1}{2}}$ $\rightarrow$  \\




PROCEED  \\

%\vspace{2\MySkip}

\begin{note}

Elementary logic functions(Boolean gates), or Disjunctive Normal Forms, Computation complexities, Algorithms, and Quantum circuits, are not part of this material.

$\bullet$ Schr$\ddot{o}$dinger Equation: General linear homogeneous solution(from translation operators), to the Time-Dependent Schr$\ddot{o}$dinger Equation $\longrightarrow$ \\

\begin{equation}
\psi(x,t) = \sum_{n = 1}^\infty C_n \psi(x) e^{- \frac{i E t}{\hbar}}, \mathbf{\textbf{\quad t = 0}}
\end{equation}

$\bullet$ Integral form, time-dependent $\longrightarrow$  \\
\begin{equation}
\psi(x,t) = \int \mathbf K[x,t:q_0,t_0]\psi(q_0,t_0)dx_0, \mathbf{\textbf{where K is the kernel or propagator}}.
\end{equation}

Schr$\ddot{o}$dinger Time Independent Equation  \\
\begin{equation}
\boxed{E \psi_0 (x) = - \frac{\hbar^2}{2m} \frac{\partial^2 \psi}{\partial x^2}  + V \psi_0 (x)}
\end{equation}


$\bullet$ The identity Zero: Zero has no multiplicative inverse!!  \\

$\bullet$ Informational Warning: All descriptions here are in the literature, and not original. DO NOT attempt to frame a physical picture!!, or take Relativity into sufficient account. Quantum Field Theory follows once these principles are mastered. \\

$\bullet$ On Rational numbers: If (a) $\neq$ 0, there is no rational number($\mathbb Q$), where (b) $\times$ 0 = (a), can be fulfilled.



\chapter{\textbf{Familiar, Recognized, And Understood Equations}}
\label{ch:Familiar, recognized, and understood equations}

 If $ \frac{d \alpha} {dt} $ = 0 $\forall$ t, then $\alpha$ is constant \\

 $(f,\phi,) = \lim_{\alpha \to 0} (f, \phi_n) = 0 $ \\

 $\frac{d \sigma} {d \Omega}$ = $|f(\theta, \phi)|^2$  \\

 $ e^{i \mathbf a \cdot \mathbf X} e^{1 \mathbf b \cdot \mathbf P} \phi_0$ is also a minimum uncertainty state  \\

 $a^\dagger |n \rangle$ = $\sqrt{n + 1} |n +1 \rangle$ \\

\begin{equation}
{\hat j_\pm} \mid j,m \rangle = \hbar \sqrt{j (j + 1) - m \mid m \pm 1)} \mid j, m \pm 1 \rangle
\end{equation} 

$\implies$ \\

Angular-momentum\index{Angular-momentum} \text{identity/Discrete orientation, with Quantum-number}\index{Quantum-number}. \quad (In units of $\frac{\hbar}{2\pi}$).  \\
  
A(t) = $e^{+ \frac{iHt}{h}} \mathbf {A}e^{- \frac{iHt}{h}}$  \\
  
$\hat{\mathbf{H}}$ = $p \hat q - \mathbf L$ \\
 
ct' = $ \gamma(ct - \beta x)$ \\
 
E = $ ( \mathbf {p^2} c^2 + m^2 c^2)^ \frac{1}{2}$  \\

$ \mathbf E$ = -$ \frac{\mathbf {\dot A}}{c} - \nabla \varphi$   \\
 
$ \forall F \in H L^2(\mathbb C,\mu_\hbar)$, we have:  \\
\begin{equation}
\Vert T_\mathbf a F \Vert^2 L^2(\mathbb C^n,\mu_\hbar)
= (\pi \hbar)^{-n} \int \mathbb C_n  e^{- \hbar \vert \mathbf{a} \vert^2}  \mathbf{e}^{-2 \mathbb{R} (\bar{\mathbf{a}} \cdot {\mathbf{a}})} \vert \mathbf F (\mathbf z + \hbar \mathbf a ) \vert^2  \mathbf e^{- \frac{\vert \mathbf x \vert^2} {\hbar}}  dx
= (\pi \hbar)^{-n} \int \mathbb C_n e^{- \frac{\vert \mathbf z + \hbar \vert^2}{\hbar}} \vert \mathbf F(\mathbf z +\hbar \mathbf a) \vert^2 dz
\end{equation}
\noindent
= $\Vert \mathbf{F} \Vert^2$ $L^2(\mathbb{C^n}$,$\mu \hbar)$  \\
  
$\prod (e^X)$ = $e^{\pi (X)}$ = $ \sum_{x=0}^\infty \frac{\pi X^m}{m!}$ (Each element is a product of Lie algebra elements. Its power series is finite dimensional, and closed.)  \\
                                                                  
The Hilbert-Schmidt-norm\index{Hilbert-Schmidt-norm} of a matrix X $\in$ $M_n(\mathbb C)$ $\longrightarrow$  \\
\begin{equation} 
\Vert X \Vert^2_{HS} = \sum_{j,k = 1}^{n} |X_{jk}|^2  
\end{equation}
  
Cauchy-Schwarz inequality\index{Cauchy-Schwarz inequality} 
\begin{equation}
\Vert {XY} \Vert_{HS} \leq \Vert X \Vert_{HS} \Vert Y \Vert_{HS} \forall X,Y \in M_n(\mathbb C) 
\end{equation}

SO(3)\index{SO(3)} $\longrightarrow$
\begin{equation} 
\sigma_{\ell} = \pi'_{\ell} \circ \phi^{-1} = \pi_\ell
\end{equation}

$\hat \otimes \longrightarrow$ Hilbert tensor product \\

A complete bounded norm-preserving orthonormal vector, and closed unitary bilinear mapping,  \\
such that, $\psi$: $\mathscr H$ $\times$ $\mathscr H$ $\rightarrow$ G.  \\

If U $\in$ $\mathscr L (\mathscr H)$ and V $\in$ $\mathscr L (\mathscr K)$ are unitary, then so is U $\otimes$ V .  \\

U $\otimes$ V $(U \otimes V )^*$ = $I$ (sesquilinear or conjugate linear mapping)  \\

If s $\rightarrow$ $U_s$ $\in$ $\mathscr L (\mathscr H)$, and t $\rightarrow$ $V_t$ $\in$ $\mathscr L (\mathscr H)$ are strongly continuous unitary representations of topological groups G and H , respectively, then (s, t) $\rightarrow$ $U_s$ $\otimes$ $V_t$ $\in$ $\mathscr L (\mathscr H \otimes \mathscr H')$ is a strongly continuous unitary representation of the product topological group G $\times$ H.  \\

$\mathbf{\textbf{Proof}}$. If (s, t) $\rightarrow$ (s0, t0) (i.e., s $\rightarrow$  $s_0$, and t $\rightarrow$ $t_0$), then ($U_s$ $\otimes$ $V_t$) (x) $\otimes$ (y) = $U_s$(x) $\otimes$ $V_t$(y) $\rightarrow$ $U_s0$(x) $\otimes$ $V_t0$(y) = ($U_s0$ $\otimes$ $V_t0$)  $\otimes$ y . Thus (s, t) $\rightarrow$ $U_s$ $\otimes$ $V_t$ is strongly continuous on H $\odot$ H' (with the set of all finite sums($\sum_{1}^n$ ($x_k$ $\otimes$ $y_k$).  \\

Feynmann Path-Integral\index{Path-Integral} $\longrightarrow$  \\
\begin{equation}
e^{- \frac{i t \hat {H}}{\hbar}} \psi (X_0) = C \int_{\text{paths with} \\ X(0)=X_0} exp\left\{ \frac{i}{\hbar}(x(\cdot),0,t) \right\} \psi(x(t)) Dx
\end{equation}
 
Gauge-transformation\index{Gauge-transformation} $\longrightarrow$ \\
\begin{equation} \label{eq:Gauge-transformation}
U_\gamma Q^{1}_{pre}(f)U^{-1}_{\gamma} = Q^{2}{_{pre}}(f), \quad where (U_\gamma) \text{is the map, and "pre" is for pre-quantization\index{pre-quantization}}.
\end{equation}

$a^{\prime}$ = $a - \nabla \psi$ $\quad$  $\phi^{\prime}$ = $\phi$ + $\frac{1}{c}$ $\frac{\partial \psi}{\partial t}$

\boxed{\text{This was David Bohms' view of the character of Electromagnetic potentials.}} 

\begin{equation}
\psi_n(r,\phi) = f(r) e^{-i n\phi},\quad \text{where n is an integer, and f is an arbitrary holomorphic function\index{holomorphic function} on} \mathbb (C), with
\int_0^{\infty}\mid f(r)\mid^2r dr < \infty
\end{equation}
 
Sympletic-potential\index{Sympletic-potential} $\longrightarrow$ $\theta$ = $\frac{1}{2}$(p dx - x dp) = $\frac{1}{{2m} \omega}$(p dy - y dp)  \\
  
Fourier-transform($\mathcal{F})$\index{Fourier-transform} $\longrightarrow$ A constant multiple of a unitary map $\longrightarrow$ $\psi(p)$ = - $\int_\mathbb R$ $\phi$(x)$e^{- \frac{{ixp}}{\hbar}}$ dx dp  \\
  
$C_c^{\infty}(\mathbb {R^n})$ $\longrightarrow$ Space of smooth, compactly supported functions\index{Space of smooth, compactly supported functions} \\
\end{note}

\section{\textbf{Common nomenclacture}}
\paragraph{Nomenclecture}
\begin{enumerate}
    \item Written in polar coordinates, $\rho$ and $\theta$ are in the unit circle; When $\rho < 1$, we have the series u = $ \sum_{\nu=1}^\infty a_\nu \rho^{\nu} \cos \nu \theta$  
    \item $\pm$ means $-1^{|r|}$ 
    \item $\nabla$ - nabla Differential operator in ordinary 3D space
    \item $\nabla$h - Gradient of scalar h 
    \item $\nabla$ $\cdot$ $\mathbf A$ - Divergence of a vector 
    \item $\nabla^2$ - Laplacian operator in ordinary 3D space spacetime $\longrightarrow$ $\frac{\partial^2}{\partial x^2} + \frac{\partial^2}{\partial y^2} + \frac{\partial^2}{\partial z^2}$ 
    \item $\Box$ - nabla operator in 3D Minskowski spacetime 
    \item $\Box^2$ - d'Alembertian operator in Minskowski 4D spacetime 
    \item $\overline{V}$ = ($V^\perp)^\perp$ is the closure of V 
    \item $g_{i,j}, g^{i,j}$ - Covariant, contraviant metric tensor of 3D space, or its components
    \item $\epsilon_0$ - Permitivity\index{Permitivity} of free space 
    \item $\mu_0$ - Permeability\index{Permeability} of free space
    \item $\rho, \phi$, z - Cylindrical corodinates in 3D space 
    \item $\sigma$ - Spacetime interval in 4D 
    \item $\phi$ - Electromagnetic scalar potential 
    \item $\omega$ - Angular-frequency\index{Angular-frequency}/speed 

\end{enumerate}

\paragraph{Structure Of The Book}

\begin{lstlisting}
The content of the book is based on spaced structuring, and iterative deepening.
\end{lstlisting}

I would like to appreciate \LaTeX{} ,  and


\includegraphics[width=2cm]{overleaf_wide_colour.pdf}, for making ""learning coding" possible, for industry standards.

\chapter{\textbf{History Of Quantum Mechanics}}

General relativity was found in 1915. \emph{Quantum mechanics in 1926}. In 1916, Einstein pointed out, that Quantum effects must lead to modifications in the Theory of General Relativity, since the Gravitational Field behaved Quantum mechanically. In 1927, Oskar Klein suggested that Quantum Gravity should modify the concepts of Space and Time.  \\


In the early thirties, Rosenfeld, wrote the first technical papers on Quantum Gravity, applying Paulis' method for the quantization of fields with Gauge Groups, to the \emph{linearized Einstein Field Equations}. The relation with a linear spin-two Quantum Field got unraveled in the works of Fierz and Pauli.
and the spin-two quantum of the Gravitational Field, first named ``\textbf{Graviton}'' in a 1934 paper by Blokhintsev, and Gal'perin. Bohr then considers the idea of identifying the \emph{Neutrino}, and the \emph{Graviton}.  \\
In 1938, Heisenberg pointed out that \emph{the Gravitational Coupling-Constant is dimensional}, and is likely to cause problems with the Quantum Theory of the Gravitational Field.  \\
Matvei Petrovich Bronstein, re-derives the Rosenfeld-Pauli quantization of the linear theory; the unique features of Gravitation require a special treatment, when the full nonlinear theory is taken into account. Field quantization techniques, must be generalized in such a way as to be applicable in the absence of a background Geometry. The limitation posed by General Relativity on \emph{mass density}, radically distinguishes the Theory from Quantum
Electrodynamics, and would ultimate lead to the need to ``reject Riemannian geometry" and perhaps also to ```reject our ordinary concepts of Space and Time".

\vskip2mm{\bf 1949}\nopagebreak

Peter Bergmann starts the program of phase space quantization of non-linear Field Theories; Physical Quantum observables must correspond to coordinate
independent quantities only. The search for these gauge independent observables, is started in the group that forms around Bergmann, Ted Newman develops a perturbation approach for finding gauge invariant observables order by order. The group studies the problems raised by systems with constraints, and on the problem of what are the observables in General Relativity.  The \textbf{Canonical} approach to Quantum Gravity is born.  \\

Bryce DeWitt completes his thesis.  He applies Schwinger's covariant quantization to the gravitational field.  \\

Dirac presents his method for treating constrained Hamiltonian systems.  \\

\vskip2mm{\bf 1952}\nopagebreak

Following the pioneering works of Rosenfeld, Fierz and Pauli, Gupta develops systematically the ``flat space quantization" of the gravitational field.  The idea is to introduce a fictitious ``flat space", that is, Minkowski metric $\eta_{\mu\nu}$, and quantize the small fluctuations of the metric around Minkowski $h_{\mu\nu}=g_{\mu\nu}-\eta_{\mu\nu}$. The \emph{Covariant approach is fully born}.  The first difficulty is immediately recognized, in searching the \textbf{propagator}, as the quadratic term of the Lagrangian is singular, as for the Electromagnetic Field, and as a consequence of Gauge Invariance. Gupta's treatment uses an indefinite \emph{norm} state-space as for the Electromagnetic Field.  \\


Charles Misner introduces the ``Feynman quantization of General Relativity".  He quotes John Wheeler for suggesting the expression
\begin{equation}
    \int exp\{(i/\hbar)({\rm Einstein\ action)}\}\ d({\rm field\
    histories)}),
    \label{eq:functional}
\end{equation}
and studies how to have a well-defined version of this idea. Misner's paper explains notions such as, why the Quantum Hamiltonian must be zero, why the individual spacetime points are not defined in the Quantum Theory, and the need of dealing with \emph{Gauge Invariance} in the
integral.  Even more remarkably, the paper opens with a discussion of the possible directions for quantizing Gravity, and lists lines of directions; \textbf{Covariant, Canonical, and Sum Over Histories}, describing them almost precisely with the same words we would today!  \\

\vskip.5cm

At the end of the fifties, all the basic ideas, and the research programs are clear.  It is only a matter of implementing them, and seeing if they work.  The implementation, however, a rather herculean task, that requires the ingenuity of people of the caliber of Feynman and DeWitt, on the \emph{Covariant side},  Dirac and DeWitt, on the \emph{Canonical side}.  \\

\vskip2mm{\bf 1959} \nopagebreak

By 1959, Dirac completely unraveled the \emph{Canonical structure of GR}.  \\

Arnowit, Deser and Misner complete formulation of GR, its Hamiltonian version in appropriate variables, simplifying the Hamiltonian formulation, and
make its geometrical reading transparent. In relation to the quantization, they present an influential argument for the \emph{finiteness} of the self \textbf{Energy of a point particle} in Classical GR, and use it to argue that \textbf{non-perturbative} Quantum Gravity should be finite.  \\


\vskip2mm{\bf 1962} \nopagebreak

Feynman showed the task of computing \textbf{transition amplitudes} in Quantum Gravity.  He shows that \emph{tree-amplitudes} lead to the Physics one expects, from the Classical Theory.

DeWitt starts developing his background field methods for the
computation of \emph{perturbative transition amplitudes}.

Bergmann and Komar clarify what one should expect from a Hilbert
space formulation of GR.

Peres writes the Hamilton-Jacobi formulation of GR
$$
G^{2}(q_{ab}q_{cd}-{1\over 2}q_{ac}q_{bd})\ {\delta S(q)\over \delta
q_{ac}}{\delta S(q)\over \delta q_{bd}} + \det q\,R[q]=0,
$$
which will lead to the Wheeler-DeWitt equation.

\vskip2mm{\bf 1963} \nopagebreak

John Wheeler shows that Quantum fluctuations of the Gravitational field, must be short scale fluctuations of the Geometry, and introduces the physical idea of spacetime foam. He suggests to study Quantum Gravity by means of a Feynman integral over a spacetime lattice.  \\

\vskip2mm{\bf 1964} \nopagebreak

Penrose introduces the idea of \emph{spin-networks}, and of a discrete
structure of space controlled by $SU(2)$ representation Theory.  The idea re-emerges 25 years later, when spin-networks will be found to label the states of \emph{Loop Quantum Qravity}.

\vskip2mm{\bf 1964} \nopagebreak

Beginning to study loop corrections to GR amplitudes, Feynman observed that \emph{unitarity}, was lost for diagrammatic rules.  DeWitt  develops the combinatorial means to correct the
quantisation (requiring independence of diagrams from the longitudinal parts of propagators).  These correction terms can be put in the form of loops of fictitious \textbf{Fermionic particles}, the Faddeev-Popov
ghosts. The key role of DeWitt in this context was emphasized by Veltman in 1974. It is in the work of Faddeev, that the key role played by the gauge orbits, (and not fields at a given point) as true \textbf{dynamical variable}s, is
fully elucidated.  \\

\vskip2mm{\bf 1967} \nopagebreak

Bryce DeWitt publishes the ``Einstein-Schr\"odinger equation".
$$
\left((\hbar G)^{2}(q_{ab}q_{cd}-{1\over 2}q_{ac}q_{bd}){\delta \over
\delta q_{ac}} {\delta \over\delta q_{bd}}-\det q\, R[q]\right)\
\Psi(q)=0,
$$
Bryce will long denote this equation as the ``Einstein-Schr\"odinger equation", attributing it to Wheeler; While John Wheeler denoted it as the DeWitt equation-- until, finally, in 1988, DeWitt calls it the way everybody else had been calling it, since the beginning: the ``Wheeler-DeWitt equation".  \\

John Wheeler discusses the idea of wave-function $\Psi(q)$ of the ``3-Geometry" $q$, and the notion of \emph{superspace}, the space of 3-Geometries.  \\

Penrose starts Twistor Theory.

The project of DeWitt and Feynman is concluded.  A set of Feynman rules for GR are written down.  \\

\vskip2mm{\bf 1969} \nopagebreak

Developing an idea in Bryce's paper on Canonical Quantum Gravity, Charles Misner starts Quantum Cosmology: Truncating the
Wheeler-DeWitt equation to a finite number of \emph{degrees of freedom}.  The idea develops into a long lasting industry from which, after a while, little new is understood.  \\

\vskip1cm

The decade closes with the main lines of the Covariant and the Canonical Theory defined. It will soon become clear that neither Theory works.  \\

\hspace{.6cm}{\bf 1970} \nopagebreak

A point made by Pauli, a paper by Zumino, suggests that the quantization of GR, may be problematic, and might make sense only by viewing GR as the low Energy limit of a more General Theory.  \\

\vskip2mm{\bf 1971} \nopagebreak

DeWitt and Feynman ideas on Gravity, t'Hooft, and Veltman, decide to study the re-normalizability of GR. They consider the re-normalization of Yang-Mills Theory, and find that the Theory is re-normalizable. One can say that the first Physical result of the research in Quantum Gravity, is the proof that Yang-Mills Theory is re-normalizable.  \\


\vskip2mm{\bf 1971} \nopagebreak

David Finkelstein writes his inspiring ``spacetime code" series of papers, which, among other ideas, discuss Quantum Groups.  \\

vskip2mm{\bf 1973} \nopagebreak

Following the program, t'Hooft finds evidence of un-renormalizable divergences in GR with matter fields. Veltman, Deser, and Van Nieuwenhuizen, confirm the evidence.  \\

\vskip2mm{\bf 1974} \nopagebreak

Hawking announces the derivation of black hole radiation. A (macroscopically) Schwarzshild black hole of
mass $M$ emits thermal radiation at the temperature
$$
       T = {\hbar c^{3}\over 8\pi k G M}
$$
The result, observed by Bekenstein, that entropy is associated to black holes, and thus they could be thought, in obscure sense, as ``hot", and
by the Bardeen-Carter-Hawking analysis of the analogy between laws of thermodynamics, and dynamical behavior of black holes. Hawking's result is not directly connected to Quantum Gravity --it is a skillful application of quantum field theory in curved spacetime-- but has a very strong impact on the field.  It fosters an intense activity in Quantum Field Theory in curved spacetime, it opens a new field of research in ``black hole thermodynamics", and opens the Quantum-Gravitational problems of understanding the \textbf{statistical} origin of the black hole (the Bekenstein-Hawking) entropy.  \\
For a Schwarzshild black hole, this is
\begin{equation}
       S = {1\over 4}\, {c^{3}  \over \hbar G}\, A
\label{entropy}
\end{equation}
where $A$ is the area of the black hole surface.  An influential, intriguing paper is written later, by Bill Unruh.  The paper points out the existence of a general relation between accelerated observers, Quantum Theory, Gravity, and Thermodynamics.  Something deep should be hidden in this tangle of problems, but we do not yet know what.  \\

\vskip2mm{\bf 1975} \nopagebreak

It becomes generally accepted, that GR \emph{coupled to matter} is not re-normalizable.  The research program started with Rosenfeld, Fierz
and Pauli stops.  \\


\vskip2mm{\bf 1976} \nopagebreak

An attempt to save the Covariant program, is made by Steven Weinberg, who explores the idea of \emph{asymptotic safety}, developing earlier ideas from Giorgio Parisi, Kenneth Wilson and others, suggesting that non-renormalizable Theories could nevertheless be meaningful.  \\

To resuscitate the Covariant Theory, even if in modified form, the path has already been indicated: Find a high Energy modification of GR. \emph{Preserving general Covariance, there is not much one can do to modify GR}. An idea that attracts much enthusiasm is \emph{super-Gravity}: It seems that by coupling a spin 3/2 particle to GR, namely with the action(in first order form)
$
 S[g,\Gamma,\psi] = \int d^{4}x\ \sqrt{-g}\ \left({1\over 2G} R -
 {i\over 2} \ \epsilon^{\mu\nu\rho\sigma}\ \psi_{\mu} \gamma_{5}
 \gamma_{\nu} D_{\rho} \psi_{\sigma} \right),
$
one can get a Theory finite, even at two loops. \\

\vskip2mm{\bf 1977} \nopagebreak

Another, independent idea, is to keep the same Kinematics, and change the Action. Add terms proportional to the divergences.  Stelle, proves that an Action, with terms quadratic in the curvature,
$
 S =
 \int d^{4}x\ \sqrt{-g}\
 \left(\alpha R+\beta R^{2}+\gamma R^{\mu\nu}R_{\mu\nu} .
 \right),
 $
is re-normalizable for appropriate values of the Coupling Constants. It has \textbf{negative Energy modes}, that make it unstable around the Minkowski vacuum, and not \emph{unitary} in the Quantum regime. The problem becomes to find a Theory re-normalizable, and \emph{unitary} at the same time, or to circumvent \emph{non-unitarity}.  \\


\vskip2mm{\bf 1978} \nopagebreak

The Hawking radiation is soon re-derived in a number of ways, strongly reinforcing its credibility.  Several of these derivations point to thermal techniques, thus motivating Hawking to revive the Wheeler-Misner ``Feynman Quantization of General Relativity" in the form of a ``Euclidean" integral over {\em Riemannian\/} 4-geometries $g$
$
  Z = \int \ Dg\ e^{-\int \sqrt{g}R}.
$
Time ordering, and the concept of positive frequency, are incorporated into the ``\emph{Analytic Continuation}" to the Euclidean sector, to deal with topology change, and that the Euclidean functional integral will prove to be a better calculation tool than the Wheeler-DeWitt equation.  \\


\vskip2mm{\bf 1980} \nopagebreak

Within the {\em Canonical\/} approach, the discussion focuses on understanding the disappearance of the \emph{Time Coordinate}, from the Wheeler-DeWitt
Theory. The problem has actually nothing to do with {\em Quantum\/} Gravity, since the \emph{Time Coordinate} disappears in the {\em Classical\/} Hamilton-Jacobi form of GR as well; The Physical Observables are \textbf{Coordinate Independent}, and thus, in particular, independent from the \emph{Time Coordinate}, in whatever correct formulation of GR. But in the Quantum context, there is no single spacetime, as there is no trajectory for a Quantum particle, and the very concepts of \textbf{Space, and Time, become fuzzy}.  This fact raises much confusion, and a vast interesting discussion on the possibility of doing meaningful \textbf{fundamental Physics}, in the absence of a fundamental notion of Time.  \\

\vskip2mm{\bf 1981} \nopagebreak

Polyakov shows that the cancellation of the conformal
anomaly in the quantization of the \emph{string action},
$
    S = {1\over 4\pi\alpha'} \int d^{2}\sigma \sqrt{g}
    \ g^{\mu\nu}
    \partial_{\mu}X^{a} \partial_{\nu}X^{b}\eta_{ab}.
$
leads to the critical dimension.  \\

\vskip2mm{\bf 1983} \nopagebreak

Two Physicists, Gary Horowitz, and Andy Strominger, summarize their contributed paper, with the words:
\begin{quote}
    In sum, \emph{higher derivative Gravity Theories} are a viable option for
    resolving the problem of Quantum Gravity \ldots
\end{quote}

Super-gravity is advertised as a likely final solution of the Quantum Gravity puzzle. It iss non-renormalizable at higher loops, and that higher derivatives Theories do not lead to viable \textbf{perturbative expansions}.  \\

In its' version in 11 dimensions, Super-gravity finds new importance in the late 1990s, in connection with \emph{String Theory}.  \\
High derivative corrections reappear, in the low Energy limit of String Theory.  \\

Hartle, and Hawking introduce the notion of the ``\textbf{wave-function of the Universe}" and the ``\emph{no-boundary condition for the Hawking integral}, opening up a new intuition on Quantum Gravity, and Quantum Cosmology.  But the Euclidean integral does not provide a way of computing genuine
field theoretical quantities in Quantum Gravity better than the Wheeler-DeWitt equation. \\
Jim Hartle, develops the idea of a \emph{Sum Over Histories} formulation of GR, into a full fledged extension of \textbf{Quantum Mechanics}, to the general Covariant setting.  The idea, will later be developed, and formalized by Chris Isham.

Sorkin introduces his \emph{poset} approach to Quantum Gravity.  \\

\hspace{.6cm}{\bf 1984} \nopagebreak

Green, and Schwarz realize that strings might describe ''our Universe''. In connection with the \emph{unexpected anomaly cancellation}, and the discovery
of the \emph{heterotic string}.

The relation between the ten dimensional Super-strings Theory, and four dimensional low Energy Physics is studied in terms of compactification
on Calaby-Yau manifolds, and orbifolds.  The dynamics of the choice of the vacuum remains unclear, but the compactification leads to 4d chiral models resembling low Energy Physics.

Belavin, Polyakov, and Zamolodchikov, publish their analysis of \emph{Conformal Field Theory}.  \\

Goroff, and Sagnotti finally compute the two loop divergences of pure GR, the divergent term is:  \\
$
 \Delta S = {209\over 737280\pi^{4}}\  {1\over\epsilon}
 \int d^{4}x\ \sqrt{-g}\
 R^{\mu\nu}{}_{\rho\sigma}R^{\rho\sigma}{}_{\epsilon\theta}
 R^{\epsilon\theta}{}_{\mu\nu} .
$

\vskip2mm{\bf 1986}  \nopagebreak

Penrose suggests that the wave-function collapse in \emph{Quantum Mechanics}, might be of \emph{Quantum Gravitational origin}.  The idea is radical, and implies a re-thinking of the basis of Mechanics.  \\

String Field Theory represents an attempt to address the main problem of String Theory: finding a fundamental, background independent, definition of the Theory.  \\

The connection formulation of GR, is developed by Abhay Ashtekar, on the basis of some results by Amitaba Sen. At the time, this is denoted the ``new variables" formulation.  It is a development in Classical General Relativity, but it has long ranging consequences on Quantum
Gravity, as the basis of Loop Quantum Gravity.  \\

\vskip2mm{\bf 1987}  \nopagebreak

Fredenhagen, and Haag explore the general Constraint that general
Covariance puts on Quantum Field Theory.  \\

vskip2mm{\bf 1987}  \nopagebreak

Green, Schwarz, and Witten, publish their book on Superstring
Theory.  In the gauge in which the metric has no super-partner, the super-string action is:
$
    S = {1\over 4\pi\alpha'} \int d^{2}\sigma \sqrt{g}
    \left( g^{\mu\nu}
    \partial_{\mu}X^{a} \partial_{\nu}X^{b} - i \psi^a
    \gamma^{\mu}\partial_{\mu}\psi^b\right) \eta_{ab}.
$
The interest in the Theory grows very rapidly. Iincreasingly, the research
in super-gravity, and higher derivative Theories, has merged into strings, and Sring Theory is increasingly viewed as a strong competing candidate for the Quantum Theory of the Gravitational Field.  As a side product, many particle Physicists begin to study General Relativity, or at least some bits of it.  \\  Strings provide a consistent \emph{perturbative theory}.  The Covariant program is fully re-born.  The problem becomes understanding why the World described by the Theory appears so different from ours.  \\


\vskip2mm{\bf 1988}  \nopagebreak

Ted Jacobson, and Lee Smolin, find loop-like solutions to the Wheeler-DeWitt equation, opening the way to Loop Quantum Gravity.  \\

The ``loop representation of Quantum General Relativity" is introduced. It is based on the new connection formulation of GR,
on the Jacobson-Smolin solutions, and on Chris Isham's ideas
on the need of non-Gaussian, or non-Fock representations in Quantum Gravity.  Loop quantization had been previously independently developed by Rodolfo Gambini, and his collaborators for Yang Mills Theories.  In the Gravitational context, the loop representation leads immediately to two results: An infinite family of exact solutions of the Wheeler-DeWitt equation is found, and Knot Theory controls the Physical Quantum states of the
Gravitational Field.  Classical Knot Theory, with its extensions, becomes a branch of Mathematics relevant to describe the different \textbf{invariant states of Quantum spacetime}.  \\

The Theory transforms the old Wheeler-DeWitt Theory in a formalism that can be concretely used, to compute Physical quantities in Quantum
Gravity.  The Canonical program is fully re-born.  Today, the Theory is called ``Loop Quantum Gravity". It is sometimes also called ``Quantum Geometry", but the expression ``Quantum Geometry" is used by a variety of other research programs as well. For a review,
complete references, and an account of the development of the Theory, see \url{https://arxiv.org/abs/gr-qc/9710008}.  \\

\vskip2mm{\bf 1988}  \nopagebreak

Ed Witten introduces the notion of Topological Quantum Field Theory
(TQFT). He uses a TQFT to give a field theoretical representation of the Jones polynomial, a Knot Theory invariant.  The expression used by Witten has an interpretation in loop quantum gravity: It can be seen as the
``loop transform" of Quantum state given by the exponential of the
Chern-Simon functional.  \\

Formalized by Atiyah, the idea of TQFT will strongly affect later development in
Quantum Gravity.  \\

Witten finds an ingenious way of quantizing GR in 2+1 spacetime dimensions, opening up a big industry of analysis of the Theory.  The quantization
method is partially a Sum Over Histories, and partially Canonical. Covariant perturbative quantization, seemed to fail for this Theory. The theory had been studied a few years earlier by Deser, Jackiw, t'Hooft, Achucarro, Townsend, and others.  \\

\vskip2mm{\bf 1989}  \nopagebreak

Amati, Ciafaloni, and Veneziano, find evidence that String Theory implies that distances smaller than the Planck scale, cannot be probed.  \\

In the string world, there is non-perturbative models of Strings ``in 0 dimension", equivalent to 2D
Quantum Gravity.  The models re-emerge in the nineties, and inspire the \textbf{Spin foam formulation of Quantum Gravity}.  \\

Turaev, and Viro define a state sum that is a rigorously defined TQFT, and can be seen as a
regulated, defined version of the Ponzano-Regge quantization of 2+1 Gravity. Turaev, and Ooguri find a 4d extension, which will have an impact on later developments.  \\

The notion of {\em weave\/} is introduced in Loop Quantum Gravity. It is a discrete structure of spacetime
emerging from Loop Quantum Gravity.  The first example of a weave, which is considered as a 3d mesh of intertwined rings. The intuition was already in Wheeler.  See Misner, Thorne, and Wheelers' book on Gravitation.  \\

\vskip2mm{\bf 1994}  \nopagebreak

Non-commutative Geometry, often indicated as a tool for describing certain aspects of Planck scale Geometry, finds a strict connection to GR, in the framework of Alain Connes' non-commutative Geometry. The Connes-Chamseddine ``Spectral Action", trace a simple function of a suitably defined Dirac-like operator $D$:
$
  S = Tr[f(D^{2}/(\hbar G))],
$
where $f$ is the characteristic function of the $[0,1]$ interval, including the standard model action, as well as the
Einstein-Hilbert action.  \\

\hspace{.6cm}{\bf 1995}  \nopagebreak

Non-perturbative aspects of String Theory begin to appear: Branes, Dualities, the matrix model formulation of M Theory \ldots .  The interest in Strings booms. Ed Witten claimed that:
\begin{quote}
    ``The Mathematics of the next millennium will be dominated
    by String Theory"
\end{quote}

The various dualities appear to relate the different versions of the
Theory, pointing to the existence of a unique fundamental Theory.  The actual construction of the fundamental background independent Theory, however, is still missing, and String Theory exists so far, only in the form of a number of (related) expansions over assigned backgrounds.  \\

\vskip2mm{\bf 1995}  \nopagebreak

Using the Spin Network orthonormal basis found on the Hilbert space of Loop Quantum Gravity, a main Physical result is obtained
within Loop Qantum Gravity: The computation of the eigenvalues of area, and volume.  The main sequence of
the eigenvalues of the area is labeled by an $n$-tuplet of half integers $\vec j=\{j_{1}\ldots j_{n}\}$, and is:
$
    A_{\vec j} = 8\pi\hbar G\ \sum_{i=1,n}\sqrt{j_{i}(j_{i}+1)}.
$
The result is rapidly extended, and derived in a number of ways.  \\

A rigorous Mathematical framework for Loop Quantum Gravity is developed.  \\

\url{https://arxiv.org/pdf/gr-qc/9901023}.  \\

\vskip2mm{\bf 1995} \nopagebreak

Ted Newman, and his collaborators, introduce the Null Surface Formulation of GR.  \\


\vskip2mm{\bf 1996}  \nopagebreak

The Bekenstein-Hawking Black Hole \emph{entropy}, is computed within Loop Quantum Gravity, as well as within String Theory.  \\

The loop result is obtained by computing the number of (Spin-Network) states, which endow a 2-sphere with a given area, as well as by loop quantizing the Classical Theory of the field outside the Hole, and studying the boundary states.  These Gravitational surface states can be identified with the states of a Chern-Simons Theory, on a surface with punctures.  The computation is valid for various realistic Black Holes, and the 1/4 factor is obtained by fixing the undetermined dimensionless parameter present in Loop Quantum Gravity (the Immirzi parameter).  \\

In String Theory, the computation exploits a strong coupling/weak coupling duality, which, in certain super-symmetric configurations, preserves the number of states: The Physical Black Hole is in a strong coupling situation, but the number of its' microstates can be computed in a weak field configuration, that has the same charges at infinity. The calculation method is thus rather indirect, and works smoothly, only for certain extremal Black Holes; One obtains precisely the 1/4 factor of equation, as well as other key aspects of the Hawking radiation phenomenology.  \\

A rigorously defined, finite, and anomaly free Hamiltonian constraint operator is constructed, by Thomas Thiemann, in Loop Quantum Gravity.  Some doubts are raised on whether the Classical limit of this Theory is in fact GR (the issue is still open), but the construction defines a consistent general covariant Quantum field Theory in 4d.  \\

Sum models obtained modifying a TQFT, are proposed by Barrett, Crane, Reisenberger, Iwasaki and others, as a tentative model for quantum GR. All these models appear as sums of ``Spin Foams": Branched surfaces carrying \emph{Spins}.  \\

The loop representation is ``exponentiated", giving rise, again, to a Spin Foam Model, corresponding to Canonical Loop Quantum Gravity.  These developments revive the \emph{Sum Over Histories} approach.  \\

There is a discussion on the difficulties of the \textbf{lattice approaches} in finding a second order phase transition.

Connections between non-commutative Gometry, and String Theory appear.  \\


\vskip2mm{\bf 1998} \nopagebreak

Juan Maldacena shows that the large $N$ limit of certain Conformal Field Theories includes a sector describing super-gravity, on the product of Anti-deSitter spacetimes, and spheres. He conjectures that the compactifications of M/String Theory on an Anti-deSitter spacetimes, is dual to a Conformal Field Theory, on the spacetime boundary.  This leads to a new proposal for defining M-theory itself, in terms of the Boundary Theory: An effort to reach
background independence (for M theory) using background dependent methods (for the Boundary Theory).

A consequence of this ``Maldacena conjecture" is an explosion of interest for an idea by Gerard t'Hooft, developed and promoted by Leonard Suskind: the ``\textbf{Holographic Principle}".  According to this principle (considered in a number of variants), the information on the Physical state in the interior of a region, can be represented on the region's boundary, and is limited by the area of this boundary.  \\

The idea is that, there is a number of different instances (the Neutral Kaon System, Gamma Ray Burst phenomenology, interferometers \ldots in which the presently operating \emph{measurement}, or observation apparata, or apparata that are going to be constructed, involve sensitivity scales comparable to --or not too far from the \textbf{Planck scale}.  If this direction fails, testing Quantum Gravity, might require the investigation of very early \textbf{Ccosmology}.  \\


For critical discussions of current direction of research in Quantum Gravity, see \url{https://link.springer.com/journal/10714}. \\


\section{Concluding remarks}

The lines of research have found many points of contact in the course of their development, andhave often intersected each other.  For instance, there is a formal way of deriving a Sum Over Histories formulation from a Canonical Theory, and viceversa; The perturbative expansion, can also be obtained expanding the Sum Over Histories; String Theory today faces the problem of a finding its' non-perturbative formulation, and thus the typical problems of a Canonical Theory, while Loop Quantum Gravity hasmutated into the Spin Foam Models, a Sum Over History formulation, using techniques that can be traced to adevelopment of String Theory of the early nineties. Lee Smolin has developed an attempt, to connect non-perturbative String Theory, and Loop Quantum
Gravity.  However, in spite of this continuous cross fertilization, the main lines of development have kept their essential separation.

The direction of investigation are identified by Charles Misner in 1959. In 1963, Peter Bergmann noted:
\begin{quote}
  ``In view of the great difficulties of this program, I consider it
    a very positive thing that so many different approaches are being
    brought to bear on the problem.  To be sure, the approaches, we
    hope, will converge to one goal."
\end{quote}
This was almost 40 years ago \ldots

The divide is particularly macroscopic between the Covariant line of research, and the Canonical/Sum Over Histories on
the other.  This divide has remained through over 70 years of research in Quantum Gravity.  The separation cannot be stronger.  Here is a
typical comparison, arbitrarily chosen among many.  On the covariant
side, Peter van Neuwenhuizen writes:
\begin{quote}
    ``\ldots Gravitons are treated on exactly the same basis as other
    particles such as Photons, and Electrons.  In particular, particles
    (including Gravitons) are always in flat Minkowski space, and move
    {\em as if\/} they followed their geodesics in curved spacetime,
    because of the dynamics of multiple Graviton exchange.  [\ldots]
    Pure Relativists often become somewhat uneasy at this point, because
    of the the following two aspects entirely peculiar to Gravitation:
    1) [\ldots].  One must decide before quantization which points are
    spacelike separated, but  it is only after
    quantization that that the fully quantized metric field can tell
    us this spacetime structure [\ldots].  2) [\ldots] In a Classical
    curved background one needs positive, and negative solutions, but in
    non-stationary spacetimes, it's not clear whether one
    can define such solutions.  {\em The strategy of particle
    Physicists has been to ignore these problems for the time being,
    in the hope that they will ultimately be resolved in the final
    Theory.  Consequently, we will not discuss them any further.}"
\end{quote}
On the Canonical side, Peter Bergmann says:
\begin{quote}
``The World point by itself, possess no Physical reality. It acquires
reality only to the extent that it becomes the bearer of specific
properties of the Physical fields imposed on the spacetime manifold."
\end{quote}
Partially, the divide reflects the different understanding of the World, that the particle Physics community, on the one hand, and the Relativity community, on the other hand, have.  The two communities have made repeated, and sincere efforts to talk to each other, and
understanding each other.  But the divide remains, and, with the divide, the feeling, on both sides, that the other side is incapable
of appreciating something basic, and essential: On the one side, the structure of Quantum Field Theory, as it has been understood in half a century of investigation; On the other side, the novel Physical
understanding of space, and time, that has appeared with General Gelativity.  Both sides expect that the point of the other will turn
out, at the end of the day, to be not very relevant.  One side, all the experience with Quantum Field Theory, is on a fixed metric spacetime, and thus is irrelevant in a genuinely background independent context.  The other because GR is only a low Energy limit of a much more complex Theory, and thus cannot be taken too seriously as an indication on the deep structure of Nature.  Hopefully, the recent successes of both lines, will force the two sides, finally, to face the problems that the other side considers prioritary: \emph{Background independence} on the one hand, control of \emph{a perturbation expansion} on the other.

\vskip.5cm

So, where are we, after 70 years of research?  There are well-developed tentative Theories, in particular Strings, Loops, and several other intriguing ideas. There is no consensus, no established Theory, and no Theory that has yet received any direct or indirect experimental support.  In the course of 70 years, many ideas have been explored, fashions have come, and gone, the discovery of the \emph{grail}, has been announced several times , with much later scorn.

However, in spite of its' age, the research in Quantum Gravity does not seem to have been meandering meaninglessly, when seen in its' entirety. On the contrary, one sees a logic that has guided the development of the research, from the early formulation of the problem, and the research directions in the fifties to nowadays. The implementation of the programs has been extremely laborious, but has been achieved.  Difficulties have appeared,
and solutions have been proposed, which, after much difficulty, have lead to the realization, at least partial, of the initial
hopes.  It was suggested in the early seventies that GR could perhaps be seen as the low Energy Theory of a \emph{Theory} without uncontrollable divergences; today, 30 years later, such a Theory --String Theory-- is known.  In 1957, Charles Misner indicated that in the Canonical framework, oneshould be able to compute eigenvalues; and in 1995, 37 years later, eigenvalues were
computed -- within Loop Quantum Gravity.  The road is not yet at the end, much remains to be understood, and some of the current
developments might lead nowhere.  But it is difficult to deny, looking at the entire development of the subject, that there has
been a \emph{linear progress}.  And the road, no doubts, is \emph{fascinating}.  \\


\chapter{Physical Ideas That Captivate}

\section{Characteristics of describing Energy}\label{sec:Characteristics of describing Energy}

We'll not expend electronic ink for philosophical arguments! i.e. How long will $\mathbb Z$ remain a bachelor, cold fusion, climate change, singularities, why wave-functions overlap, determinism(e.g. Wigners' friend, Schr$\ddot{o}$dingers' cat), why periodicity of Quantum numbers lead to periodicity of chemical elements, Spin-flip transitions(oscillating magnetic monopoles/distance transverse magnetic fields($\mathbf B$), produced by an oscillating magnetic dipole, etc....  \\
Ultra-cold atoms(lock-stepped with the same spin), have been realized in the laboratory(Bose-Einstein condensate[BEC]), through observing the center of mass system(a trick) of wave-like properties of atoms, so that the Pauli exclusion principle is NOT violated.  \\
This technology, although limited(interaction with the external envinroment, depolarizes particles), has applications in Quantum computers(cryptography - Quantum key distribution ; Refer to figure \ref{fig:Binary Entropy Diagram.png}), and is realized only for very low temperature Kelvin scales, not for Classical computers.  \\
Even the speed of light (C) - has been stopped!, using slow light pulses. The wave-particle duality(using a super-computer), was proved by an experiment in 1974, by John Clauser( concludes that the wave-function collapse is an objective $\mathbb R$ process - not a physical process, and not Relativistic Invariant - agreed). The idea of measurement of the wave-function, with any apparatus(alters the azimuthal Quantum number/angular momentum of the orbital plane - resulting to a breakup of Eigenvalues into a one-dimensional unitary representation cylindrical symmetry of 2j + 1, from spherical symmetry of j.  \\
Why, currently no answer/solutions...maybe its' even a wrong question. Considering Faradays' Law, knowing the Fourier transform, invoke a linear, and orthogonal transformation of Stokes' theorem, generalizing how the Electric Field remains Constant, the Magnetic Field too (curl$\mathbf E$) = - ($\frac{\partial \mathbf B}{\partial t}$), with $\iint_s \mathbf F dS$ = $\int_c \mathbf F dt$, where (C) is a closed curve, and $\mathbf F$ is a Vector Field defined in $\mathbb C$. It can be calculated, by making the Electric Field Constant, and varying the Magnetic Field e.g.
$\mathbf B(t)$ = $\langle tx, ty, -2tz \rangle$, $0<t<\infty$. Note that in general, the z component of angular momentum is $\hbar$m, with probability 1, and (m) is the state of the Magnetic Quantum number.  \\
This analysis is fiducial, and is based on the Theory of Complex Numbers, Maxwells' equations, Fourier transform, homogeneous linear differential equations, the double-slit experiment(the results on the screen is a transmitted split beam intensity profile), Stern-Gerlach experiment(represents an entangled Magnetic moment of Silver atoms), and Schr$\ddot{o}$dingers' equation.  \\
There are no visible sources of some forms though, they're(strange, and arbitrary) schematic representation of Physical ideas(Force, Momentum, Energy, Mass, Work, Acceleration. Length Of Interval, Time, and Velocity)...which are all approximations!  \\
The simultaneous probability distributions of various observables, given differential equations, which may be integrated, show how these observables change in time.  \\
Keep in mind, we've evolved from thinking the Sun was static, (It's a linear combination of a Stationary State[ \textbf{NOT} static], its' \emph{probabilities} don't change!!, and that it revolves around the Earth{defined also a linear combination of a Stationary State, with no changing probabilities), it's now the other way around.  \\

\begin{figure}
\begin{center}
\includegraphics[width=\textwidth, height=\textheight, keepaspectratio]{Binary Entropy Diagram.png}
\end{center}
\caption{Binary Entropy}
\label{fig:Binary Entropy Diagram.png}
\end{figure}

. For any questions, suggestion or comment, feel free to contact me on GitHub \href{https://github.com/StevenOwino}{issues} or email me at \email{stowino21@gmail.com}.

\paragraph{Axioms of Quantum Mechanics}
\noindent

In Dirac notation, an element of space state is a Ket($\mid \quad \rangle$). Time at $t_0$ is defined by specifying a Ket $\mid \Psi(t_0) \rangle$, belonging to  a Vector space.  \\

\begin{enumerate}
  \item \textbf{A1}: The wave-function is the complete description of a quantum system of an abstract vector space(state space $\epsilon$).
  \item \textbf{A2}: Hermitian operators are observables.
  \item \textbf{A3}: The measurement axiom; A measurable physical quantity A, is described by an observable A, acting on $\epsilon$.
  \item \textbf{A4}: Time evolution via the Schr$\ddot{o}$dinger equation.            
\end{enumerate}

\paragraph{Illustrating Vectors\index{Vectors}} \label{para:Illustrating Vectors}

\begin{figure}
\begin{center}
\includegraphics[width=12cm]{Vector pointing towards the Sun.png}
\end{center}
\caption{Vector pointing towards the Sun}
\end{figure}

\section{Some Conceptual Frameworks; \texttt{The names in color are Nobel prize receipients, and nomenees.}}

$\bullet$ Copenhagen interpretation(s) - $\color{blue}{\text{Bohr, Heisenberg, Pauli}}$, von Weis$\ddot{a}$cker, $\color{blue}{\text{Peierls}}$. \\  

$\bullet$ Non-local Hidden Variables - $\color{blue}{\text{de Broglie}}$, Bohm, Maudlin, Goldstein, Hardy. \\  

$\bullet$ Convival Solipism - Zwirn \\ 

$\bullet$ Modal Interpretations - Kochen, Dieks, van Fraassen, Bub, Sipe.  \\

$\bullet$ Conscienceness-induced Collapse - $\color{blue}{\text{Wigner}}$, von Neumann, $\color{blue}{\text{Josephson}}$.  \\

$\bullet$ Objective Collapse Models - Ghiradi-Rimini-Weber, Pearl, $\color{blue}{\text{Penrose}}$, Tumulka.  \\

$\bullet$ Consistent/Decoherent Histories - Griffiths, Omn$\acute{e}$s, $\color{blue}{\text{Gell-Mann}}$, Hartle.  \\

$\bullet$ Super-determinsm - $\color{blue}{\text{t'Hooft}}$, Hossenfelder.  \\
                                                
$\bullet$ Rational Interpretations(LQG) - Rovelli (\emph{ Approximates the curvature of Space, using the geometry of a Tetrahedron}) \\
                                                
$\bullet$ Ensemble Interpretations - Ballentine, $\color{blue}{\text{Einstein}}$  \\
                                                
$\bullet$ Informational Interpretation - $\color{blue}{\text{Zellinger}}$, Brukner, Bulb.  \\
                                                
$\bullet$ Many-Worlds Interpretation - Everett, Hawking, Science-fiction fans. \\


\begin{remark}  \\

$\mathbf{\textbf{Required reading:}}$ \\

$\bullet$ Spectrum of the harmonic oscillator \\
  
$\bullet$ Spectral Theorem \\
 
$\bullet$ Linear Algebra(particularly Matrix Theory)  \\

$\bullet$ Continuous Functional Calculus(systems of 2nd order differential equations for vector-valued functions, and factoring) \\
 
$\bullet$ Measure Theory  \\

$\bullet$ Equi-partition theorem(law of equi-partition of Energy)  \\

$\bullet$ Electromagnetic theory of light. \\

\end{remark}

\chapter{Terminology}
\section{Definitions} \label{sec:Definitions}

\begin{figure}
\begin{center}
\includegraphics[width=10cm]{qho.png}
\end{center}
\caption{This is a figure of the Quantum Harmonic Oscillator. It is defined by $E_n$ = ($n + \frac{1}{2}) \hbar \omega$}
\label{fig:qho.png}
\end{figure}

\begin{figure}
\begin{center}
\includegraphics[width=15cm]{relativistic mass increase.png}
\end{center}
\caption{Limit concept of Mass with the Speed of Light}
\label{fig:relativistic mass increase.png}
\end{figure}

\index{Jacobi Identity|(}

Quantum-Theory\index{Quantum-Theory} $\rightarrow$ A generalization/systematic refinement of classical mechanics \\
 
Quantum-communication/computation $\longrightarrow$ Manipulating qubits with quantum gates; Qubits being transmitted between different parties; Exploring applications based on sending information, rather than processing it.  \\

Quantum Fourier Transform\index{Quantum Fourier Transform} $\longrightarrow$ 
\begin{equation}
\phi = \sum_{k = 0}^{n - 1} y_k \mid k> = \frac{1}{\sqrt 2} \sum_{k = 0}^{n - 1} e^{ \frac{2 \pi {ijk}}{N}} \mid k> = \frac{1}{\sqrt 2} \sum_{k = 0}^{N - 1} \omega^{jk} \mid k> \text{\quad For N = 2, \quad k = 0,1 \quad and} \quad \omega = e^{\frac{2 \pi}{N}})
\end{equation}

  
An-entangled-pair\index{An-entangled-pair} $\longrightarrow$ Applying one qubit by manipulating the larger set of the four possible states.  \\

Von Neumann entropy\index{entropy} $\longrightarrow$
\begin{equation}
S(\rho) = - tr(\rho \ln \rho) = - \sum_j \lambda_j \ln \lambda_j \text{(These are diagonalized values);} S(\rho) = 0 \iff \rho^2 = \rho \text{(case of a pure state/projector)}
\end{equation}
 
Maximally mixed state $\longrightarrow$ $S(\rho_m)$ = - N $\frac{1}{N} \ln \frac{1}{N}$ = $\ln N$  \\

Gibbs entropy(classical statistical mechanics) $\longrightarrow$
\begin{equation}
S_G = - k_B \sum_i {P_i \ln P_i}, where \quad k_B \text{is the Bolzmann constant, and P is probabilities}
\end{equation}
  
Shannon entropy(classical information theory) $\longrightarrow$
\begin{equation}
H(X) = \mathbf E(\log_b P(X)) = - \sum_x P(x) \log_b P(x), \text{where H is for Shannon,} \mathbf E \text{is the expectation/average value, X is a variable, and P is probabilities}
\end{equation}
  
State-reduction\index{State-reduction} $\longrightarrow$

The effective two-level Hamiltonian( $\hat H_{eff}$ = $\begin{pmatrix}0 & \Omega \cr \Omega & \vartriangle - i \frac{\Gamma}{2} \end{pmatrix} )$, dissipating via spontaneous emission, during "measurement", included in the dynamics(A Quantum trajectory is NOT reversible).
Note that $[\hat H_{eff}, \hat H^\dagger_{eff}]$ $\neq$ 0, this Hamiltonian also generates a non-unitary dynamics $e^{- i \hbar_{eff}(t)}$ $\neq$ $e^{i \hat H^\dagger_{eff}(t)}$  \\

Measurement-of-a-state\index{Measurement-of-a-state}(indicates a loss of energy) $\longrightarrow$
\begin{equation}
\langle \psi \mid \psi \rangle = \langle \psi_0 \mid e^{- i \hat H_{eff}(t)} e^{i \hat H_{eff}(t)} \mid \psi_0 \rangle \rightarrow e^{- \Gamma(t)}
\end{equation}
\noindent

$\log_b$ $\text{can take any of the values 2[Classical bits], e[the Natural Logarithm], and Base 10[orders of magnitude]}$ \\

$H_{bin}(p)$ = $ -p \log p - (1 - p) \log (1 - p)$ $\rightarrow$ e.g. $H_{bin}(\frac{1}{2})$ = $ - \frac{1}{2} \log \frac{1}{2} - \frac{1}{2} \log \frac{1}{2}$ = $\log 2$  \\

Normalization\index{Normalization} - technical manipulation of Mathematical symbols(to
eliminate the physical vacuum contribution, and to simplify the one electron contribution).  \\
 
Vacuum-region\index{Vacuum-region} $\rightarrow$ discontinuity of fields in the small. \\ 

See figure \ref{fig:relativistic mass increase.png} \quad Limit\index{Limit} - $\forall$ total infinities, the degrees of freedom ceases to increase
without limit as shorter and shorter time intervals are considered. \\

Wave-equation\index{Wave-equation} $\longrightarrow$
\begin{equation}
\frac{\partial^2u(x,t)}{\partial t^2} = c^2 \frac{\partial^2 u(x,t)}{\partial x^2}  \text{[Relates acceleration and curvature].}
\end{equation}
  
Decoherence - The inevitable interaction of a system and its environment(quantum-to-classical transition; Its a fundamental limiting factor in the useful operation of Quantum Computers).  \\
                                                                         
Spin(total-angular-momentum)\index{Spin(total-angular-momentum)}/directional quantization) - a particle $\ell$ of non-negative integer or half integer. The Hilbert space for such a particle is
$\mathscr L^2(\mathbb R^3) \hat \otimes V_{\ell}$, where $V_{\ell}$ is an irreducible projective representation of SO(3)
of dimension 2$\ell$ + 1

Photons - A collection of an indefinite number of indistingishable particles, that obey Bose-Einstein statistics, with integer spin $\mathbb 1$.  \\

Boson(Bose-Einstein statistics) - A particle with an integer spin  \\

Fermion(Fermi-Dirac statistics) - A particle with half-integer spin e.g. an electron  \\

For j,k,l $\in$ {1,2,3} $\text{we define}$ $\epsilon_{j,k,l}$ $\text{by the formula:}$  \\
\begin{equation}
\epsilon_{j,k,l} = \left\{
\begin{array}{ll}
   1 \text{ if {j,k,l} is an even permutation of (1,2,3)}  \\
  -1 \text{ if {j,k,l} is  an odd permutation of (1,2,3)}  \\
   0 \text{ if any two of {j,k,l} are equal}  \\
\end{array}
\right.
\end{equation}

e.g $\epsilon_{3,2,1}$ = -1 and $\epsilon _{2,1,2}$ = 0.  \\

The commutation relations for the basis ${F_1,F_2,F_3}$ for SO(3) may be written(using the summation convention!) as:  \\

$[F_j,F_k]$ = $\epsilon_{j,k,l}F_l$  \\

i.e. if we take j=1 and k=2, then the sum on $\ell$ gives a nonzero value only when $\ell$ = 3, and we recover the relation $[F_1,F_2]$ = $F_3$  \\

Invariant\index{Invariant} $\rightarrow$ The density $\nu$ is irreducible under translations in the real directions. \\
 
Polarization\index{Polarization}\label{sec:Polarization} - The set of directions in which the elements of the quantum subspace are co-variantly constant.  \\
 
Covariant-derivative\index{Covariant-derivative} -$\rightarrow$ A conection $\nabla$ operator for a vector field.  \\
  
Integrability\index{Integrability} $\rightarrow$ Two complex vector fields X and Y lie in $P_z$ at each point z, then so does [X,Y] \\
 
Hermitian-line-bundle\index{Hermitian-line-bundle} $\longrightarrow$ (L,$\nabla$) connection of $L^{{\otimes}{k}}$ having a curvature $\frac{\omega}{\hbar}$ for any positive integer n(things behave nicely when k tends to $\infty$)  \\

Half-form-space\index{Half-form-space} $\longrightarrow$
\begin{equation}
s = (\psi \circ q ) \otimes \sqrt{q^*(\beta)}
\end{equation}

A measurable-function\index{measurable-function} $\psi$: X $\longrightarrow$ $\mathbb C$ is said to be integrable,
if $\int_X$ $\mid \psi \mid$ d$\mu$ < $\infty$    \\

Density of $\mu$ w.r.t. $\nu$ $\longrightarrow$
\begin{equation}
\mu(E) = \int_E \rho \quad d\mu \quad \forall E \in \omega
\end{equation}

supposing $\mu$ and $\nu$ are two $\sigma$-finite measures on a measure space(X,$\Omega$), 
and that $\mu$ is absolutely continuos w.r.t. $\nu$: then there exists a non-negative, measurable function $\rho$ on X. \\
 
Monotone-class\index{Monotone-class} $\longrightarrow$ Suppose $\mathscr M$ is a monotone class of subsets of a set X, and suppose $\mathscr M$ contains an
algebra $\mathscr A$ of subsets X, then $\mathscr M$ contain the $\sigma$-algebra generated by $\mathscr A$ \\
 
Dirac $\delta$ function\index{Dirac $\delta$ function} $\longrightarrow$
\begin{equation} \label{eq:Dirac Delta Function}
\delta_\chi = \chi(0),
    roughly \quad \delta(x - X) = \left\{
    \begin{array}{ll}
    \infty \text{ if x = X} \\
    0      \text{ otherwise} \\
    \end{array} \\
    \right.
\end{equation}
 
\textcolor{blue}{The function illustrates the limit of a distribution.}  \\

If $\Vert \cdot \Vert$: V $\longrightarrow$ $\mathbb R$, then

\begin{equation}
\Vert \psi \Vert = \sqrt{\langle \psi,\psi \rangle}
\end{equation}

$\Vert \cdot \Vert$ is a norm on V  \\

Hilbert-space\index{Hilbert-space} $\longrightarrow$ A vector space $\mathscr H$ over $\mathbb R$ or $\mathbb C$, equipped with an inner product $\langle \cdot, \cdot \rangle$,  
 such that $\mathscr H$ is complete in the norm above.       
 [For a seperable infinite dimensional $\mathscr H$, the partially ordered set of all questions in QM(Quantum Mechanics), is isomorphic
 to the partially ordered set of all closed sub-spaces].  \\

Orthogonal-space\index{Orthogonal-space} $\longrightarrow$ If V is any closed subspace of $\mathscr H$, a subspace $V^\perp$ of $\mathscr H$, is:  \\

\begin{equation}
V^\perp = {\phi \in \mathscr H \mid \langle \phi, \psi \rangle = 0 \quad \forall \psi \in V}\\
\end{equation}

Skew-self-adjoint\index{Skew-self-adjoint} $\longrightarrow$

\begin{equation}
A^* = - A \quad \text{One-parameter group of orthogonal/unitary transformation.}
\end{equation}

Orthogonal-projection $\longrightarrow$ For a closed subspace V, where V = range(P), P is the projection if it is any bounded operator
on $\mathscr H$ satisfying $P^2$ = P, and $P^*$ = P \\
 
A definite-state\index{definite-state}(random sequence) $\longrightarrow$
\begin{equation}
\Vert \bar{\sigma}_z^{N} \mid \psi \Vert^2 = \frac{1}{N^2} \langle \psi \mid \sum_{r=1}^N \sum_{s=1}^N \sigma_z^{(r)}\sigma_z^{(s)} \mid \psi \rangle
\end{equation}

\begin{equation}
\langle \psi \mid \sigma_z^{(r)} \sigma_z^{(s)} \mid \psi \rangle = \delta^{rs}
\implies \lim_{N \to \infty} \Vert \bar{\sigma_z}^{N} \Vert^2 = \lim_{N \to \infty} \frac{1}{N^2}N = 0
\end{equation}

Likewise for $\bar{\sigma}_z^a$ etc.  \\

Zorns'-Lemma\index{Zorns'-Lemma} $\longrightarrow$  \\

Let P be a set of elements a,b...Suppose there is a binary relation defined between certain pairs (a,b) of elements P,
expressed by a $\prec$ b, with the properties:  $\left\{
            \begin{array}{ll}
            a \prec a \\
            \text{ if a $\prec$ b and b $\prec$ a, then a = b,} \\
            \text{ if a $\prec$ b and b $\prec$ c, then a $\prec$ c(transitivity)} \\
            \end{array} \\
            \right.$  $\text{Then P is said to be partially ordered or semi-ordered, by the relation $\prec$}$.  \\

Curl\index{Curl} $\longrightarrow$ Measure of the rotation in the vector field about the points in the direction of the normal vector $\mathbf N$, and Stokes theorem justifies this interpretation. \\
 
Amperes Law(line integral along a closed path) $\longrightarrow$ $\oint \mathbf B \cdot d \mathbf l$ = $\mu_0 \mathbf l_{encl}$, 
where $\mu_0$ is the magnetic constant, and $l_{encl}$ is the net current enclosed by the path.  \\
  
Equi-partition-principle\index{Equi-partition-principle} $\longrightarrow$
\begin{equation}
E = \frac{3}{2}kT, \text{where k is the Boltzmans' constant.}
\end{equation}

See page \pageref{Sampled wave-function distribution} $\text{for Ground-state}$ $\longrightarrow$ An electron in a stable innermost orbit/lowest energy state(most negative), and cannot fall into the nucleus.  \\

Equilibrium position $\longrightarrow$ An electron in a stable innermost orbit/lowest energy state(most negative), and cannot fall into the nucleus.
(Its positive definite at its local minimum, hence its Energy is $\frac{\hbar \omega}{2 \pi}$ at 0). The equality $\frac{\hbar \omega}{2 \pi}$ is the Minimum Energy. $\psi_0$(x) = $(2 \pi \sigma^2)^{- \frac{1}{4}}$ $e^{- \frac{x^2}{4 \sigma^2}}$, with $\sigma^2$ = $\frac{\hbar}{2m \omega}$, and E = $\frac{\hbar \omega}{2}$. \\
 
Pauli-exclusion-principle\index{Pauli-exclusion-principle} $\longrightarrow$ No two or more electrons in an atom(Fermions), can have the same value, for any of the four quantum numbers. \\
 
Self-adjoint $\longrightarrow$ $T(F(\phi), \psi)$ = $T(F(\psi), \phi)$ = $T(\phi, F(\psi))$  \\
 
Eigenvalue-equation\index{Eigenvalue-equation} $\longrightarrow$ An equation, where the operator, operating on a function, produces a constant, times the function. [The function is called 
an Eigenfunction, and the resulting numeric is called the Eigenvalue]. Eigen here is the German word meaning self, or own.  \\
               
Hamiltonian-energy\index{Hamiltonian-energy}(discrete)   $\longrightarrow$
\begin{equation}
\hat {\mathbf H} = \frac{1}{2} m \omega^2_c \hat {r}^2_c
\end{equation}

\text{This is the Kinetic Energy of a particle of Mass m with Cyclotron Frequency} $\omega_c$, $\text{and radius}$ $\hat {\mathbf {r}_c}$.

Dirac equation $\longrightarrow$
\begin{equation}
i \hbar \frac{\partial \psi}{\partial t} =  \left[ c {\mathbf \alpha} \cdot \hat {\mathbf p} + \frac{e}{c} {\mathbf A} + {\beta m} {c^2} + V(r) \right] \psi
\end{equation}

Where the coupling of the electron to the scalar potential $\Phi(r)$ is included via V(r) = ${- e \Phi(r)}$.  \\

Gibbs Inequality $\longrightarrow$ $H(f_1,...,f_n)$ $\leq$ ${- \sum_{j=1}^{m} f_j log(u_j)}$ \\
 
Projection-valued-measure\index{Projection-valued-measure}associated with f(A) $\longrightarrow$ f(A')  \\
  
[ f(A) is the self-adjoint\index{self-adjoint} operator, corresponding to the observable $f(A')$ ]. \\ 

E $\rightarrow$ $P^{A}_{f^{-1}(E)}$. Note that the operators corresponding to the Energy and Momentum observables\index{observables} of a
quantum system are obtainable as infinitesimal generators of certain one-parameter groups\index{groups} of unitary operators. \\
                                                                        
Eigenvector\index{Eigenvector} of a self-adjoint operator $\longrightarrow$ A state in which the corresponding observable takes
on the corresponding Eigenvalues\index{Eigenvalues} with probability $\mathbf 1$ \\
 
Von Neumann-density-matrix\index{Von Neumann-density-matrix} $\longrightarrow$ The matrix of an operator\index{operator} defining a mixed state.  \\
  
Differentiable-function\index{Differentiable-function} $\psi$ [Schr$\ddot{o}$dingers equation] $\longrightarrow$
\begin{equation}
\frac{\partial^2 \psi}{\partial x^2} + \frac{\partial^2 \psi}{\partial y^2} + \frac{\partial^2 \psi}{\partial z^2} + \frac{2m}{\hbar^2} \left(\mathbf E + \frac{e^2}{\sqrt{x^2 + y^2 + z^2}}\right) = 0, with \quad \hbar = \mathbf K.
\end{equation}

Canonical-commutation-relation\index{Canonical-commutation-relation} $\longrightarrow$ [$\hat X, \hat P$] = $i \hbar \delta_{ij}$, i,j = x,y,z.  \\

Heisenberg-commutation-relation\index{Heisenberg-commutation-relation} $\longrightarrow$  \\

$(X_{\phi_1}, X_{\phi_2} - X_{\phi_2}, X_{\phi_1})$ = $(\frac{\hbar}{i})$ [$f_{\phi_1}, f_{\phi_2}]$,
where $f_\phi$ is an observable, and $X_\phi$ the self-adjoint operator.  \\

  
Discrete\index{Discrete} $\longrightarrow$ A representation which can be decomposed as a direct sum of irreducibles. \\
  
Schurs' Lemma\index{Schurs' Lemma}(only true in $\mathbb C$ representations, NOT $\mathbb R$ vector spaces) $\longrightarrow$ Let T be $\in$ R(L,M), then L restricted to the orthogonal complement of the null space
of T, is $\equiv$ to M restricted to the Closure\index{Closure} the Range of T.  \\
  
Degeneracy\index{Degeneracy} $\longrightarrow$ Occurrence of multiple eigenvalues/states sharing the same energy. \\ 

Zeeman-effect\index{Zeeman-effect} $\longrightarrow$ Splitting apart of eigenvalues, hence of spectral lines, in a magnetic field. \\
  
Kronocker-delta\index{Kronocker-delta}(Orthonormality condition\index{Orthonormality condition}) $\longrightarrow$
$\delta_{ij}$ = $\left\{
\begin{array}{ll}
1 \text{ if i = j}  \\
0 \text{ if i $\neq$ j} \\
\end{array} \\
\right.$ \\
 
Tensor/Kronecker-product\index{Tensor/Kronecker-product} of two Spin\index{Spin} one-half state space[needed to calculate the hyperfine-splitting\index{hyperfine-splitting} in the Hydrogen-atom\index{Hydrogen-atom}] $\longrightarrow$  \\

$ \mid 0, 0 \rangle$ =  \\
$\frac{1}{\sqrt 2}$ $\left( \mid \uparrow \rangle_1 \otimes \mid \downarrow \rangle_2 - \mid \downarrow \rangle_1 \otimes \mid \uparrow_2 \right)$ = $\frac{1}{\sqrt 2}$ $\left( \mid \uparrow \downarrow \rangle - \mid \downarrow \uparrow \rangle \right)$  \\
  
Clebsch-Gordan-coefficient\index{Clebsch-Gordan-coefficient}(expansion)/vector coupling/Wigner $\longrightarrow$ \\
\begin{equation}
\mid J,M \rangle = \sum_{m_1 = {-j_1}}^{j_1} \sum_{m_2 = {-j_2}}^{j_2} \langle{JM} \mid {m_1,m_2} \rangle \langle{m_1,m_2} \mid J^{'},M^{'} \rangle = \delta_{jj^{'}} \delta_{MM^{'}}
\end{equation}
    
Hydrogen\index{Hydrogen} hyperfine split $\longrightarrow$ $\mathbf{\text{Energy shift/difference}}$ of the lowest antiparallel state of the spin of the protons magnetic moment and the electrons magnetic field, pointing in the same direction(corresponds to the fine-structure\index{fine-structure} splitting).  \\  
    
Collision-broadening\index{Collision-broadening} $\longrightarrow$ Distortion of an emitted wave during radiation, thus smearing the spectral line. \\
 
See Eq. \eqref{eq:adjoint-action} \quad \textcolor{blue}{Inner-product}\index{Inner-product} $\longrightarrow$ <A|B> = $\frac{1}{2}tr(A,B^\dagger)$ \\

A Matrix changes  a derivative, so A = $frac{d}{dt}$. To find the transpose of this (A) we need to define
the Inner Product between two functions x(t) and y(t). \\

The Inner Product changes from the sum of $X_k Y_k$ to the integral of x(t) y(t). \\

Inner product of functions $X^{T}y$ = (X, y) = $\int_{- \infty}^\infty$ x(t)y(t)dt  \\

From this Inner Product, we know the requirement on $A^T$ . The word "adjoint" is more
correct than "transpose" when we are working with derivatives.  \\

The transpose of a matrix has $(Ax)^{T} \mathbf y$ = $\mathbf{x}^T$ $(A^T \mathbf y)$. The adjoint of A = $\frac{d}{dt}$ has:

(Ax,y) = $\int_{- \infty}^\infty$ $\frac{d}{dt}$ y(t) dt = $(x, A^{T}y)$

Recognize integration by parts. The derivative moves from the first
function x(t), to the second function y(t). During that move, a minus sign appears.
This tells us that the transpose of the derivative is minus the derivative.

The derivative is antisymmetric: $\mathbf{A}$ = $\frac{d}{dt}$, and $\mathbf{A^T}$ = ${- \frac{d}{dt}}$. Symmetric matrices have $S^T$ = S, antisymmetric matrices have $A^T$ = -A. Linear algebra includes derivatives and integrals, because those are both linear.  \\

\textcolor{blue}{This anti-symmetry of the derivative applies to centered difference matrices.} A =  \\

$
\begin{bmatrix}
0 & 1 & 0 & 0 \\
-1 & 0 & 1 & 0 \\
0 & -1 & 0 & 1 \\
0 & 0 & -1 & 0
\end{bmatrix}
$
\noindent

\textcolor{blue}{;transposes to} $A^T$ =  \\

$
\begin{bmatrix}
0 & -1 & 0 & 0 \\
1 & 0 & -1 & 0 \\
0 & 1 & 0 & -1 \\
0 & 0 & 1 & 0
\end{bmatrix}
$

= -A.  \\

And a forward difference matrix transposes to a backward difference matrix, multiplied
by -1. In differential equations, the second derivative (Acceleration) is symmetric. The
first derivative (Velocity) is antisymmetric.  \\

 
Entanglement\index{Entanglement} $\longrightarrow$ A Bell state that cannot be expressed as a tensor product of a single qubit state with another single qubit state. \\
 
No-cloning-theorem\index{No-cloning-theorem} $\longrightarrow$ One cannot create a copy of quantum states, although teleporting is permitted.  \\
  
Teleportation\index{Teleportation} $\longrightarrow$ Same quantum state in different locations.  \\

See \ref{fig:qho.png} for the Quantum-Harmonic-Oscillator\index{Quantum-Harmonic-Oscillator} $\longrightarrow$  \\
\noindent
$\hat H$ = $\hbar \omega$  $\left(\hat {a}^\dagger \hat {a} + \frac{1}{2} \right)$ = $\hbar \omega$ $\left(\hat N + \frac{1}{2} \right)$, where $\hat {a}$ = ${\sqrt \frac{m \omega}{2 \hbar}}$ $\left( \hat x + \frac{i}{m \omega} \hat p \right)$; \\
$\hat {a}^\dagger$ = ${\sqrt \frac{m \omega}{2 \hbar}}$ $\left( \hat x - \frac{i}{m \omega} \hat p \right)$, and  \\
$\hat N$ = $\hat {a}^\dagger \hat a$ = $\frac{m \omega}{2 \hbar}$ $\left(\hat x - \frac{i}{m \omega} \hat p \right)$ $ \left( \hat x + \frac{i}{m \omega} \hat p \right)$  \\
 
Flux-quantum\index{Flux-quantum}(a dimensionless quantity) $\longrightarrow$ $\Phi$ = $\frac{\hbar}{2e}$, where 2e is the charge of the electron(Cooper pairs). \\
     
Von Neumann-equation\index{Von Neumann-equation} $\longrightarrow$ $\partial_t$ $\rho(t)$ = $- {\frac{i}{\hbar}}$ $[H,\rho(t)]$, where $\rho$ is the density matrix, and H is the Hamiltonian;
similar to the Schr$\ddot{o}dinger$ equation, $i\hbar \partial_t \rho(t)$ = $[H, \rho(t)]$(describes the time-evolution\index{time-evolution} of the density operator).  \\
  
Resolution-of-an-identity\index{Resolution-of-an-identity} $\longrightarrow$ $\sum_{k_0} \sum_{\mathbf k_1}$ $<k_0|A_0| {\mathbf k_1}><{\mathbf k_1}|A_1|k_0>$ \\

Analytic definition of an angle $\longrightarrow$


%\end{Definitions}

\section{Ideas}

\begin{figure}
\begin{center}
\includegraphics[width=10cm]{Anharmonic vs simple Harmonic Oscillator.png}
\end{center}
\caption{Representation of Anharmonic vs simple Harmonic Oscillator}
\begin{center}
\label{fig:Anharmonic vs simple Harmonic Oscillator.png}
\end{center}
\end{figure}

$\mathbf{\textbf{Ideas that work[linear, and angular momentum of a field}}$

$\bullet$ A perfect crystal lattice( i.e. a One-Dimensional system) is effectively empty, for an electron in the lowest band,
even though the space is full of atoms.  \\

$\bullet$ The absolute value of Energy is equivalent to the fraction of a reduced mass of a proton-electron system, times Charge to the fourth power,
over twice an integer squared times the square of the Plancks' constant $\longrightarrow$ $|E|$ = $\frac{\mu Q^4}{{2n^2} \hbar^2}$  \\
  
$\bullet$ There are two special classes of polarizations(See the definition on Subsection~\ref{sec:Polarization}), those that are purely real(i.e, $\bar P_z$ = $P_z$ $\forall$ z $\in$ N),
and those that are purely complex(i.e, $P_z \cap \bar P$ = {0} $\forall$ z $\in$ N). \\
  
$\bullet$ The half-form of the Hilbert space is the completion w.r.t. $\vert s \vert^2$ = $\int_\Xi$ ${\widetilde{(s1,s2)}}$  of the space
of polarized for which $\vert s \vert^2$ < $\infty$  \\
  
$\bullet$ Pairing maps $\longrightarrow$ $\int_\mathbb R$ $\overline {\int_\mathbb {R} \phi (x) e^{- \frac{{ixp}}{\hbar}} dx}$ $\psi$ (p) dp   \\
 
$\bullet$ Dirac $\delta$ function $\longrightarrow$ $\delta_xi$ = $\xi$(0)  \\

$\bullet$ Faradays' Law(The Electric field is the negative of the rate of change of the corresponding Magnetic field w.r.t time) $\longrightarrow$ $curl \mathbf E$ = - $\frac{\partial \mathbf B}{\partial t}$ \\
 
Linear-system\index{Linear-system} $\longrightarrow$ Its properties satisfy a finite-dimensional vector space, has a posive-definite quadratic form,
and is independent. \\
 
Configuration-vector\index{Configuration-vector} $\phi$ $\longrightarrow$ $\frac{d}{dt}(\psi, \phi)$ = $\psi, {-I}^*(\phi)$ or $\frac{d\phi}{dt}$ = $\psi$ $\frac{d\psi}{dt}$ = $-I^*(\phi)$,
satisfies the second-order equation $\frac{d^2\phi}{dt}$ = ${-I^*(\phi)}$, where $\phi$ is a scalar
or vector-valued function on 3-space, and ${I^*}$ is a differential operator. \\
 
$\bullet$ $\mathbf{\textbf{Physical-law\index{Physical-law}}}$ (Seiberg-Witten-equations\index{Seiberg-Witten-equations}(useful on sympletic 4-manifolds)), is an example of a class of $\phi$s''\index{class of $\phi$'s}(a scaling-factor\index{scaling-factor}). proves one suitable existance and uniqueness theorems of the class of $\phi$s') $\longrightarrow$ When $\phi$ is regarded as a function of x,y,z and t, then it becomes a partial differential equation of the form $\frac{\partial^2 \phi}{dt^2}$ = ${-I^*}(\phi)$. \\
 
$\bullet$ Light $\longrightarrow$ Vibrations in the electromagnetic-field\index{electromagnetic-field}  \\

Energy $\longrightarrow$ An observable corresponding to $\hbar$, times the dynamical operator, and that its an integral
so that we have an $\mathbf{\textbf{energy conservation law}}$. \\
 
$\bullet$ Heisenberg-commutation-relation\index{Heisenberg-commutation-relation} $\longrightarrow$ ${Q_j}{Q_k} - {P_j}{P_k}$ = $i$ $\hbar$ $\delta_j^k$ \\
 
$\bullet$ Wave-function of a system\index{Wave-function of a system} $\longrightarrow$ A state vector $\psi$, is a square summable function on $\mathscr M$, and
trajectory $e^{-i Ht}(\psi)$ is a function on $\mathscr M \times \mathbf R$, where $\mathbf R$ is the real line(This is a function of $\mathbf{3n + 1}$ real variable).  \\
  
$\bullet$ Expected-value\index{Expected-value} of the j-th Momentum component $\longrightarrow$ $\frac{\hbar}{i}$ $\int$ $\left(\frac{\partial \psi}{\partial q_j} \overline {\psi} \right)$ $dq_1...dq_{3n}$.  \\

Negative Energy $\longrightarrow$ Bound State


  
\chapter{Physics of the Universe}

\section{Historical Timeline of the Quantization of Light}

\begin{table}[ht]
\centering
\caption{Historical Timeline of the Quantization}
\label{tab:quantization}
\begin{tabular}{|l|p{6cm}|}
\hline
\textbf{Year} & \textbf{Description} \\ \hline
\textbf{1670--1678}          &  \textbf{Isaac Newton} proposes the corpuscular theory of light, suggesting light consists of particles.  \\ \hline    \textbf{1801}                &  \textbf{Thomas Young} performs the double-slit experiment, demonstrating the wave nature of light. \\    \hline
\textbf{1864}                &  \textbf{James Clerk Maxwell} formulates the theory of electromagnetism, describing light as an electromagnetic wave.  \\ \hline
\textbf{1900}                &  \textbf{Max Planck} introduces the concept of quantization to explain blackbody radiation, proposing that energy is emitted or absorbed in discrete packets called \textit{quanta}. \\  \hline
\textbf{1905}                &  \textbf{Albert Einstein} explains the photoelectric effect using Plancks' quantum hypothesis, suggesting that light consists of discrete packets of energy called \textit{photons}  \\  \hline
\end{tabular}
\end{table}

where:  \\
\noindent

E = h $\nu$ \\


\subsection{Key Takeaways}

The quantization of light marked a paradigm shift in Physics, leading to the development of Quantum Mechanics. Light exhibits both wave-like and particle-like properties, a concept known as $\textbf{wave-particle duality\index{wave-particle duality}}$.

\section{Manifolds of pairing maps on an invariant self-adjoint subspace} \label{sec:Manifolds of pairing maps on an invariant self-adjoint subspace}

\subsection{Constants\index{Constants}}
$\mathbf{\textbf{A few constants:}}$  \\

$\hbar$*c $\simeq$197.327 MeV $ \cdot$ fm  \\

Coupling/Fine structure constant\index{Coupling/Fine structure constant} $\alpha$ $\rightarrow$ $\frac{e^2}{4 \pi \epsilon_0 \hbar c}$ $\simeq$ $\frac{1}{137.036}$, $m_e$ $c^2$ $\simeq$ 0.511 MeV, $m_p$ $c^2$ $\simeq$ 9.38 MeV. 
Described as deviations in the electron spin/removal of energy degeneracy of the orbital quantum number $\ell$:  
E = $E_{n,\ell}$[Energy shifts, and splitting of spectral lines]. \\

Bohr Radius $\longrightarrow$ $r_B$ = $\frac{4 \pi \epsilon_0 \hbar^2}{e^2{m_e}}$ $\approx$ 0.53 $\times$ $10^{-10}$m  

\subsection{Useful Formulae}

$\bullet$  (1) Relativity: p = mv$\gamma$, E = $\gamma$ m $c^2$, $E^2$ = $p^2$ * $c^2$ + $m^2$ $c^4$, $\gamma$ = $\frac{1}{\sqrt1 - \beta^2}$,
$\beta$ = $\frac{v}{c}$  \\

$\bullet$ (2) de Broglie: $\lambda$ = $\frac{h}{p}$,  \\
Compton: $\lambda_C$ = $\frac{2 \pi h}{mc}(1 - \cos \theta)$,  \\

\noindent
$\hat{p}$ = $\frac{h}{i}$ $\frac{\partial}{\partial x}$,  \\
$\hat{\mathbf p}$ = $\frac{\hbar}{i}$ $\nabla$,  \\
$[\hat{x_i},\hat{p_j}]$ = $i\hbar \delta_{ij}$,  \\
$[\hat{p_i},f(\hat{\mathbf x})]$ = $\frac{\hbar}{i}$ $\frac{\partial{f}} {\partial{x_i}}$.  \\

The equations above illustrate Photo-Electron Energy and interaction free measurements.  \\

$\bullet$ (3) Schr$\ddot{o}$dinger-equation\index{Schr$\ddot{o}$dinger-equation}:  \\
i $\hbar$ $\frac{\partial \psi}{\partial(t)}$ (x,t) =  $(\frac{- i \hbar^2}{2 m}$ $\nabla^2$ + V (x,t)) $\psi (x,t)$,  \\
$\frac{\partial}{\partial(t)}$ $\rho (x,t)$ + $\nabla$ $\cdot$ $\mathbf J$ (x,t) = 0,   \\
$\rho$ (x,t) = | $\Psi (x,t) |^2$,  \\
$\mathbf J$ (x,t) = $\frac{\hbar} {m}$ Im $[\psi^*$ $\nabla \psi]$. \\

$\bullet$ (4) Fourier transforms: $\psi x$ = $\frac{1}{\sqrt2 \pi}$ $\int{dk}$ $\phi$ (k) $e^{ikx}$,  \\
$\Phi (k)$ = $\frac{1}{\sqrt2\pi}$ $\int{dx}$ $\Psi$ (x) $e^{-ikx}$,
$\int${dx}| $\Psi (x) |^2$ = $\int {dk}$ | $\Phi k |^2$.
$\frac{1}{2\pi}$ $\int_ {-\infty}^{\infty}$ $e^{ikx}{dx}$ = $\delta(k)$, $\frac{1}{(2\pi)^3}$ $\int$ $e^{ik}$ $\cdot$ ${x}(d)^3{x}$ = $\delta^{(3)}(k)$.  \\
  
From the distribution-integral\index{distribution-integral} $\rightarrow$
\begin{equation}
\int_{-\infty}^{\infty} e^{(-ax^2 + bx)} dx = \sqrt{ \frac{\pi}{a}} e^{\frac{b^2}{4a}}, \quad when \quad \Re{(a)} > 0.
\end{equation}

$\bullet$ (5) expectation value:  \\
$\langle\hat{\mathbf Q} \rangle$ $\equiv$ $\langle\psi,  \\
\hat{\mathbf Q}\psi \rangle$, $i \hbar \frac{d}{dt} \langle \hat{\mathbf Q \rangle}$ = $\langle[ \hat{\mathbf Q},{\mathbf H}] \rangle$,
$\hat{\mathbf Q}$ time independent.  \\

$\bullet$ (6) Uncertainty:  \\

Consider Schwarz Inequality. Since $[\sigma^2_Q, \sigma^2_R]$ $\geq$ $(\frac{1}{2i} \langle [\hat{Q}, \hat{R}])^2$,
\begin{equation}
[\sigma^2_x, \sigma^2_p] = (\frac{1}{2i} \langle i \hbar \rangle)^2 \implies (\frac{[i \hbar]}{2i})^2
\end{equation}

 Resulting to [$\frac{\hbar^2}{4}$] \\

So, $\boxed{[\sigma^2_x, \sigma^2_p$] = $\frac{\hbar}{2}}$  \\

$\vartriangle{A} \equiv \vert {A} - \langle \hat{A} \rangle \mathbf{1} \psi \vert$  \\

$(\vartriangle{A})^2$ = $\langle \hat{A}^2 \rangle - \langle \hat{A} \rangle^2 \geq$ 0.  \\

$\vartriangle{A} \vartriangle{B}$ $\geq$ $\langle \psi$ \quad $\mid$ $\frac{1}{2i}$ $[\hat{A},\hat{B}]$ \quad $\mid$ $\psi$ $\rangle$  \\

$\bullet$ (7) saturation\index{saturation}: $(\hat{B} - \langle \hat{B} \rangle \mathbb1) \mid \psi \rangle$ = i $\lambda (\hat{A} - \langle \hat{A} \mathbb1)\mid \psi \rangle$, $\lambda \in \Re$.  \\
$\vartriangle{x} \vartriangle{p} \geq \frac{\hbar}{2}$. For $\psi$ ~ $\exp(- \frac{1}{4} \frac{x^2}{\vartriangle^2})$, \\
$\vartriangle{x} = \vartriangle$ and $\vartriangle{p} = \frac{\hbar}{2 \vartriangle}$. \\
$\vartriangle{H} \vartriangle{t} \geq \frac{\hbar}{2}$, $\vartriangle{t}$ $\equiv$ $\frac{\vartriangle Q } {\mid \frac{d Q}{dt} \mid}$.  \\

$\bullet$ (8) stationary state: $\psi$(x,t) = $\psi (x)e^{- \frac{iEt}{\hbar}}$  \\
\begin{equation}
- \frac{\hbar^2}{2m} \nabla^2 \psi(x) + \mathbf V(\mathbf x) \Psi(\mathbf x) = E \psi(\mathbf x)
\end{equation}

$\bullet$ (9) variational-principle\index{variational-principle}: $E_{gs} \leq \langle \psi, \hat{H} \psi \rangle$, for all normalized $\psi$. \\

$\bullet$ (10) Hellmann-Feynmann: $\hat{H}$ $(\lambda)$ $\psi_n$ $(\lambda)$ = $E_n$ $(\lambda)$ $\psi_n$ $(\lambda)$ $\longrightarrow$ $\frac{dE_n (\lambda) } {d \lambda} $ = $\langle \psi_n (\lambda), \frac{d \hat{\mathbf H}{\lambda}} {d \lambda} \Psi_n(\lambda) \rangle$.  \\

$\bullet$ (11) virial Theorem: One dimension/Single particle moving under a central potential:
\begin{equation} \label{eq:virial Theorem}
\langle \frac{\hat{p}^2}{2m} \rangle = \frac{1}{2} \langle \hat{x}  \frac{dV}{dx} \rangle
\end{equation}
\noindent

The virial Theorem is the principal tool, in Classical Kinetic Theory, for calculating the equation of state corresponding to an imperfect gases; It can be used to derive ideal gas law for perfect gases.  \\

$\bullet$ (12) Central potential: $\langle \frac{\hat{p}^2}{2m} \rangle$ = $\frac{1}{2}$ $\langle{r}$ $\frac{\partial{V}}{\partial{r}} \rangle $.  \\

$\bullet$ (13) \textcolor{blue}{Commutator Identities:} \\

[A,BC] = [A,B]C + B[A,C] \\

\begin{equation} \label{eq:adjoint-action}
e^A Be^{-A} = e^{ad_A}B = B + [A,B] + \frac{1}{2!} [A,[A,B]] + \frac{1}{3!} [A,[A,[A,B]]] +...,
\end{equation}
 
$\text{This is Hadamards' Lemma using Taylors' expansion, with}$
\begin{equation}
(ad_A) \text{as an adjoint action;}\quad (ad_A(X) = [A,X])
\end{equation}

Note that $e^A$ $e^B$ $\neq$ $e^{A + B}$ = $e^C$, with C = A + B + $\frac{1}{2}$[A,B] + $\frac{1}{12}$([A,[A,B]] - $\frac{1}{12}$[B[A,B]]) + ...,  
Its' not evident what the higher terms(...) refer to, although its' been expanded and expressed in terms of commutators. \\
 
$e^A Be^{-A}$ = B + [A,B], if  [A,[A,B]] = 0, \\ 

$[B,e^A]$ = [$B,A]e^A$, if [ [A,[A,B]] = 0, \\
 
$e^{A + B}$ = $e^A e^B e^{-\frac{1}{2}|A,B|}$ = $e^B e^A e^{\frac{1}{2}[A,B]}$ if [[A,B],A] = [[A,B],B] = 0. \\

$\bullet$ (14) matrix exponential\index{matrix exponential}: $e^{i M \theta}$ = $\mathbf{1} \cos \theta$ + i M $\sin \theta$,  \\

$\bullet$ (15) index-manipulation\index{index-manipulation}:  \\
$\mathbf a \cdot \mathbf b$ = $a_i b_i$,  \\
$\delta_{ij} B_j$ = $B_i$, $\delta_{ii}$ = 3, \\
${(\mathbf a \times \mathbf b)}_i$ = $\epsilon_{ijk} \mathbf a_j \mathbf b_k$, $\epsilon$ 123 = 1  \\
$\epsilon_{ijk} \epsilon_{ipq}$ = $\delta_{jp} \delta_{kq}$ - $\delta_{jp} \delta_{kp}$,
\begin{equation}
\epsilon_{ijk} \epsilon_{ijq} = 2 \delta_{kq}
\end{equation}

$\bullet$ (16) linear algebra: matrix representation in the basis\index{matrix representation in the basis} $(v_{1,...,}v_n)$: $Tv_i$ = $\sum_k $ $T_{kj}$ $v_k$.  \\

$\bullet$ (17) basis change\index{basis change}:  \\

$u_k$ = $\sum_j A_{jk}v_j$,\\
\begin{equation}
T({u}) = A^{-1} T ({v}) A
\end{equation}

$\bullet$ (18) Inner-product\index{Inner-product}:  \\

\begin{equation} \label{eq:Inner-Product}
\langle{v,v} \rangle \geq 0, for all v, \langle{v,v} \rangle = 0 \iff v = 0, \langle{u,v} \rangle = \langle{vu} \rangle^*
\end{equation}

$\bullet$ (19) Schwarz inequality:  \\
\begin{equation}
\mid (u,v) \mid \leq \Vert u \Vert \quad  \Vert v \Vert
\end{equation}

$\bullet$ (20) Complex-vector-space\index{Complex-vector-space}:  \\

$\langle v,Tv \rangle$ = 0,  \\

$\forall$ $\nu$ $\in$ V $\longrightarrow$ T = 0.  \\

$\bullet$ (21) adjoint\index{adjoint}: \\

$\langle{u,Tv} \rangle$ = $\langle T^\dagger{u,v} \rangle$,  \\

$(T^\dagger)^\dagger$ = T,  \\

$(ST)^\dagger$ = $T^\dagger S^\dagger$. \\

$\bullet$ (22) Hermitian operator T\index{Hermitian operator T}: $T^\dagger$ = T \\ 

$\bullet$ (23) unitary operator U\index{unitary operator U}: $U^\dagger U$ = $UU^\dagger$ = $\mathbf{1}$.  \\

$\bullet$ (24) orthogonal-projector\index{orthogonal-projector}: $P_u$: V $\longrightarrow$ U,  \\

${P_uP_u}$ = $P_u$, \\

$P_u^\dagger$ = $P_u$,  \\

V  = $\text{range}$ $P_u$ $\oplus$ $\text{null}$ $P_u$. \\

$\bullet$ (25) bra-kets\index{bra-kets}:

$\langle$ $u \mid v \rangle$ $\equiv$ $\langle { u,v} \rangle$,  \\

$\mid Tv \rangle$ $\equiv$ $T \mid v \rangle$,  \\

$\langle Tv \mid$ = $\langle v \mid T^\dagger$, \\

$\langle u \mid T^\dagger v \rangle$ = $\langle Tu \mid v \rangle$. \\ 

$\mid \alpha_1 v_1$ + $\alpha_2 v_2 \rangle$ = $\alpha_1 \mid v_1 \rangle$ + $\alpha_2 \mid v_2 \rangle$,  \\ 

$\langle \alpha_1 v_1$ + $\alpha_2 v_2 \mid$ = $\alpha_1^* \langle v_1 \mid $ + $\alpha_2^* \langle v_2 \mid $.  \\  

T = $\mid u \rangle$ $\langle w \mid$ $\longrightarrow$ $T^\dagger$ = $\mid w \rangle$ $\langle u \mid$, \\ 

$\mathbf{\textbf{tr}}$T = $\langle \mathbf w \mid \mathbf u \rangle$.  \\

$\bullet$ (26) orthonormal-basis\index{orthonormal-basis} $\mid i \rangle$: \\

$\langle i \mid j \rangle$ = $\delta_{ij}$,  \\

$\mathbf{1}$  = $\sum_i \mid i \rangle \langle i \mid$.  \\

$T_{ij}$ = $\langle i\mid T \mid j \rangle$ $\longleftrightarrow$ T = $\sum_{i,j} T_{i,j}$ $\mid i \rangle$ $\langle j \mid$,  \\

$\langle i \mid T^\dagger \mid j \rangle^*$,  \\ 

$(T^\dagger)_{ij}$ = $(T_{ji})^*$.  \\  

$[M,M^\dagger]$ = 0 $\longleftrightarrow$ $\text{M is unitarily}$ diagonizable\index{unitarily diagonizable}.  \\

$\bullet$ (27) Spectral-Theorem\index{Spectral-Theorem}:   \\

$T^\dagger$ = T $\longrightarrow$ T = $\sum_k \lambda_k P_k$, \\

$P^\dagger_k$ = $P_k$,  \\

$P_k P_l$ = $\delta_{kl} P_l$,  \\

$\sum_k P_k$ = $\mathbf{1}$.  \\

$\bullet$ (28) Schr$\ddot{o}$dinger picture:  \\

$\mid$ $\psi,t \rangle$ = $U(t,t_0) | \psi,t_0 \rangle$, \\

$\hat{H}(t)$ = $i \hbar \frac{\partial \mathcal U(t, t_0)}{\partial t} \mathcal U(t_0,t)$,  \\

U(t,0) = $\exp^{( - \frac{i \hat{\mathbf H} t} {\hbar} ) }$, for $\hat{H}$ $\text{time}$ independent\index{time independent}. \\

$\bullet$ (29) Heisenberg picture: \\

$\hat{A}_H(t)$ $\equiv$ $U^\dagger(t,0) \hat{A}_S U(t,0)$,  \\

$i \hbar \frac{d \hat{A}_H(t)}{dt}$ = $[\hat{A_\mathbf H}(t)]$,   \\

$\hat{H}_H(t)]$ + $i \hbar$ $\frac{\partial \hat{A}_H(t)}{dt}$.  \\

$\bullet$ (30) harmonic oscillator\index{harmonic oscillator}: \\

$\hat{H}$ = $\frac{1}{2m} \hat{p}^2$ + $\frac{1}{2}m \omega^2 \hat{x}^2$ = $ \hbar \omega(\hat{N}$ + $\frac{1}{2})$,  \\

$\hat{N}$ = $\hat{a}^\dagger \hat{a}$,  \\  

$\hat{a}$ = $\frac{1}{\sqrt{2}L_0}$ $(\hat{x}$ + $\frac{i \hat{p}}{m \omega})$,  \\

$\hat{a}^\dagger$ = $\frac{1}{\sqrt{2}L_0}(\hat{x} - \frac{i \hat{p}}{m \omega})$, $L^2_0$ = $\frac{\hbar}{m \omega}$,  \\

$\hat{x}$ = $\frac{L_0}{\sqrt{2}}(\hat{a} + \hat{a}^\dagger)$,  \\

$\hat{p}$ = $\frac{i}{\sqrt{2}} \frac{\hbar}{L_0}(\hat{a}^\dagger - \hat{a})$,  \\

$[\hat{a}, \hat{a}^\dagger]$ = 1,  \\

$[\hat{N}, \hat{a}]$ = $\hat{-a}$,  \\

$\hat{a} \varphi_0$ = 0, \\

$\varphi_0(x)$ = $N_0 \exp(- \frac{1}{2} \frac{x^2}{L^2_0})$, \\

$N^2_0$ = $\frac{1}{\sqrt{\pi}L_0}$,  \\

$\varphi_n$ = $\frac{1}{\sqrt{n!}} (a^\dagger)^n \varphi0 $, \\

$\hat{H} \varphi_n$ = $ \hbar \omega( n + \frac{1}{2}) \varphi_n$, \\

$ \hat{\mathbf N} \varphi_n$ = $n \varphi_n$,  \\

$\langle \varphi{m}, \varphi{n} \rangle$ = $\delta_{mn}$, \\ 

$\hat{a}^\dagger \varphi_n$ = $\sqrt{n + 1} \varphi_n + 1$,  \\

$\hat{a} \varphi_n$ = $\sqrt{n} \varphi_{n - 1}$, \\

$\hat{x}_H(t)$ = $\hat{x} \cos \omega{t} + \frac{\hat{p}}{m \omega} \sin{\omega{t}}$,  \\  

$\hat{p}_H(t)$ = $\hat{p} \cos \omega{t} - m \omega \hat{x} \sin \omega{t}$.  \\

$\bullet$ (31) coherent-states\index{coherent-states}:\\

$\mid \alpha \rangle$ $\equiv$ $e^{\alpha{a}^\dagger - \alpha^*{a}} \mid 0 \rangle$,  \\

\begin{equation}
\alpha = \frac{\langle \hat{x} \rangle} {\sqrt2{L_0}} + i \frac{\langle \hat{p} \rangle L_0} {\sqrt2 \hbar} \in \mathbb{C}
\end{equation}

$\hat{a} \mid \alpha \rangle$ = $\alpha \mid \alpha \rangle$,  \\

$\mid \alpha \rangle$ = $e^{-{\frac{1}{2} \mid \alpha \mid^2}}$ $e^\alpha$ $a^\dagger$ $\mid 0 \rangle$,  \\

$\mid \alpha, t \rangle$ = $e^{-i \omega \frac{t}{2}}$ $\mid e^{-i \omega t} \alpha \rangle$.  \\

$\bullet$ See \eqref{eq:orbital-angular-momentum} for the quantization of $\hat{L_z}$, and $\hat{L^2}$ for angular momentum Eigenvalues\index{Eigenvalues}. \quad (32) orbital-angular-momentum\index{orbital-angular-momentum} \\
\begin{equation}
\hat{L}_x = \hat{y} \hat{p}_z - \hat{z} \hat{p}_y, \quad \hat{L}_y = \hat{z} \hat{p}_x - \hat{x} \hat{p}_z, \quad \hat{L}_z = \hat{x} \hat{p}_y - \hat{y} \hat{p}_x
\end{equation}
\label{eq:orbital-angular-momentum}
\noindent

$[\hat{L}_x, \hat{L}_y]$ = $i \hbar \hat{L}_z$, \\

$[\hat{L}_y, \hat{L}_z]$ = $i \hbar \hat{L}_x$, \\

$[\hat{L}_z, \hat{L}_x]$ = $i \hbar \hat{L}_y$,  \\

$\hat{L}^2$ $\equiv$ $\hat{L}_x \hat{L}_x + \hat{L}_y \hat{L}_y + \hat{L}_z \hat{L}_z$, $[\hat{L}^2$, $\hat{L}_i] = 0$. \\ 

$\nabla^2$ = $\frac{\mathbf 1} {r}$ $\frac{\partial^2}{\partial \mathbf r^2}$ $\mathbf r$ - $\frac{\mathbf 1} {\mathbf r^2}$ $\frac{\hat L^2} {\hbar^2}$,  \\

$\hat{L}^2$ = $- \hbar^2$ $(\frac{\partial^2}{\partial \theta^2}$ + $\cot$ $\theta$ $\frac{\partial}{\partial \theta}$ + $\frac{1}{\sin^2 \theta}$ $\frac{\partial^2}{\partial \phi^2})$,  \\

$\hat{L}_z$ = $\frac{\hbar}{\mathbb i}$ $\frac{\partial}{\partial \theta}$,  \\

$\hat{L}_\pm$ = $\hbar e^{\pm i \phi}$ $(\pm\frac{\partial}{\partial \theta}$ + i $\cot \theta$ $\frac{\partial}{\partial \theta})$.  \\

$Y_{\ell m}$ $(\theta, \phi)$ = $\langle \theta \phi$ $\mid$  ${\ell m}$ $\rangle$ = $Y_{\ell m}$ $(\Omega)$ $\equiv$ $\mathcal N_{\ell m}$ $P^m_\ell$ $(\cos$ $\theta)$ $e^{i m \theta}$.  \\

$L_z Y_{\ell m}$ = $\hbar$ $\mathbf m$ $\mathbf Y_{\ell m}$, \\

$\mathbf L^\angle$ $Y_{\ell m}$ = $\hbar^\angle$ $\ell$ ( $\ell + 1 )$ $Y_{\ell m}$, \\

$\int d \Omega \mathbf Y^*_{\ell' m'} (\Omega)$ = $\delta_{\ell' \ell} \delta_{m' m}$,  \\

$\sum_{\ell = 0}^\infty \sum_{m = - \ell}^\ell Y^*_{\ell m} (\Omega)$ = $\delta$ $(\cos\theta - \cos\theta')$ $\delta$ $(\phi - \phi')$.  \\

$Y_{0,0}(\theta, \phi)$ = $\frac{1} {\sqrt{4 \pi}}$,  \\

$Y_1$, $\pm 1 (\theta, \phi)$ = $\mp {\sqrt \frac{3}{8 \pi}}$ $\sin \theta$ $e^{\pm i \phi}$, \\

$Y_{1,0}$ $(\theta, \phi)$ = ${\sqrt \frac{3}{4 \pi}}$ $\cos$ $\theta$. \\

$\bullet$ (33) central potentials\index{central potentials}:  \\

$\psi(r, \theta, \phi)$ = $\frac{u(r)}{r}$ $Y_{\ell m}(\theta, \phi)$,  \\  

$ (- {\frac{\hbar^2}{2m}}$ $ \frac{d^2}{dr^2}$ + V(r) + $\frac{\hbar^2 \ell (\ell + 1)} {2mr^2}) u(r) $ = u(r), u(r) $\approx$ $r^{\ell + 1}$, as r $\longrightarrow$ 0.  \\

$\bullet$ (34) Spin one-half\index{Spin one-half} \\

$\hat{H}$ = $- \mu \cdot B$, $\hat{\mu}$ = $ g \frac{e \hbar}{2mc}$ $\frac{1}{\hbar} \hat{S}$,  \\  

$\mu_B$ = $\frac{e \hbar}{2m_eC}$, $ \hat{\mu_e}$ = -2 $ \mu_B$ $\frac{\hat{\mathbf S}}{\hbar}$.  \\

$\mid 1 \rangle$ $\equiv$ $\mid Z; + \rangle$ = $\mid + \rangle$ = $\begin{pmatrix}1\cr0\cr\end{pmatrix}$,  \\ 

$\mid 2 \rangle$ $\equiv$ $\mid Z; - \rangle$ = $\mid - \rangle$ = $\begin{pmatrix}0\cr1\cr\end{pmatrix}$, \\ 

$\hat{S}_i$ = $\frac{\hbar}{2} \sigma_i$, $\sigma_x$ = $\begin{pmatrix}0 & 1\cr 1 & 0\cr\end{pmatrix}$, \\

$\sigma_y$ = $\begin{pmatrix}0 & -i\cr i & 0\cr\end{pmatrix}$, $\sigma_z$ = $\begin{pmatrix}1 & 0\cr 0 & -1\cr\end{pmatrix}$  \\

$[\sigma_i, \sigma_j]$ = $2i \epsilon_{ijk} \sigma_k$,  \\ 

$[\hat{S_i}, \hat{S_j}]$ = $i \hbar \epsilon_{ijk} \hat{S_k}$,  \\  

$\sigma_i \sigma_j$ = $\delta_{ij} \mathbf{1} + i \epsilon_{ijk} \sigma_k$, $(\sigma \cdot \mathbf{a} ) (\sigma \cdot \mathbf{b} )$ = $\mathbf{a} \cdot \mathbf{b} \mathbf{1} + i \sigma \cdot (\mathbf{a} \times \mathbf{b})$, \\

$e^{i \mathbf a \cdot \sigma}$ = $\mathbf{1} \cos \mathbf a + i {\mathbf \sigma}  \cdot  \hat{\mathbf a} \sin a$,  \\

$ \mathbf a$ = $\hat{\mathbf a} a$, a = $|\mathbf a|$.  \\  

$\hat{S_n}$ $\equiv$ $\mathbf n \cdot \hat{S}$ = $\frac{\hbar}{2} \mathbf n \cdot \sigma$, $\hat{S_n} \mid \mathbf n; \pm \rangle$ =  $\pm \frac{\hbar}{2} \mid \mathbf n; \pm \rangle$, \\ 

$\mid \mathbf n \rangle$ $\equiv$ $\mid \mathbf n ; + \rangle$,  \\

 $\langle \mathbf n \mid  \hat{\mathbf S} \mid \mathbf n \rangle$ = $\frac{\hbar}{2} \mathbf n$,  \\

$\mid \mathbf n; _ \rangle$ = $\cos \frac{\theta}{2} \mid + \rangle + \sin \frac{\theta}{2} \mathbf{e}^{\mathbf{i} \theta} \mid - \rangle$, \\

$\mid \mathbf n; - \rangle$ = $- \sin \frac{\theta}{2} \mathbf{e}^- {\mathbf{i} \theta} \mid + \rangle$ + $\cos \frac{\theta}{2} \mid - \rangle$, \\

$\hat{R_n} (\alpha)$ $\equiv$ $\exp(- \frac{i}{\hbar} \alpha \hat{S_n})$, \\ 

$\hat{R_n} (\alpha) \mid \mathbf n' \rangle$ = $\mid n'' \rangle$, with $ \mathbf n''$ = $\mathcal{R_\mathbf n} (\alpha) \mathbf n'$ \\

$\hat{R_n}(\alpha) \hat{\mathbf S} \hat{R_n}(\alpha)$ = ${\mathcal{R}}$ ($\alpha$) $\hat{\mathbf{S}}$ \\

$\bullet$ (35) Spin precession\index{Spin precession}: $ \hat{H}$ = $ {\mathbf \omega_L} \cdot \hat{S}$,  \\

$ {\mathbf \omega_L}$ = $ - \gamma{\mathbf B}$, $ \hat{\mathbf \mu}$ = $ \gamma \hat{\mathbf S}$.  \\

$\bullet$ (36) general angular momentum\index{general angular momentum}: \\

$ [\hat{J_i}, \hat{J_j}]$ = $i \hbar \epsilon_{ijk} \hat{\mathbf J_k}$ $\longleftrightarrow$ $\hat{\mathbf J} \times \hat{\mathbf J}$ = $ i \hbar \hat{\mathbf J}$,  \\

$[\hat{\mathbf J^2},\hat{\mathbf J_i}] $ = 0.  \\  

$ \hat{\mathbf J_ \pm}$ = $\hat{\mathbf J_x} \pm i \hat{\mathbf J_y}$, $ \hat{(J_ \pm)^\dagger}$ = $\hat{\mathbf J_\mp}$, \\

$\hat{\mathbf J_x}$ = $\frac{1}{2} (\hat{\mathbf J_+} + \hat{\mathbf J_-})$, $\hat{\mathbf J_y}$ = $\frac{1}{2i}(\hat{\mathbf J_+} - \hat{\mathbf J_-})$,  \\  

$[\hat{\mathbf J_z}, \hat{\mathbf J_ \pm}]$ = $ \pm \hbar \hat{\mathbf J_ \pm}$, $[\hat{\mathbf J_+}, \hat{\mathbf J_-}]$ = $2 \hbar \hat{\mathbf J_z}$,  \\

$[\hat{\mathbf J^2}, \hat{\mathbf J_\pm}]$ = 0,  \\  

$\hat{\mathbf J^2}$ = $\hat{\mathbf J_+} \hat{\mathbf J_-}$ + $\hat{\mathbf J^2_z}$ - $\hbar$ $\hat{\mathbf J_z}$ = $\hat{\mathbf J_-} \hat{\mathbf J_+}$ + $\hat{\mathbf J^2_z}$ + $\hbar$ $\hat{\mathbf J_z}$.  \\

$\hat{\mathbf J^2} |{\mathbf {jm}} \rangle$ = $\hbar^2 \mathbf j(\mathbf j + 1) | \mathbf {jm} \rangle$ = $ \hbar \mathbf m | \mathbf {jm} \rangle$,  \\

$\mathbf m$ = $ - \mathbf j,...,\mathbf j$  \\

$\hat{\mathbf j_\pm} | {\mathbf jm} \rangle$ = $\hbar$ $\sqrt{{{\mathbf j} {(\mathbf j + 1)}} - {\mathbf m}{(\mathbf m \pm 1)}}$ $ | \mathbf j, \mathbf m \pm 1 \rangle$. \\

$\hat{R^\dagger_n}(\alpha) \hat{\mathbf J} \hat{R_n}(\alpha)$ = $\mathcal R_n (\alpha) \hat{\mathbf J}$,  \\ 

$\hat{\mathbf R_n}(\alpha)$ = $e^{ - i \frac{\alpha}{\hbar} {\mathbf n} \cdot {\mathbf J}}$. \\

$\bullet$ (37) addition of angular momentum\index{addition of angular momentum}:  \\ 

$\hat{\mathbf J}$ = $\hat{\mathbf J_1} + \hat{\mathbf J_2}$.  \\

$\bullet$ (38) uncoupled basis: $ | j_1 j_2; \mathbf m_1 \mathbf m_2 \rangle$ $\text{CSCO:}$ $(\hat{\mathbf J^2_1}, \hat{\mathbf J^2_2},
\hat J_{1z}, \hat{j_{2z}})$  \\

$\bullet$ (39) coupled basis\index{coupled basis}:  \\

$\mid {j_1} {j_2 ;} \mathbf{jm} \rangle$ CSCO\index{CSCO}: \\

$(\hat{\mathbf J^2_1}, \hat{J^2_2}, \hat{\mathbf j^2}, \hat{j_z})$.  \\

$j_1 \otimes j_2$ = $(j_1 + j_2) \oplus  (j_1 + j_2 - 1)\oplus...\oplus \mid j_1 - j_2 \rangle$.  \\  

\begin{equation}
\mid  \mathbf j_1 \mathbf j_2; \mathbf{jm} \rangle = \sum_{\mathbf m_1 + \mathbf m_2 = \mathbf m} \mid {\mathbf j_1}{\mathbf j_2}; {\mathbf m_1}{\mathbf m_2} \rangle \underbrace{\langle{\mathbf j_1 \mathbf j_2}; {\mathbf m_1 \mathbf m_2} \mid {\mathbf j_1 \mathbf j_2}; \mathbf{jm} \rangle}_{\text{Clebsh-Gordan- coefficient\index{Clebsch-Gordan-coefficient}}}
\end{equation}

$\hat{\mathbf J_1} \cdot \hat{\mathbf J_2}$ = $\frac{1}{2}$ $( \hat{\mathit J_{1 +}} \hat{\mathit J_{2 -}} + \hat{\mathit J_{1 -}} \hat{\mathit J_{2 +}})$ + $ \hat{\mathit J_{1z}}$ $\hat{\mathit J_{2z}}$ = $ \frac{1}{2}$ $(\hat{\mathbf J^2} - \hat{\mathbf J^2_1} - \hat{\mathbf J^2_2})$. \\

$\bullet$ (40) Hydrogen-atom:  \\ 

$\hat{H}$ = $ \frac{\hat{\mathbf p^2}} {2 {\mathbf m}}$ - $ \frac{\mathit Z e^2}{r}$, $ \mathit Z$ = 1 $\text{for hydrogen.}$ \\

$ \mathit E_n$ = $ - \frac{\mathit {Z^2 e^2}}{2a_0} \frac{1}{\mathbf n^2}$, $a_0$ = $ \frac{\hbar^2}{me^2}$ $ \simeq$ 52.9 $\pm$,  \\

$\frac{e^2}{2a_0}$ = $ \frac{1}{2} \mathbf {mc^2} \alpha^2$ $ \equiv$ Ry $ \simeq$ 13.6 eV. \\

$ \psi_{n \ell m} (\mathbf X)$ = $ \mathcal N ( \frac{r}{a_0})^\ell$ $ ( polynomial \in {\frac{r}{a_0}} of degree [{\mathbf {n} - (\ell + 1)}])$ $ e^{- \frac{\mathit Zr} {\mathbf {n a_0}}}$ $\mathbf Y_{\ell m}$ $(\theta, \phi)$, \\ 

$\psi_{100} ( \mathbf r, \theta, \phi)$ = ${\sqrt{\frac{\mathit{Z^3}}{\pi {a^3_0}}}} e^{- \frac{ \mathit{Z}}{a_0}}$ \\

$\mathit Z$ = 1: $ \langle r \rangle$ = $ \frac{1}{2} a_0(3n^2 - \ell(\ell + 1))$, $ \langle \frac{1}{r} \rangle$ = $\frac{1}{a_0 n^2}$,  \\

$\langle \frac{1}{r^2} \rangle$ = $ \frac{1}{a^2_0 n^3 (\ell + \frac{1}{2})}$, \\

$ \langle \frac{1}{r^3} \rangle$ = $ \frac{1}{\mathbf {a^3 n^3} \ell (\ell + \frac{1}{2}) (\ell + 1)}$.  \\

$\bullet$ (41) Fine structure corrections\index{Fine structure corrections} of the Hydrogen atom:  \\

$\mathit E_{n \ell {\mathbf j m_j}}$ = $ - \frac{e^2}{2 \mathbf a_0}$ $ \frac{1}{\mathbf n^2}$ $ (1 + \frac{\alpha^2}{\mathbf n^2} [\frac{\mathbf n}{\mathbf j + \frac{1}{2}} - \frac{3}{4}])$.  \\

$\bullet$ (42) Density matrix: \\ 
 
$\mathit E$ = $[ (\mathbf p_1, \mid \psi_1 \rangle), ..., ( \mathit p_n, \mid \psi_n \rangle)]$,  
$\mathit p_1,...,\mathit p_n$ > 0, \\ 

$p_1 +...+ p_n$ = 1,  \\ 
   
$\rho_\mathit E$ $\equiv$ $\sum_{a = 1}^n$ $\mathit p_a \mid \psi_a \rangle \langle \psi_a \mid$, 
$\langle \hat{\mathit Q \rangle_\mathit E}$ = tr $ (\mathit Q_{\rho_{\mathit E}})$.   
General $\rho$ is positive semi-definite, and $\mathbf{tr}(\rho)$ = 1. \\

Pure state $\longleftrightarrow$ $\mathbf \mathbf{tr}(\rho^2)$ = 1.  \\
   
Spin-one-half-density-matrix\index{Spin-one-half-density-matrix}: $\rho$ = $\frac{1}{2} ( \mathbf{1} + \mathbf a \cdot \sigma)$,  \\

$\mid \mathbf a \mid$ $\leq$ 1.  \\

$\bullet$ (43) Time evolution\index{time evolution}:  \\

$i \hbar \frac{\partial \rho}{\partial t}$ = $ [ \hat{H}, \rho]$.  \\

$\bullet$ (44) Schmidt decomposition\index{Schmidt decomposition}: \\ 

$\mid  \psi AB \rangle$ = $ \sum_{k = 1}^r$ $\sqrt {\mathit p_k}$ $\mid \mathbf k_A \rangle$ $\otimes$ $\mid  \mathit k_B \rangle$,  \\

$r \leq$ $\mathbf d_A$ $\leq$ $\mathbf d_B$, \\

$\rho_A$ = $ \sum_{k =1}^r$ $ \mathit p_k \mid \mathit k_A \rangle \langle \mathit k_A \mid$,  \\  

$\rho_B$ = $ \sum_{k =1}^r$ $ \mathit p_k \mid \mathit k_B  \rangle \langle \mathbf k_B \mid$,  \\  

$\langle k_A \mid k'_A \rangle$ = $\delta_{k,k'}$, \\  

$\langle k_B \mid k'_B \rangle$ = $\delta_{k,k'}$  \\

$\bullet$ (45) Lindblad equation\index{Lindblad equation}:  \\
\begin{equation}
\frac{\partial \rho} {\partial t} = \frac{1} {i \hbar} [ \mathbf H, \rho ] +  \sum_k ( \mathbf L_k \rho \mathbf L^\dagger_k - \frac{1} {2}[ \mathbf L^\dagger_k \mathbf L_k, \rho])
\end{equation}

$\bullet$ (46) electromagnetic couplings\index{electromagnetic couplings}:  \\

$\hat{H}$ = $\frac{1}{2m}(\hat{\mathbf p}) - \frac{q}{c}$ $\mathbf A$ $(\hat{x}, t)^2$ + q $\phi(\hat{\mathbf x, t})$ (no spin).  \\ 

$\bullet$ (47) Go to equation \eqref{eq:Gauge-transformation} \quad Gauge-transformations\index{Gauge-transformations}:  \\

$\mathbf A'$ = $\mathbf A + \nabla \wedge$,   \\  

$\phi'$ = $\phi$ - $\frac{1}{c}$ $ \frac{\partial \wedge}{\partial {\mathbf t}}$, \\  

$\Psi'$ = $\exp(i \frac{q \wedge}{\hbar c})$ $\Psi$.  \\

$\bullet$ (48) Pauli Hamiltonian (electron\index{electron}):  \\

$\hat{\mathbf H}_{Pauli}$ = $ \frac{1}{2m_e} (\hat{\mathbf p} + \frac{e}{c} \mathbf A)^2$ + $\frac{e \hbar}{2m_e c}$ $\sigma$ $\cdot$ $\mathbf B$ - e $\phi$ $(\hat x,t)$. \\

$\bullet$ (49) Time-independent Perturbation:  non-degenerate:  \\

$\mid n \rangle_\lambda$ = $\mid n^{(0)} \rangle - \lambda$ $\sum_{k \neq n} \frac{\delta \mathbf H_{kn}} {\mathbf E^{(0)}_k - \mathbf E^{(0)}_n}$ $\mid k^{(0)} \rangle$ + $\mathcal O (\lambda^2)$,  \\

$E_n (\lambda)$ = $ E^{(0)}_n$ + $ \lambda \delta \mathbf H_{nm}$ - $\lambda^2$ $\sum_{k \neq n} \frac{ | \delta \mathbf H_{kn} |^2} {\mathbf E^{(0)}_k - \mathbf E^{(0)}_n}$ + $\mathcal O (\lambda^3)$. \\ 

 $\bullet$ (50) Degeneracy\index{Degeneracy} lifted at $\mathcal O (\lambda)$: good-basis\index{good-basis}: \\

$\delta \mathbf H_{\mathcal {ij}}$ $\equiv$ $\langle \Psi^{(0)}_\mathcal I \mid \delta \mathbf H \mid \Psi^{(0)}_\mathcal J \rangle$ = $\mathbf E^{(1)}_\mathcal {nJ}$ $\delta_\mathcal {IJ}$,  \\

$\mid \Psi_\mathcal I \rangle_\lambda$ = $\mid \Psi^{(0)}_{\mathcal I} \rangle - \lambda \sum_p \frac{\delta \mathbf H_{p \mathcal I}} {\mathbf E^{(0)}_p - \mathbf E^{(0)}_n} \mid \mathbf p^{(0)} \rangle + \sum_{K \neq \mathcal I} \frac{ \mid \Psi^{(0)}_K \rangle} {\mathbf E^{(1)}_{n \mathcal I} - \mathbf E^{(0)}_{n \mathcal K}} \sum_p \frac{ \delta \mathbf H_\mathcal K p \delta \mathbf H_pl} {\mathbf E^(0)_p - \mathbf E^(0)_n}$ + $\mathcal O \lambda^2$,  \\

$\mathbf E_{n \mathcal I} (\lambda)$ = $\mathbf E^{(0)}_n + \lambda \delta \mathbf H_{\mathcal II} -\lambda^2 \sum_p \frac{ \mid \delta \mathbf H_p \mathcal I \mid^2}{\mathbf E^{(0)}_p - \mathbf E^{(0)}_n}$ + $\mathcal O (\lambda^3)$. \\

$\bullet$ (51) WKB-quantization\index{WKB-quantization} with (a, b) turning-points\index{turning-points} or hard walls:  \\

$\int_a^b k(x') dx$ = $(\mathbf n + \beta) \pi$,  \\

$\mathbf n$ = $0,1,2,...,$  \\

$\beta$ = $1 - \frac{1}{4}$ (Number of turning points). \\

$\bullet$ (52) Time-dependent perturbations\index{Time-dependent perturbations}:

$\hat{\mathbf H} (t)$ = $\hat{\mathbf H^{(0)}} + \delta \mathbf H (t)$, $ | \Psi (t) \rangle$ $\equiv$ $e^{i \hat{\mathbf H}^{(0)} {\frac{t}{\hbar}}} \mid \Psi (t) \rangle$,  \\

$i \hbar \frac{ \mathbf d}{\mathbf dt} \mid {\tilde \Psi} (t) \rangle$ = $ {\tilde \delta} \mathbf H (t) \mid \Psi (t) \rangle$,  \\

$\delta \mathbf H (t)$ $\equiv$ $e^{i \hat{\mathbf H}^{(0)}{\frac{t}{\hbar}}} \delta \mathbf H (t) e^{- i \hat{\mathbf H^{(0)}} {\frac{t}{\hbar}}}$,  \\

$\mid \tilde \Psi (t) \rangle$ = $ \mid \Psi(0) \rangle + \int_0^t \frac{ \delta \mathbf H(t')} {i \hbar} \mid \Psi(0) \rangle \mathbf dt' + \mathcal O( \delta \mathbf H^2)$  \\

$\bullet$ (53) Fermis'-golden-rule\index{Fermis'-golden-rule}:  \\

$\hat{\mathbf H}$ = $\hat{\mathbf H^{(0)}}$ + V,

$\hat{\mathbf H}$ = $\hat{\mathbf H^{(0)}}$ + $2 \mathbf H' \cos \omega t$,  \\

w = $\frac{2 \pi}{\hbar} \mid \mathbf V_f \mid^2 \rho(\mathbf E_f)$, \\

w = $\frac{2 \pi}{\hbar} \mid \mathbf H'_f \mid 2 \rho(\mathbf E_f)$  \\

The constant transition rate of Energy to a final state.  \\

%$\hat{\mathbf{\textbf{H}}}$ = $\hat{\mathbf{\textbf{H}^{(0)}}}$ + $V$, \\

%$\hat{\mathbf{\textbf{H}}}$ = $\hat{\mathbf{\textbf{H}^{(0)}}}$ + $2$ $\mathbf{\textbf{H}^{\prime}}$ $\cos$ $\omega$ $t$,  \\

%$w$ = $\frac{2 \pi}{\hbar}$ $\mid$ $\textbf{V_f}}$ $\mid^2$ $\rho$ ($\textbf{E_f}}$),   \\

%$w$ = $\frac{2 \pi}{\hbar}$ $\mid$ $\mathbf{\textbf{H}^{\prime}_f}$ $\mid$ $2$ $\rho$ ($\m%athbf{\textbf{E_f})$    \\

$\bullet$ (54) Adiabatic-approximation\index{Adiabatic-approximation}: \\ 

$\mid \Psi(t) \rangle$ $\simeq$ $e^{i \theta_k(t)} e^{i \gamma_k(t)} \mid \Psi_k(t) \rangle$,  \\  

$\hat{\mathbf H} (t) \mid \Psi_k(t) \rangle$ = $\mathbf E_k(t) \mid \Psi_k(t) \rangle$, \\   

$\theta_k(t)$ = $ -\frac{1}{\hbar} \int_0^t \mathbf E_k(\mathbf t') {\mathbf dt'}$, \\ 

$\gamma_k(t)$ = $\int_0^t \nu_k(t') {\mathbf dt'}$, \\ 

$\nu_k(t)$ = $i \langle  \Psi_k(\mathbf t) \mid {\dot \Psi_k}(t) \rangle$. \\

$\bullet$ (55) Berrys' phase\index{Berrys' phase}:  \\

$\gamma_n(\Gamma_{if})$ = $\int_{\Gamma if} \mathbf A_n(\mathbf R) \cdot  d{\mathbf R}$, $ \mathbf A_n(\mathbf R)$ $\equiv$ $ i \langle \Psi_n(\mathbf R) \mid \nabla_{\mathbf R} \mid \Psi_n(\mathbf R) \rangle$. \\


$\bullet$ (56) Scattering-on-the-half-line\index{scattering-on-the-half-line}:  \\

$\Psi(\mathbf x)$ = $e^{1 \delta} \sin(\mathbf {kx} + \delta)$,  \\   

$\Psi_s(\mathbf x)$ = $e^{i \delta} \sin \delta e^{ikx}$,  \\  

$\mathbf x$ > $ \mathbf R$  \\

$\bullet$ (57) Time-delay\index{Time-delay}:  \\  

$\vartriangle t$ = $2\hbar \delta'( \mathbf E_0)$.  \\


$\bullet$ (58) Levinsons' Theorem:  \\  

$\mathit N_b$ = $\frac{1}{\pi} ((\delta)(0) - \delta(\infty))$.  \\


$\bullet$ (59) Resonance\index{Resonance}, Breit-Wigner:  \\  

$\mid \Psi_s \mid^2$ $\simeq$ $\frac{ \frac{1}{4} \Gamma^2} {(\mathbf E - \mathbf E_\alpha)^2 + \frac{1}{4} \Gamma^2}$.  \\

Energy dependence of the squared scattering amplitude $ \mid \Psi_s \mid^2$ near resonance.  \\

$\bullet$ (60) Scattering in 3D:  \\  
\begin{equation}
\Psi(\mathbf r) = \varphi(\mathbf r) + \Psi_s(\mathbf r) \simeq e^{ikz} + f_k(\theta, \phi) \frac{e^{ikr}}{\mathbf r}, \mathbf r \gg a.
\end{equation}

$\frac{d \sigma}{d \Omega}$ = $\mid f_k(\theta, \phi) \mid^2$,  \\ 

$\sigma$ = $\int \mid f_k(\theta, \phi) \mid^2 d \Omega$.  \\

$\bullet$ (61) Rayleigh:  \\

$e^{ikz}$ = ${\sqrt 4 \pi} \sum_{\ell = 0}^\infty {\sqrt{2 \ell + 1}}$ $i^\ell \mathbf Y_{\ell 0}(\theta) j_\ell(kr)$.  \\

$\bullet$ (62) Phase shifts\index{Phase shifts}: \\ 

$f_k(\theta)$ = $\frac{\sqrt(4 \pi)} {\mathbf k} \sum_{l = 0}^\infty \sqrt{2l + 1} \mathbf Y_{l 0}( \theta) e^{i \delta_l} \sin \delta_l$, \\

$\sigma$ = $\frac{4 \pi}{\mathbf k^2} \sum_{l = 0}^\infty(2 l + 1) \sin^2 \delta_{l}$.  \\

$\Psi \mathbf r \mid_l$ = $\mathbf A_{l} \mathbf j_{l} \mathbf(kr) + \mathbf B_{l} \mathbf n_{l} \mathbf{kr} \mathbf Y_{l 0} (\theta)$,  \\  

$\mathbf r > \mathbf a$,  \\

$\tan \delta_{l}$ $\equiv$  - $\frac{\mathbf B_{(l)}}{\mathbf A_{(l)}}$.  \\

\chapter{Stationary-State\index{Stationary-State}}

There is time-dependence on a stationary-state. If we separate the variables of the wave-function $\Psi(x,t)$, to $\psi(x)$, and $\phi(t)$, we have a constant on the (t) component: $\psi(x)e^{\frac{-i E}{\hbar}t}$; The probability density of the wave-function(or the magnitude squared), complex conjugated, guarantees oscillation in establishing a \emph{time-independent wave-function:} $\Big| \Psi(x,t) \Big|^2$ = $\psi^{*}(x) e^{ + \frac{i E}{\hbar}t}$ \quad $\psi(x) e^{ - \frac{i E}{\hbar}t}$ = $e^{0} (\psi^{*}(x) \psi(x))$ = $\Big| \psi(x) \Big|^2$, where $\mathbf{e^{0}}$ = $\mathbf{1}$.  \\
$\therefore$ with an index (n) on the $\psi$'s, and (E), we can obtain the ($n^{th}$) (E) of the time-independent Schr$\ddot{o}$dinger wave-function.  \\
Stationary, and static aren't the same! \\

In general, $\Psi(x,t)$ = $\sum_{n = 1}^\infty$ $C_n$ $\Psi(x,t)$ = $\sum_{n = 1}^\infty$ $\psi(x) e^{\frac{i E}{\hbar}t}$, with constant ($C_n$ = $a_n + i b_n$). The evolution of the wave-function is Linear, and collapses, when there's a measurement.  \\
At $t = 0$, the \emph{Initial-State}\index{Initial-State}, $\psi(x,0)$ = $\sum_{n = 1}^\infty$ $C_n$ $\psi_{n} (x)$.  \\
Any $t^{'}$, can be approximated, with this set of wave-function.







\chapter{The Simple Harmonic Oscillator}


When Energy is quantized, the only allowed values for
E = $\epsilon \hbar \omega$, are values that yield solutions in the physical Hilbert space: $E_n$ = $(n + \frac{1}{2}) \hbar \omega$, $(n) = 0, 1, 2,...$.  \\
For each value of (n), the corresponding polynomials is of ($n_{th}$) order, called Hermite polynomials. We know that: $\boxed{[x, p] = i \hbar}$, and that the Hamiltonian $\hat H$ = $\frac{p^2}{2m}$ $\frac{1}{2}$ $m$ $\omega^2 \hat x^2$.The Quantum Harmonic Oscillator is an extension of the Classical Oscillator. We expect that Energy eigenstates, only exist for certain quantized values of (E). This is a differential equation:
\begin{equation}
- \frac{\hbar^2}{2m}\frac{d^2 \psi (x)}{d x^2} + \frac{1}{2} m \omega^2 x^2 \psi (x) = E \psi (x)
\end{equation}

The length scale $L_0$ = $\sqrt{\frac{\hbar}{m \omega}}$, to be called the $\mathbf{\textbf{oscillator size\index{oscillator size}}}$, is as relevant to the harmonic oscillator as the Bohr radius ($a_0$), is to the Hydrogen atom. The ground state energy of the quantum oscillator, is ($\frac{\hbar \omega}{2}$), and ($L_0$), is precisely the Classical oscillation amplitude of an oscillator with that Energy. A recursion relation of ($\epsilon$ = $2n + 1$) must terminate, for us to get an Energy eigenstate. The corresponding physical Energies $E_n$ = $\frac{\hbar \omega}{2}$ $(2n + 1)$ = $\boxed{\hbar \omega$ $(n + \frac{1}{2})$, $n = 0,1,2,...}$.The spectrum of this result is non-degenerate, and all energy eigenstates of the harmonic oscillator are bound states of a one-dimensional potential. The polynomial solutions are normalized Hermite polynomials and satisfy $2^{nd}$ order differential equation. We can write the eigenstate as a function of the physical coordinate (x), from Hermite polynomials. \emph{The Energy eigenstates are real functions of (x)}. We have full orthonormality of the Energy eigenstates $\langle \varphi_n, \varphi \rangle$ = $\delta_{mn}$, from the normalization
constant ($N_n$), determined from the condition $\langle \varphi_n, \varphi_n \rangle$ = $\int_{- \infty}^\infty$ dx $\varphi_n$ (x) = $1$. ($\hat x)$  and ($\hat p$) are Hermitian operators. All Energy eigenstates satisfy Energy (E),  so E $\geq$ $\frac{1}{2}$ $\hbar \omega$, since $\langle \hat H \rangle_\psi$ = $E$ $\hat p$, is $\frac{\hbar}{i}$ $\frac{\partial}{\partial x}$ in \emph{Position} space, and $\hat x$ is $i \hbar$ $\frac{\partial}{\partial p}$ in \emph{Momentum} space. The normalized ($n_th$) excited state of the simple harmonic oscillator is:
\begin{equation}
\varphi_n \equiv \frac{1}{\sqrt N!} \underbrace{\hat a^{+}...\hat a^{+}}_\text{n} \psi_0 = \frac{1}{\sqrt N!} (\hat a^{+})^n \psi_0
\end{equation}
Since $\hat N \varphi_n$ = $\varphi_n$, means ($\hat N$) counts the number of creation operators acting on the ground state, and its' Energy is given by:
\begin{equation}
E_n = \hbar \omega (n + \frac{1}{2})
\end{equation}
$\therefore$ $(\vartriangle x)^2_n$ = $L^2_0$ $(n + \frac{1}{2})$ $\approx$ $L_0 \sqrt{n}$, comparable to the amplitude of the Classical oscillation of a particle with Energy $\hbar \omega n$. Only for the ground state $\phi_0$, does the product of uncertainties saturate the lower bound given by the Heisenberg uncertainty inequality:
\begin{equation}
(\vartriangle(x), \vartriangle (p)) \geq \frac{\hbar}{2}
\end{equation}

\chapter{Sample Practice}

\section{More Useful Equations}

$\bullet$  d'Alemberts' Principle $\longrightarrow$  $\sum_i \langle \mathbf {F}_i(\mathbf c(t),t) - \mathbf {m}_i {c_i(t)}, \mathbf {v}_{i} \rangle$ = 0 $\forall$ tangent vectors $\mathbf v$ at $\mathbf {M}_{\mathbf c(t)}$  \\

$\bullet$  $n^{th}$ Taylor Polynomial about x = $x_0$ for f,  \\

$p_n(x)$ = $f(x_0) + f'(x_0)(x - x_0) + \frac{f''(x_0)}{2!}(x - x_0)^2 +...\frac{f^{(n)}{(x_0)}}{n!}(x - x_0)^n $,  \\    
$\sum_{k=0}^\infty$ $\frac{ ({-i} * {\frac{\theta}{2} * {X} )^k}} {k!}$  \\

$\bullet$ Airy equation $\longrightarrow$ $\frac{d^2 \psi} {dU^2}$ - $u \psi (u) = 0$  \\

$\bullet$ Jacobi identity $\longrightarrow$ $[X,[Y,Z]] + [Y,[X,Z]] + [Z, [X,Y]] = 0$  \\

$\bullet$ Stokes theorem(Line-Integrals\index{Line-Integrals})  \\

$\iint_s curl \mathbf F \cdot d \mathbf S$ = $\int_c \mathbf F \cdot d \mathbf r$, where C is a closed curve, and $\mathbf F$ is a vector field defined on C. \\

$\bullet$ A sympletic-manifold\index{sympletic-manifold} (N, $\omega$), is quantizable (for a particular value of $\hbar$), if
$\frac{1}{{2 \pi \hbar}}$ $\int_s \omega \in \mathbb Z$, for every closed surface S in N.  \\

$\bullet$ Probability-density-function(Gaussian-distribution\index{Gaussian-distribution}):  \\
\begin{equation}
f(x) = \frac{1}{\sigma \sqrt{2 \pi}} e^{- \frac{(x - \mu)^2}{2 \sigma^2}}
\end{equation}

$\bullet$ Poisson bracket of f and g;  Moyal bracket $\longrightarrow$ [f,g] = \\

$\frac{\partial f}{\partial q_1}$ $\frac{\partial g}{\partial p_1}$ +...+ $\frac{\partial f}{\partial q_n}$ $\frac{\partial g}{\partial p_n}$ - $\frac{\partial f}{\partial p_1}$ $\frac{\partial g}{\partial q_1}$ - ... - $\frac{\partial f}{\partial p_n}$ $\frac{\partial g}{\partial q_n}$ \\

$\bullet$ Laguerre Polynomials $\longrightarrow$  \\

$f_{n,l}(r) := L_{n - l - 1}^{2l +  1}$$ \left( \frac{2r}{nr_B} \right)$, where $r_B$ is the Borh-radius\index{Borh-radius}.  \\

$\bullet$ Rydberg formula $\longrightarrow$\\
\begin{equation}
\frac{1}{\lambda} = R_H \left({\frac{1}{n^2} - \frac{1}{m^2}} \right)
\end{equation}

where $R_H$ is the Rydberg constant.  \\


$\bullet$ Rayleigh formula $\longrightarrow$  \\ 

 $\frac{8\pi\nu^2{kT}\vartriangle\nu}{c^3}$ $\hbar$ $e^{-\left(\frac{\hbar}{k}\right) \left(\frac{\nu}{T} \right)}$,  \\

Wiens' formula $\longrightarrow$ \\

$\frac{A \nu^3 e^{-b \nu}}{T \vartriangle \nu}$, where A and b are constants(does not fit experimental facts for small $\frac{ \nu}{T}$, although it avoids the paradox of infinite total energy.  \\


$\bullet$ Radon-Nikodym Theorem(derivative of $\nu$ w.r.t $\mu$) $\longrightarrow$  \\

$\nu (E)$ = $\int_E \rho(s) d\mu(s)$ \\
There exists a measurabe function $\rho$ which is the density of $\nu$ w.r.t $\mu$ \\

$\bullet$ Schr$\ddot{o}$dinger operator $\longrightarrow$  \\ 

$\lim_{X \to \infty} \frac{\beta([0,X])}{\beta_q([0,X])}$ = $(2 \pi \hbar)^N$, where \\

$\beta([0,X])$ is a continuous measure space of the set of classic Hamiltonian, $\beta_q([0,X])$ is the discrete
measure of the number of eigenvalues < (X), and (N) is the number of dimensions in configuration space.  \\

$\bullet$ Schurs' Theorem $\longrightarrow$ $U^*{\mathbf A}U$ = T = $[T_{ij}] \in M_n(\mathbb C)$, $T_{ij}$ = 0(i > j). \\

$\bullet$ Angular momentum in the tensor product space\index{tensor product space} $\longrightarrow$  \\

$\hat j_i$ = ${\hat j_i^{(1)}} \otimes \mathbf{1} + \mathbf{1} \otimes  {\hat j_i^{(2)}}$  satisfy [$\hat j_i, \hat j_j$] = ${i \hbar \epsilon_{ijk} {\hat j_k}}$ acting on $V_1 \otimes V_2$ \\

$\bullet$ Maxwells'-equations\index{Maxwells'-equations} $\longrightarrow$  \\
\noindent
$\nabla \cdot \mathbf E$ = $ \frac{\rho}{\epsilon_0}$ \\
$\nabla \cdot \mathbf B$ = 0  \\
$\nabla \times \mathbf E$ = $ - \frac{\partial B}{\partial t}$ \\
$\nabla \times \mathbf B$ = $\mu_0$  $\left( \mathbf {J} + \epsilon_0 \frac{\partial \mathbf E}{\partial t} \right)$  \\

$\bullet$ Abelian/Commutatitive\index{Abelian/Commutatitive} $\longrightarrow$
\begin{equation} \label{eq:Abelian/Commutatitive}
\forall a,b \in G, \quad a \cdot b = b \cdot a
\end{equation}
$\bullet$ Matrix representation of 2-D Hilbert space\index{Matrix representation of 2-D Hilbert space} $\longrightarrow$  \\ 

$H^2$ = $\frac{1}{2}$ $\begin{pmatrix}1 & 1 \cr 1 & -1 \end{pmatrix}$ $\begin{pmatrix}1 & 1 \cr 1 & -1 \end{pmatrix}$ = $\begin{pmatrix}1 & 0 \cr 0 & 1 \end{pmatrix}$ = I,  \\  

where H is for the Hadamard gate, and (I) is the identity.  \\

$\bullet$ Liebniz-Rule\index{Liebniz-Rule} $\longrightarrow$  \\
\begin{equation}
\frac{d}{dt} \int_{- \infty}^\infty \mid \psi{x,t} \mid^2 dx = \int_{- \infty}^\infty \frac{\partial}{\partial t} \mid \psi{x,t} \mid^2 dx = 0
\end{equation}

$\bullet$ Wave function for a particle in a box: $\longrightarrow$
\begin{equation}
\Psi(x,t) = \sum_{n=1}^\infty C_n \sqrt{\frac{2}{a}} \sin(\frac{n \pi}{a}x) e^{- i\left( \frac{n^2 \pi^2 \hbar}{2m {a}^2} \right)t}
\end{equation}

$\bullet$ Ground-state\index{Ground-state} for a particle in a box $\longrightarrow$ $\Psi_n$(x) = ${\sqrt \frac{2}{a}}$ $\sin \left( \frac{n \pi}{a}x \right)$  \\

$\bullet$ Gaussian-integral\index{Gaussian-integral} $\longrightarrow$  \\

If $z^2$ = ${a}x^2$, and dz = $\sqrt{a}$dx, then $\int_{- \infty}^\infty$ $e^{- a x^2}$ dx = $\sqrt{ \frac{\pi}{a}}$[Highly simplified] \\

$\bullet$ $\mathbf{\textbf{Hookes' Law}\index{Hookes' Law}}$ $\longrightarrow$ \\

F = - $\frac{dV}{dx}$ = - $\frac{d}{dx}$ $\left( \frac{1}{2} k {x^2} \right)$ = -k x = m a = m $\frac{d^2(x)}{dt^2}$ \\

$\bullet$ Energy-spectrum\index{Energy-spectrum} of the Transmon qubit $\longrightarrow$  \\ 

$\omega_j$ = $\left( \omega - \frac{\delta}{2} \right) j$ + $\frac{\delta}{2} j^2$ \\

$\bullet$ Pure density operator, and state\index{state} $\longrightarrow$ \\
\begin{equation}
\mathbf{tr} \rho^2 = tr \rho = \mathbf{1} ; \frac{1}{N} < \mathbf{tr} \rho < \mathbf{1}
\end{equation}

$\rightarrow$  \\

positive-definiteness, where $\lambda_j$ $\geq$ 0 ($\lambda_j$ are the eigenvalues, and 'N' is the Hilbert space dimension. \\

$\bullet$ Traceless-Pauli-matrices\index{Traceless-Pauli-matrices} $\longrightarrow$  \\

$\rho$ = $\frac{1}{2}$ $\left[ \begin{pmatrix}1 & 0 \cr 0 & 1 \end{pmatrix} + x \begin{pmatrix}0 & 1 \cr 1 & 0 \end{pmatrix} + y \begin{pmatrix}0 & -i \cr i & 0  \\ \end{pmatrix} + z \begin{pmatrix}1 & 0 \cr 0 & -1 \end{pmatrix} \right]$  = $\frac{1}{2}$ ( $\mathbf{I}$ + $\vec{r} \cdot \vec{\sigma})$, where $\vec{r} \in \mathbb R^3$,\\ and $\vec{\sigma}$ = (x,y,z). \\

\subsection{Test your comprehension}

Every equation has a meaning, and can be logically explained.

\paragraph{Test Case A}
\noindent

Q1. Time-dependent Schr$\ddot{o}$dinger Equation: \\
  
Describes how a quantum system evolves $\longrightarrow$
\begin{equation}
i \hbar \frac{\partial \psi(x,t)}{\partial t} = \left( -\frac{\hbar^2}{2m} \frac{ \partial^2}{ \partial x^2} + V(x,t) \right) \psi(x,t) = \hat {H} \psi
\end{equation}


   
$\bullet$ $\psi$(x,t): The above equation is a wave-function (Probability amplitude).  \\

One-Dimensional Time-independent Schr$\ddot{o}$dinger Equation $\longrightarrow$ \\

-$\frac{i \hbar^2}{2m}$ $\frac{d^2 \psi}{d {X}^2}$ + V$\psi$ = $\hat {H} \psi$  \\

Q2. Heisenberg Uncertainty Principle/Minimum uncertainty wave-packet $\longrightarrow$ \\ 

$\vartriangle{x} \vartriangle{p}$ $\geq$ $\frac{\hbar}{2}$(Note that $| \frac{i \hbar \langle \psi \mid\psi \rangle}{2i} \mid^2$ = $ \sqrt{\frac{\hbar^2}{4}}$ = $\frac{\hbar}{2}$). \\

$\bullet$ You cannot measure Position and Momentum precisely at the same time \\
 
$\bullet$ More precision in one, less precision in the other \\

Q4. Particle in a box Energies $\longrightarrow$ \\

$E_n$ = $\frac{n^2 \pi^2 \hbar^2}{{2m}{L^2}}$, n = 1,2,3,..  \\
    
$\bullet$ $\hbar$: reduced Plancks' constant($\frac{\hbar}{2 \pi}$) [$\text{h is the Plancks constant}$].  \\
  
$\bullet$ i: imaginary unit/Lateral-unit\index{Lateral-unit}  \\
      
$\bullet$ $\hat {H}$: Hamiltonian (total energy operator)  \\

Q5. Commutation-relation $\longrightarrow$ [ $\hat {x}, \hat {p}$ ] = i$\hbar$ \\

Significance: \\ 

$\bullet$ Leads to uncertainty principle  \\

$\bullet$ Shows position and momentum, are incompatible observables \\
 
\paragraph{Test Case B}
\noindent

Q6. Particle in a box \\

(a) Wave-function  \\
 
Boundary conditions:  \\
  
$\bullet$ $\psi$(0) = 0 $\implies$ $\mathbf {B}$ = 0  \\

$\bullet$ $\psi ( \mathbf{L}$) = 0 $\sin ({k}{\mathbf L})$ = 0  \\
  
So:  \\

    k = $\frac{{n} \pi}{L}$  \\

Final wave-function: \\
 
$\psi_n$(x) = $\sqrt{\frac{2}{L}}$ $\sin$ $\left( \frac{{n} \pi {x}}{L} \right)$ \\
 
  
(b) Energy-Level\index{Energy-Level} \\

$E_n$ = $\frac{n^2 \pi^2 \hbar^2}{2{m}{L}^2}$ \\
 
(c) Probability\index{Probability} \\
 
$\mathbf {P}$ = $\int_0^{ \frac{L}{2}}|\psi_1|^2$ dx  
              = $\frac{2}{L}$ $\int_0^{\frac{L}{2}}$ $\sin^2$ $\left( \frac{\pi {x}}{L} \right)$ dx  \\
  
Result:  \\
  
$\mathbf {P}$ = $\frac{1}{2}$  \\

Q7. Commutator and Uncertainty  \\
  
(a) Show: \\ 

    [$\hat {x}, \hat {p}]$ = i$\hbar$  \\

    Apply to function f(x): \\
 
        $\hat {x} \hat {p}$f = x(-i$\hbar$f') \\
 
                                 $\hat {p} \hat {x}$f = -i$\hbar$ $\frac{d}{d{x}}$(xf) = -i$\hbar$(f + xf')  \\
   Subtract:  \\
  
        [$\hat x, \hat p]$f = i$\hbar$f \\
 
   Thus: \\
 
        [$\hat x, \hat p]$ = i$\hbar$ \\
 
(b) Uncertainty-Principle\index{Uncertainty-Principle} \\
                                                                                                
    General relation: \\
 
       $\vartriangle \mathbf A \vartriangle \mathbf B$ $\geq$ $\frac{1}{2}$ $\mid \langle [ \mathbf A, \mathbf B] \rangle \mid$ \\
  
    Substitute: \\
 
       $\vartriangle x \vartriangle p$ $\geq$ $\frac{\hbar}{2}$  \\

Q9. Hydrogen-Atom\index{Hydrogen-Atom} \\
                                                                                       
    (a) Energy-Level\index{Energy-Level} \\

    $E_n$ = - $\frac{13.6}{n^2}$ eV \\
 
    (b) Quantum-number\index{Quantum-number}  \\
                                                                                                 
    $\bullet$ n: Energy-Level\index{Energy-Level}  \\
  
    $\bullet$ $\ell$: orbital-shape\index{orbital-shape} \\
 
    $\bullet$ m: orientation \\
 
    (c) Why Energy is Negative\index{Why Energy is Negative}  \\
  
    $\bullet$ Electron is bound to nucleus  \\
  
    $\bullet$ Energy is required to remove electron   \\
 
    $\bullet$ Bound-states\index{Bound-states} \\

$\mathbf{\textbf{Wave-function-collapse\index{Wave-function-collapse}}}$   \\
 
$\bullet$ Before measurement, we have superposition  \\
  
$\bullet$ Measurement leads to a definite state  \\

A collapse = Transition from probability to reality$\mathbb R$


\paragraph{Additional Notes}
\noindent

I.(a) Completeness  \\

A Hilbert space is complete if:  \\

    Every Cauchy sequence $\psi_n$ $\subset$ $\mathscr H$ converges to a limit in $\mathscr H$  \\
  
Formally:  \\
  
    $\Vert \psi_n - \psi_m \Vert$ $\longrightarrow$ 0 $\implies$ $\exists \psi$ $\in \mathscr H$, $\psi_n$ $\longrightarrow$ $\psi$ \\
   

(b) Extension to Orthonormal basis  \\
  

Given orthonormal-set\index{orthonormal-set} $e_1$:  \\
  
$\bullet$ If not complete, there exists $\psi$ $\neq$ 0 orthogonal to all $e_1$  \\
  
$\bullet$ Add $\psi$, normalize enlarge set \\ 

Zorns'-Lemma\index{Zorns'-Lemma} argument:  \\

Maximal orthonormal-set\index{orthonormal-set} = orthonormal basis  \\

(c) Projection is bounded\index{bounded}  \\

Projection-operator\index{Projection-operator} $\mathbf P$:  \\  
    
    $\mathbf P^2$ = $\mathbf P$, $\mathbf {P}\dagger$ = $\mathbf P$ \\ 
    
For any $\psi$: \\  
    
    $\Vert \mathbf P \psi \Vert$ $\leq$ $\Vert \psi \Vert$ \\ 
    
    Thus: \\
 
        $\Vert \mathbf P \Vert$ $\leq$ 1. therefore bounded\index{bounded}. \\

II. Spectral-Theorem\index{Spectral-Theorem} \\
 
(a) Statement \\ 

For self-adjoint-operator\index{self-adjoint-operator} $\hat A$:  \\
  
    $\hat A$ = $\int \lambda$ dE($\lambda$)  \\
  
$\bullet$ E($\lambda$): Projection-valued-measure\index{Projection-valued-measure}  \\
  
$\bullet$ Generalizes diagonalization \\

(b) Spectrum-Types\index{Spectrum-Types} \\
 
$\bullet$ $\mathbf {Dicrete}$ eigenvalues (Bound-states\index{Bound-states}) \\
 
$\bullet$ $\mathbf {Continuous}$ scattering-states\index{scattering-states}  

(c) Measurement-Interpretation\index{Measurement-Interpretation} \\ 

Probability of outcome in set\index{Probability of outcome in set} $\vartriangle$: \\
 
    $\mathbf {P}(\vartriangle)$ = $\langle \psi$|E($\vartriangle$)|$\psi \rangle$  \\

    Measurement-outcome\index{Measurement-outcome} = spectrum of operator\index{spectrum of operator}  \\

III. Operators and Dynamics  \\  

Time Evolution  \\  

(a) Derivation \\ 

Start: \\
 
    $\langle \mathbf {A} \rangle$ = $\langle \psi$|$\mathbf {\hat A}$|$\psi \rangle$ \\ 

Differentiate: \\
 
$\frac{d}{dt} \langle \mathbf {A} \rangle$ = $\langle$ $\frac{\partial \psi}{\partial t}$|$\mathbf A$|$\psi \rangle$ + $\langle \psi$|$\mathbf A$|$\frac{\partial \psi}{\partial t} \rangle$ + $\langle \psi$|$\frac{\partial \mathbf A}{\partial t}$|$\psi \rangle$   \\   
        
Use Schr$\ddot{o}dinger$ equation:  \\
  
    i$\hbar \frac{\partial \psi}{\partial t}$ = $\mathbf H \psi$ \\ 
    
Substitute: \\
 
$\frac{d \langle \mathbf A \rangle}{dt}$ = $\frac{i}{\hbar}$ $\langle$[$\mathbf {\hat H}$, $\mathbf {\hat A}$] $\rangle$ + $\langle \frac{\partial \mathbf A}{\partial t} \rangle$  \\  
                              
    
 (b) Apply to position \\  

$\mathbf {\hat H}$ = $\frac{p^2}{2m}$ \\
 
$\mathbf{\hat  H}$,x = $\frac{1}{2m}$ [$p^2,x]$  \\ 

Use:  \\
  
    [$\hat p$,x] = -i$\hbar$  \\
    
Result:  \\

    $\frac{d \langle x \rangle}{dt}$ = $\frac{\langle p \rangle}{m}$ \\
 
    Velocity operator emerges  \\ 
    
(c) Conservation condition \\ 

If: \\
 
    [$\mathbf {\hat H}, \mathbf {\hat A}$] = 0 and $\frac{\partial A}{\partial t}$ = 0 \\
 
    Then: \\
 
        $\frac{d \langle A \rangle}{dt}$ = 0 \\
 
        A is conserved  \\

IV. Harmonic Oscillator Algebra \\ 

(a) Commutator\index{Commutator}  \\
  
$\hat a$ = $\frac{1}{\sqrt{2 \hbar {m} \omega}}$(m $\omega$ x + ip)  \\  

Compute: \\
 
    [$\hat a, \hat a{\dagger}]$ = 1 \\ 
    
(b) Hamiltonian\index{Hamiltonian}
\begin{equation} \label{eq:Hamiltonian}
\mathbf{\hat H} = \hbar \omega (\hat a{\dagger} \hat a + \frac{1}{2})
\end{equation}

(c) Orthonomality\index{Orthonomality}  \\  

$\langle n \mid m \rangle$ = $\delta_{nm}$ \\ 

Proof:  \\
  
    $\bullet$ Ladder-Operators\index{Ladder-Operators} generate states  \\
  
    $\bullet$ Inner-product\index{Inner-product} preserved recursively  \\

V. Advanced Systems \\ 

Angular-Momentum\index{Angular-Momentum}  \\

(a) Lie Algebra \\ 

[$J_i,J_j$] = i$\hbar \epsilon_{ijk} J_k$  \\
  
Derived from rotational symmetry generators \\ 

(b) Ladder-Operators\index{Ladder-Operators} \\

$J_\pm$ = $J_x \pm$ i$J_y$ \\
 
[$J_z,J_\pm$] = $\pm \hbar J_\pm$  \\  

(c) Eigenvalues\index{Eigenvalues}  \\

$J^2$ $\mid$ ${jm} \rangle$ = $\hbar^2j(j + 1)| {jm}$ $\rangle$ \\
 
$J_z \mid {jm} \rangle$ = $\hbar m \mid{jm} \rangle$  \\

VI. Perturbation-Theory\index{Perturbation-Theory}  \\  

(a) First Order  \\  

$E_{n}^{(1)}$ = $\langle n^{(0)}$|V|$n^{(0)} \rangle$ \\ 

(b) Second Order  \\
  
$E_{n}^{(2)}$ = $\sum_{k \neq n}$ $\frac{| \langle {k}|{V}|{n} \rangle |^2}{E_{n}^{(0)} - E_{k}^{(0)}}$ \\ 
    
(c) Degeneracy\index{Degeneracy} \\
 
$\bullet$ Degenerate-states\index{Degenerate-states} mix \\
 
$\bullet$ Solve matrix: \\
  
    $\mathbf V_{ij}$ = $\langle i\mid \mathbf {V}|j \rangle$ \\
 
    Diagonalize\index{Diagonalize} within subspace \\

VII. Conceptual \\ 


Decoherence\index{Decoherence} \\  

$\bullet$ System interacts with environment \\
 
$\bullet$ Phase information lost \\
 
$\bullet$ Superposition(classical mixture)  \\
  
Explains Classical world emergence  \\

VIII. Path-Integral\index{Path-Integral} \\ 

Core Idea  

Amplitude\index{Amplitude} = $\sum_{paths}$ $e^{\frac{i}{\hbar} S[x(t)]}$  \\  

Action\index{Action} \\ 

S = $\int$L dt  \\  

Classical-Limit\index{Classical-Limit}  

$\bullet$ Does not account for a dominant path/stationary action  \\

$\bullet$ Quantum Mechanics recovers Classical Mechanics \\

Quantum Mechanics becomes Linear Algebra, Operator Theory, and Symmetry, NOT just formulas.  \\

\chapter{Example Functions}

\subsection{Equations}

Plane-waves\index{Plane-waves} $\longrightarrow$  \\

$\psi(x,t)$ = $e^{(ikx - i\omega t)}$; we then differentiate this function to obtain the
Energy, and Momentum operator, using the Chain Rule...and applying Plancks' rule\index{Plancks' rule}(multiplying by i$\hbar$).  \\

$i \hbar \frac{\partial}{\partial t} \psi(x,t)$ = $i \omega e^{(ikx - i \omega t)}$ = $i \omega \psi(x,t)$, seeing the eigenfunction on the left side of the equation(taking the derivative), and the resulting constant eigenvalue on the right hand side of the equation.  \\

Energy-operator\index{Energy-operator}($\hat E$) $\rightarrow$
\begin{equation}
i \hbar \frac{\partial}{\partial t} \psi(x,t) = i \hbar(-i) \omega \psi(x,t) = \hbar \psi(x,t)
\end{equation}

Momentum operator($\hat P$) $\rightarrow$  \\
\begin{equation}
(-i \hbar) \frac{\partial}{\partial x} \psi(x,t) = ( -i \hbar) {ik}e^{ikx - i \omega t} =  \hbar k \psi(x,t)
\end{equation}
(Note -i cancels i), k = $\frac{2 \pi}{\lambda}$; $\omega$ = kc  \\

Vector and Matrix multiplication $\rightarrow$ A$\vec{v}$ = $ \lambda$ $\vec{v}$, where $\lambda$ is an eigenvalue, and $\vec{v}$ is the eigenvector. \\

Energy of a photon $\rightarrow$ E = $\hbar \omega$ ; $\vec{P}$ = $\hbar$k ; $\vartriangle$ m = 0(m is for mass); E = Pc \\

Discrete Fourier Transform $\rightarrow$  \\

Consider the space of 2$\pi$ periodic functions f: $\mathbb R$ $\rightarrow$ $\mathbb C$, with a continuous 1st derivative;
We have $E_k(x)$ = $e^{ikx}$, k $\in$ $\mathbb Z$. Using Eulers' formula($e^{i \alpha}$ = $\cos \alpha + i \sin \alpha)$; $e^{2 \pi i}$ = 1,
$E_k(x) \text{by substitution}$, $\equiv$ $E_k(x)$ = $\cos(kx) + i \sin(kx)$;  
Here $E_k$ is an orthonormal set.  \\

Now, for the Hermitian Scalar Product $<f \mid$ $\frac{1}{2 \pi}$ $\int_0^{2 \pi}$ $\overline {f(x)}$ $\cdot$ g(x) dx, we have $<E_k \mid E_s>$ =  $\frac{1}{2 \pi}$ $\int_0^{2 \pi}$ $e^{- ikx}$ $\cdot$ $e^{isx}$ = $\frac{1}{2 \pi}$ $\int_0^{2 \pi}$ $e^{i(s - k)x} dx$(factored),
= $\frac{1}{2 \pi}$ $\frac{1}{i(s - k)}$ $e^{i(s - k)x}$ $|_{x = 0}^{x = 2 \pi}$.  \\

Note that $\iff$ s and k = 0,
$\left\{
\begin{array}{ll}
0,  \text{if s $\neq$ k} \\
1,  \text{if s = k}       \\
\end{array}              \\
\right.$  \\
    
Since its periodic($\frac{1}{2 \pi}$ $\int_0^{2 \pi}$1 dx = 1, if s = k); With the resulting anti-derivative of the function.  
So, we can pass from ($f_0, f_1,...,f_{n-1})$ $\rightarrow$ Fourier coefficients ($c_0,c_1,...,c_{n - 1}$). This is the Discrete Fourier series; A discrete signal capturing frequency information, given that
$C_k$ = $\frac{1}{\sqrt N}$ $\sum_{j = 0}^{N - 1}$ $f_j$ $e^{ \frac{i 2 \pi {kj}}{N}}$(DFT); Note that $C_k$ = $<E_k \mid f>$ = $\sum_{j = 0}^{N-1}$ $\overline {E_{kj}}$ $f_{j}$, where N is the  number of samples.  \\

Inverse Discrete Fourier Transform(IDFT) $\rightarrow$  \\
 
$f_j$ = $\sum_{k = 0}^{N - 1}$ $C_k$ $E_{kj}$ = $\frac{1}{\sqrt N}$ $\sum_{k = 0}^{N - 1}$ $C_k$ $e^{- \frac{i 2 \pi {kj}}{N}}$.  \\

Langranges' Theorem $\rightarrow$  \\ 

Let G be a group of an order $|G|_n$; Then, $\forall$ a $\in$ G, $a^n$ = e; 
So, $\mathbb Z_p^*$ = G. It has a multiplicative inverse.  \\

Fermats' Little Theorem $\rightarrow$ Let p be a prime; Then $\forall$ a $\in$ $\mathbb Z_p^*$, $a^{p - 1}$ = 1 mod p, (a) is not divisible by p. \\

In an RSA setting, m = $p \cdot q$, p and q are Primes; For n = $|\mathbb Z_m^*|$ = (p -1) $\cdot$ (q - 1); 
If (a) is not divisible by p or q, then $a^{(p - 1)(q - 1)}$ = 1 mod {pq}. \\

Hadamard $\longrightarrow$ H; $ \mid 0>$ $\rightarrow$  \\

$\frac{1}{\sqrt 2}$ $ \mid 0>$ + $\mid 1>$, and $ \mid 1>$ $\rightarrow$ $ - \frac{1}{\sqrt 2}$ $ \mid 0> + \mid 1>$  \\

H =
$
\begin{bmatrix}
\frac{1}{\sqrt 2} &  - \frac{1}{\sqrt 2} \\
\frac{1}{\sqrt 2} &    \frac{1}{\sqrt 2}
\end{bmatrix}
$

Eulers' Identity $\rightarrow$ $e^{i \pi}$ = -1  \\

Adiabatic-approximation\index{Adiabatic-approximation} $\rightarrow$ $\Psi(t)$ = $e^i {\theta_k (t)}$ $e^i {\gamma_k (t)}$ $|\mid  \psi_k(t)>$  \\

Lowest-Landau-level\index{Lowest-Landau-level} $\rightarrow$  \\
\begin{equation}
\Psi_m(z) \propto z^m e^{-|z|^2}
\end{equation}

\begin{remark}  
Quantum Fourier Transform $\rightarrow$ $U_{QFT}$ is not a Hermitian operator,
$U_{QFT}$ $\neq$ $U_{QFT}^{\dagger}$. It creates an equally weighted superposition of states, leaving amplitudes as complex numbers,
while Hadamard leves them as $\mathbb R$. QFT also allows derivation of Fourier Basis, used in Quantum Mechanics inorder to ease and describe  various calculations involved in the Theory.
\end{remark}

Baker-Campbell-Hausedorff Theorem\index{Baker-Campbell-Hausedorff Theorem} $\rightarrow$ \\
\begin{equation}
\log(e^Xe^Y) = X + \int_{0}^1 g(e^{\text{ad X}} e^{\text{t ad Y}})(Y)dt.
\end{equation}

For a complete derivation, see Dynkins' analysis of Lie groups.

Perturbed-System\index{Perturbed-System} $\rightarrow$  \\

$\hat H$ = $\frac{\hat p^2}{2m}$ + $\frac{1}{2} m^2 \omega^2 x^2$ - $\frac{\lambda}{2} x^2$,  
with an exact solution of E $\approx$ $\frac{\hbar}{2}$ - $\frac{\hbar \lambda}{4m \omega}$  \\

Perturbation-Theory\index{Perturbation-Theory} $\rightarrow$ \\
\begin{equation}
E = E^{(0)} + \langle \psi^{(0)} \mid \lambda \hat w \mid \psi^{(0)} \rangle = \frac{\hbar \omega}{2} + \langle \psi^{(0)} \mid {-(\frac{\lambda}{2}) x^2} \mid  \psi^{(0)} \rangle =  \frac{\hbar \omega}{2} - \lambda \frac{\lambda \hbar}{4m \omega} =  - \frac{\lambda}{2} \langle x^2 \rangle =  - \frac{\lambda}{2} \sigma^2.
\end{equation}

Note that $\psi^{(0)}$(x) = c$e^{- \frac{x^2}{4 \sigma^2}}$ and $\sigma^2$ = $\frac{\hbar}{2m \omega}$.  \\

Cross-Product\index{Cross-Product} $\rightarrow$ \\ 
    
The cross product of u = $(u_1,u_2,u_3)$ and v = $(v_1,v_2,v_3)$ is a vector                                                                                            
(u $\times$ v) = \\

$\begin{pmatrix} \mathbf{i} & \mathbf{j} & \mathbf{k} \cr u_1 & u_2 & u_3 \cr v_1 & v_2 & v_3 \end{pmatrix}$ =  
($u_2v_3 - u_3v_2)\mathbf{i}$ + ($u_3v_1 - v_1u_3)\mathbf{j}$ + ($u_1v_2 - u_2v_1)\mathbf{k}$.  \\
    
The vector (u $\times$ v) is perpendicular to u and v.  \\

The cross product (v $\times$ u) is -(u $\times$ v).  \\ 

Note that (u $\times$ u) = 0,  \\
 
$\Vert u  \times v \Vert$ = $\Vert u \Vert$ $\Vert v \Vert \mid \sin \theta \mid$, and $\mid u \cdot v \mid$ = $\Vert u \Vert$ $\Vert v \Vert$ $\mid \cos \theta 
\mid$.  \\

Time evolution Hamiltonian operator\index{Time evolution Hamiltonian operator} \label{eq:
Time evolution Hamiltonian operator}, form a complete basis $\rightarrow$  \\

\begin{equation}
i \frac{d}{dt} \mid \psi \rangle(t) = i \frac{d}{dt} \hat U (t) \mid \psi (0) \rangle = \hat H \hat U(t) \mid \psi(0) \rangle = \hat H \mid \psi (t) \rangle
\end{equation}
\noindent

Note that $\frac{d}{dt}$ $\hat U$(t) = $ - \frac{i}{\hbar}$ $\hat H$ $\hat U$(t), its like the Schr$\ddot{o}$inger equation!, and
$\hat U$(t)|$\psi_n \rangle$ = $e^ \frac{ - iE_nt}{\hbar}$ $\psi_n \rangle$ too. \\

Quantum-jump\index{Quantum-jump} $\rightarrow$ \\
    
An abrupt projection of a system to an Eigenstate $\psi_0$, with a dissipative process(spontaneous emission). \\

The evolution of the system starts from zero, ruled by the effective Hamiltonian. The modification of $\psi (t)$ by $\mathbf{\textbf{non-observation}}$ of spontaneous emission, reduces the population of the state excited by 1 - $\frac{1}{2} \Gamma$dt, while the ground-state remains unchanged. Every $\mathbf{\textbf{quantum jump}}$ projecting the system into the ground-state constitutes a $\mathbf{\textbf{measurement\index{measurement}}}$.  \\

Monte Carlo Wave-function(MCWF) simulation\index{Monte Carlo Wave-function(MCWF) simulation} $\rightarrow$  \\ 

$\psi$(t) $\hookrightarrow$ $\mid \psi(t + dt) \rangle$ $\equiv$  \\

$\left\{
\begin{array} {ll}
\frac{(1 - i \hat H dt) \mid  \psi(t + dt) \rangle}{\sqrt{ \langle \psi(t) \mid \psi(t) \rangle}} \text{ if $\zeta$ > 1 - $\langle \psi(t) \mid\psi(t) \rangle$} \\
\mid \psi_0 \rangle \text{ if $\zeta$ < 1 - $\langle \psi(t) \mid \psi(t) \rangle$} \\
\end{array}
\right\}.$

Concept of Magnitude\index{Concept of Magnitude} $\rightarrow$ If w = 3 + 4i, then $\mid w \mid$ = $\sqrt{\Re(w^2) + \Im(w^2)}$ = $\sqrt{3^2 + 4^2}$ = 5  \\

Cube roots of 1 $\rightarrow$  \\

$z^3$ = 1 $\implies$ $(x + iy)^3$ = 1 $\implies$ $x^3$ + $3x^2(iy)$ + $3x(iy)^2$ + $(iy)^3$ = 1 $\implies$ $(x^3 - 3xy^2)$ + $i(3x^2y - y^3)$ = 1;  \\

So, $x^3 - 3xy^2$ = 1, and $3x^2y - y^3$ = 0.  \\

x = 1, y = 0, and z = 1, is a solution. \\ 

If x and y are allowed to be complex, then there are three solutions (if x is non-zero). It has
three cube roots. A real number has one real cube root, and two further cube roots which form a
complex conjugate pair. For instance, the cube roots of 1 are: 1, $\frac{1}{2}$ + $\frac{\sqrt 3}{2}$, $ - \frac{1}{2}$ - $\frac{\sqrt 3}{2}$.  \\ 

Note that,  \\ 
the principal cube root of a negative number is a complex number, instance $\sqrt[3]{-8}$ is NOT -2, but rather (1 + i$\sqrt{3}$). Quartic equations can also be solved in terms of cube roots and square roots. \\

Seiberg-Witten functional\index{Seiberg-Witten functional} $\rightarrow$ $D_{A \varphi}$ = 0,  \\
\begin{equation}
F^+_A = \frac{1}{4} \langle e_j \cdot e_k \cdot \varphi, \varphi \rangle e^j \wedge e^k
\end{equation}

Fundamental electron velocity of the Hydrogen states\index{velocity of the Hydrogen states} $\rightarrow$  \\
\begin{equation}
E_1 = \frac{- \hbar^2}{2m_e a^2_B}
\end{equation}

Relativistic velocity of the electron $\rightarrow$  \\

 v = $\sqrt{\frac{2E_1}{m_e}}$ = $\frac{\hbar}{m_e a_B}$ = $\frac{e^2}{4 \pi \epsilon_0 \hbar}$ = $\alpha$c \\

Velocity-operator\index{Velocity-operator} $\langle \hat{v} \rangle$ $\rightarrow$ \\

$\frac{i \hbar}{m}$ $\int \psi^*$ $\frac{\partial \psi}{\partial x}$ dx \\
 
[Velocity as motion of expectation, applies to Momentum(p), $\mathbf H$(Hamiltonian total Energy of a system), and Kinetic Energy + Potential Energy\quad(T). \\

Bells' Inequality $\rightarrow$ $\mid P\mathbf{\textbf{(a,b)}}$ - $P(\mathbf{\textbf{a,c)}} \mid$ $\leq$ $1 + P\mathbf{\textbf{(b,c)}}$  \\

Klein-Gordon equation $\rightarrow$ \\

$\frac{1}{c^2}$ $\frac{\partial^2}{\partial t^2}$ - $\nabla^2$ +  $ ( {\frac{m_e c}{\hbar}})^2$] $\psi$ = 0  \\
    
Note that [$\partial^\mu, \partial^nu$ = [$\partial^{\mu \nu}$]. The matching-index-height rule\index{matching-index-height rule} is satisfied.  \\

\section{Sampled wave-function distribution}
\label{Sampled wave-function distribution}

\begin{figure}[!h]
\includegraphics[width=10cm]{Ground state search.png}
\caption{Mnimizing Energy for a ground-state}.
\label{fig:Ground state search.png}
\end{figure}

\begin{figure}[!h]
\includegraphics[width=10cm]{Sampled wave function distribution.png}
\caption{Estimated Ground State Energy, with its' analytical solution}.
\label{fig:Sampled wave function distribution.png}
\end{figure}

\begin{figure}[!h]
\includegraphics[width=10cm]{BW.png}
\caption{Breit-Wigner distribution}.
\label{fig:BW.png}
\end{figure} 

\cite{https://citizendium.org/wiki/index.php?title=Particle_in_a_box&oldid=984214}

\chapter{Angular Momentum}

The Angular Momentum of a rigid body rotating with an angular velocity ($\vec{\omega}$) about a fixed axis, and having a moment of inertia (I) matrix elements, with respect to that axis, is given by L = $\mathit{I} \vec{\omega}$ \quad ($\vec{\omega}$ in rad/s).
The rate of change of angular momentum of a particle equals the torque of the net force acting on it. The moment of inertia about some given axis is related to the moment about a parallel axis through the center of mass. It's represented on the SO(3) unitary group.  \\

\section{Rigid-Body Rotation}

Any general displacement of a rigid body, can be represented by a translation plus a rotation. We split the problem of rigid body motion into two separate phases, one concerned solely with the translational motion of the body, the other, with its' rotational motion. The total
angular momentum divides naturally into contributions from the translation of the center of mass and, from the rotation about the center of mass.

The total angular momentum will involve the Cartesian coordinates of the center of mass, and
the Euler angle coordinates. The Potential Energy in a uniform gravitational field will depend only upon the Cartesian vertical coordinate of the center of gravity. The two groups of coordinates will then be completely separated, and the translational and rotational problems can be solved independently of each other. We will link derivatives relative to a coordinate system fixed at some point in the rigid body.

The derivatives relative to any rigid body axes must vanish, and with $\vec{\omega}$ as being the angular velocity vector appropriate to the first coordinate system.The angular momentum vector is related to the angular velocity by a linear transformation. The operator (I), a Tensor of the second rank, acting upon the vector $\vec{\omega}$, results in the physically new vector (L). This rotation can be expressed in terms of the Euler angles ($\phi$), ($\theta$), and ($\psi$). A proper choice of these angles will transform (I) into its diagonal form:
$(I_D)$ = $\begin{pmatrix}\mathit{I_1} & 0 & 0 \cr 0 & \mathit{I_2} & 0 \cr 0 & 0 & \mathit{I_3} \end{pmatrix}$, where $(I_1)$, $(I_2)$, and $(I_3)$, are the eigenvalues of (I), referred to as the components of the principal moment of inertia tensor. Three roots are the desired principal moments.
All directions perpendicular to the axis of symmetry are then alike, which is the mark of a double root to the equation. The principal axes are then the symmetry axis, and any two perpendicular axes in the plane normal to the symmetry axis. The principal moments of inertia cannot be negative, because as the diagonal elements in the principal axes system have the form of sums of squares. For one of the principal moments to vanish, all points of the body must be such that two coordinates of each particle are zero. See Figure~\ref{fig:Spinning top.pdf}  \\
This can happen only if all points of the body are collinear with the principal axis corresponding to the zero principal moment. Any two axes perpendicular to the line of the body will then be the other principal axes. Indeed, this is a limiting case of a body with an axis of symmetry passing through the origin. If two of the roots of the equation are equal, the inertia ellipsoid thus has two equal axes, and is an ellipsoid of revolution. If all three principal moments are equal, the inertia ellipsoid is a sphere.  \\
It is worth emphasizing that the inertia tensor (I), and all the quantities associated with it(principal axes, principal moments, inertia ellipsoid...), are only relative to some particular point fixed in the body. If the point is shifted elsewhere in the body, all the quantities will in general be changed.The $\vec{\omega}$ vector, will always move such that the corresponding normal to the inertia ellipsoid is in the direction of the angular momentum. In particular, the direction of (L) is fixed in space, and it is the inertia ellipsoid (fixed with respect to the body), that must move in space in order to preserve this connection between $\omega$ and (L). Both $\mathit{T}$, the Kinetic Energy, and $\mathit{L}$, the Angular Momentum, are constants of the motion, and the tangent plane is always a fixed distance from the origin of the ellipsoid. The total angular velocity $\vec{\omega}$ is constant in magnitude, and precesses about the (z) axis with the same frequency. The symmetry axis, is one of the principal axes, and will be chosen as the (z) axis of the coordinate system fixed in the body. Since one point is stationary, the configuration of the top is completely specified by the three Euler angles: ($\theta$) gives the inclination of the (z) axis from the vertical, ($\phi$) measures the azimuth of the top about the vertical, while ($\psi$) is the rotation angle of the top about its own (z) axis. $\mathbf{\textbf{The distance of the center of gravity (located on the symmetry axis), from the fixed point, will be denoted by (I).}}$.The angular velocity of precession, is given by $\dot \phi$ = $\frac{1}{\mathit I_3 \omega_3}$ $\frac{\partial V}{\partial (\cos \theta)}$.There are two examples of precession phenomena in nature: The types of precession that arise from the torques induced by Earths' equatorial “bulge," and the precession of moving charges in
a magnetic field.  \\
The Earth is a top whose figure axis, is precessing about the normal to the ecliptic, the plane of Earths' orbit, a motion known astronomically as the precession of the equinoxes. The precession rate, relative to the orbital angular velocity, $\omega_0$, is then: $\frac{\dot \phi}{\omega_0}$ = - $\frac{3}{2}$ $\frac{\omega_0}{\omega_3}$ $\frac{\mathit I_3 - \mathit I_1}{\mathit I_3} \cos \theta$.  \\
$\theta$ = $23^{\circ} 27^{\prime}$, from the solar-induced precession, such as to cause a complete rotation of the figure axis, about the normal to the ecliptic (plane of Earths' orbit).\\
Magnetic moment, and the angular momentum also have a similar form. For the Classical gyromagnetic ratio, the precession angular velocity is known as the Larmor-frequency\index{Larmor-frequency}. To first order in $\mathbf{\textbf{B}}$, the effect of a constant magnetic field on a Classical
system, is to superimpose on its' normal motion, a uniform $\mathbf{\textbf{precession\index{precession}}}$, with $\mathbf{\textbf{angular
frequency}}$.

\begin{figure}
\begin{center}
\includegraphics[width=10cm]{Spinning top.pdf}
\end{center}
\caption{The relation between $\phi$, $\theta$, and $\psi$}
\label{fig:Spinning top.pdf}
\end{figure}

\chapter{Spin}

This is a property of the universe.  \\
A two-particle wave-function is a difficult proposition, calculated with $\hat H$, integrated on all dimensions of an indistinguishable particle! See Eq.~ \eqref{eq:Hamiltonian}. Electrons are confined, and DON'T interact.  \\

\paragraph{Quantum Projection}
\noindent

\textcolor{blue}{Now, for the contradiction of the proposal:}  \\
Suppose that we have a molecule containing two atoms in a state in which the total spin is zero, and that the spin of each atom is $\frac{\hbar}{2}$, and does not change the total angular momentum. The two atoms will begin to separate, and will soon cease to interact. Their combined spin angular momentum, however, remains equal to zero, because by hypothesis, no torques have acted on the system. If the spin were a Classical angular momentum variable, the interpretation of this process would be as follows: While the two atoms were together in the form of a molecule, each component of the angular momentum of each atom would have a definite value that was always opposite to that of the other, thus making the total angular momentum equal to zero. When the atoms separated, each atom would continue to have every component of its spin angular momentum opposite to that of the other. The two spin-angular-momentum vectors would therefore be $\mathbf{correlated}$. These correlations were originally produced when the atoms interacted in such a way as to $\mathbf{\textbf{form a molecule of zero total spin}}$, but after the atoms separate, the correlations are maintained by the deterministic equations of motion of each spin vector separately, which bring about $\mathbf{\textbf{conservation of each component of the separate spin-angular momentum vectors}}$.
This means that, a measurement of any component of the spin of atom No. 1 provides, as in Classical Theory, an indirect measurement of the same component of the spin of atom No. 2, without in any way disturbing that particle! For since there is no interaction with particle No. 2, the process of measurement cannot have affected this particle in any way.
But now let us remember that, in each case, the observer is always free to reorient the apparatus in an arbitrary direction while the atoms are still in flight, and thus to obtain a definite (but unpredictable) value of the spin component in any direction that he chooses. Since this can be accomplished without in any way disturbing the second atom, we conclude that if ERP assumptions are applicable, precisely defined elements must exist in the second atom, corresponding to the simultaneous definition of all three components of its spin. Because the wave-function can specify, at most, only one of these components at a time with complete precision, we are then led to the conclusion that the wave-function does not provide a complete description of all elements of reality existing in the second atom.
If this conclusion were valid, then we should have to look for a new Theory in terms of which a more nearly complete description was possible. We shall see, however, that the analysis given by ERP involves in an integral way the implicit assumptions that the world is actually made up of separately existing and precisely defined “elements of reality.” Quantum Theory, however, implies a quite different picture of the structure of the world at the microscopic level. This picture leads, as we shall see, to a perfectly rational interpretation of the hypothetical experiment of ERP within the present framework of the Theory. Such correlations can be demonstrated, for example, in a process in which the (z) component of the spin of each atom is measured by allowing each atom to pass through a separate Stern-Gerlach apparatus. After the measurement, the system has been transformed from one that has a indefinite combined angular momentum, $\mathbf{\textbf{The wave-function for a system of zero total spin is the same}}$. Both particles have a positive value together, or that both particles have a negative value together, with the constant phase factor of [$e^\frac{\pm \mu \mathscr H \alpha^{\prime} \vartriangle t(x_1 - x_2)}{\hbar}$], where each atom obtains an opposite momentum corresponding to its opposite spin. Because the apparatus is Classically describable, the apparatus wave-functions [which depend on the axes of $(z_1)$ and $(z2)$], multiply the spin wave-function by phase factors, so that we finally obtain: $\psi$ = $\frac{1}{\sqrt 2}$ $(\psi_{x_1}$ $e^{i \alpha x_1 }$ + $\psi _{x_2}$ $e^{i \alpha x_2})$, where $\alpha x_1$ and $\alpha x_2$ are seperate in the phase factor. We conclude that the value of $\sigma_x$, for each particle is also correlated to that of the other, in such a way that the sum of the two is zero. We see therefore, that the type of correlation of $\sigma_x$ which can develop depends on the sign with which $\psi_{x_1}$ and $\psi{x_2}$ are added, and therefore also on the combined $\mathbf{\textbf{angular momentum}}$.
We assert that the wave-function (in principle), can provide the most complete possible description of a system,
that is consistent with the actual structure of matter.
There are actually no precisely defined “elements of reality” belonging to the $\mathbf{\textbf{electron}}$.
The $\mathbf{\textbf{probability}}$ for the development of a definite value of any spin component in a suitable process of measurement is proportional to the square of the amplitude of the coefficient of the part of the wave-function corresponding to this component.
Thus, we must give up the Classical picture of a precisely defined spin variable associated with each atom, and replace it by our Quantum concept of a potentiality, the probability of whose development is given by the wave-function.
As long as definite phase relations exist between $(x_1)+$ and $(x_2)-$, one cannot categorize (or classify) the system as having a spin which corresponds either entirely to $(x_1)+$ or entirely to $(x_2)-$, with respective probabilities $\mid {x_1+} \mid^2$, and $\mid {x_2-} \mid^2$.
It is only when the wave-function is an eigenfunction of a given spin component, that the system is certain (in interaction with a suitable apparatus) to develop a predictable value of that spin component.The wave-function therefore describes the propagation of correlated potentialities.
The concept of spin is a one-to-one correlation, and precisely define Mathematical abstractions.

\section{Self-rotation of the Electron}

Spin is an abstract property(A small rigid particle, rotating about an axis), exceeding the speed of light (C).

\chapter{Spherical Harmonics}

Space-Inversion, or Parity\index{Parity} transformation property of the wave-function is:  \\
$Y^m_\ell (\theta, \phi)$ $\rightarrow$ $\boxed{Y^m_\ell(\pi - \theta, \phi + \pi) = (-1)^{\ell}$ $Y^m_\ell$ $(\theta, \phi)}$, $\boxed{\ell = \nu_1 + \nu_2 + \nu_3}$, which is a homogeneous polynomial of rank $\ell$ in the unit vector $\vec{x}$.  \\
Since $\vec{x}$ = - $\vec{x}$, a spherical harmonic change of sign factor,  \\
$\implies$
$\left\{
\begin{array}{ll}
   \theta^{'} = \pi - \theta  \\
   \phi^{'}   = \phi + \pi  \\
\end{array}
\right.$


\chapter{Scattering}



\chapter{Transitions}

Consider Dyson Series, up to the $2_{nd}$ term. The probability of a perturbed wave-function from its' initial state(i), to the final state(f), is given by:
\begin{equation}
P_{fi} = \left| - \frac{i}{\hbar} \int_{0}^t \langle \psi_f \mid \hat V (t^{\prime}) \mid \psi_i \rangle e^{i \omega_{fi} (t^{\prime})} dt^{\prime} \right|^2
\end{equation}
Where $\hat V(t)$, is a time-dependent operator, and $\omega_{fi}$ = $\frac{E_f -  E_i}{\hbar}$ = $\frac{i}{\hbar}$ $\left[\langle \psi_f \mid \hat H_0 \mid \psi_f \rangle - \langle \psi_i \mid \hat H_0 \mid \psi_i \rangle \right]$ = Transition frequency. As a function of $\omega_{fi}$, the transition probability has a diffraction pattern. See Figure~\ref{fig:Interference-pattern-intensity-recorded-by-the-3000-pixel-line-camera.pdf}.  \\
The transition probability of finding the system in a state $\mid \psi_f \rangle$, of Energy ($E_f$), is \emph{maximum only when} $E_i$ = $E_f$, or $\omega_{fi} \simeq 0$. The central peak becomes narrower, and higher as time increases(A property of the Dirac Delta Function)See Equation~\ref{eq:Dirac Delta Function}, hence yields a transition probability of $\mathbf{\textbf{unit time}}$, and guarantees conservation of Energy, in the limit (t) $\rightarrow$ $\infty$. Note that $\iff$ $E_f$ $\neq$ $E_i$ = $0$. A constant \emph{Perturbation}\index{Perturbation} (time-independent), neither removes, nor adds Energy to a system, it causes Energy conservation transitions.

\begin{figure}
\begin{center}
\includegraphics[width=10cm]{Interference-pattern-intensity-recorded-by-the-3000-pixel-line-camera.pdf}
\end{center}
\caption{(A) Interference-pattern intensity recorded by the 3000-pixel line camera, averaged over 10,000 camera frames. (B) Log of spatial spectral power with double-slit power peaking around wave-number 45 (equivalent to a wavelength of about 69 pixels) and single-slit power at 1. The peak around wave-number 90 is a harmonic of the double-slit frequency.}
\label{fig:Interference-pattern-intensity-recorded-by-the-3000-pixel-line-camera.pdf}
\end{figure}

\paragraph{Harmonic [Periodic] Perturbation}
\noindent

The time between the moments of turning the perturbation on or off, is given by:
\begin{equation}
\hat V (t) = \hat v e^{i \omega t} + \hat v^{\dagger} e^{ - i \omega t}
\end{equation}

We have transition probability peaks, when the frequency of the system is $\pm \omega$ $\implies$ RESONANCE. So,  \\
$\boxed{p_{fi} (t)$ = $\frac{2 \pi t}{\hbar}$ $\Big| \langle \psi_{f} \mid \hat{v} \mid \psi_{i} \rangle \Big|^2$ $\delta(E_f - E_i + \hbar \omega)}$
and,
$\boxed{p_{fi} (t)$ = $\frac{2 \pi t}{\hbar}$ $\Big| \langle \psi_f \mid \hat{v}^{\dagger} \mid \psi_{i} \rangle \Big|^2$ $\delta(E_f - E_i - \hbar \omega)}$.  \\
The \emph{Physical} meaning of the above equations $\implies$ that when the system is excited through perturbation, and if $E_f$ < $E_i$, there's photon energy release $\hbar \omega$, to the potential $\hat V (t)$. This process is called $\mathbf{\textbf{STIMULATED EMISSION}}$. When $E_f$ > $E_i$, we have $\mathbf{\textbf{ABSORPTION}}$. The objective of constant harmonic perturbation, is the result of the transmission rate; The relation of the rate of absorption, and emission is known as the $\emph{condition of Detailed-Balancing}$\index{Detailed-Balancing}.

\paragraph{Rapid Approximations}
\noindent

As $e^{i \omega_{fi} t}$ = $\frac{1}{i \omega_{fi}}$ $\frac{\partial e^{i \omega_{fi} t}}{\partial t}$ $\therefore$ \\
-$\frac{i}{\hbar}$ $\int_0^{\tau}$ $\langle \psi_f \mid \hat V \mid \psi_i \rangle$ $e^{i \omega_{fi} t}$ dt = -$\frac{i}{\hbar \omega_{fi}}$ $\int_0^{\tau}$ $\langle \psi_f \mid \hat V \mid \psi_i \rangle$ $\frac{\partial}{\partial t}$ $e^{i \omega_{fi} t}$ dt  \\
The approximation is valid only when the time change in the Energy of perturbation during one period of oscillation. If the rate of change ($\hat V$), and hence ($\hat H(t)$) is very small, then $\frac{\partial}{\partial t}$ $\langle \psi_f \mid \hat V \mid \psi_i \rangle$ $\rightarrow$ 0, $\implies$ $p_{fi}$ $\rightarrow$ 0, hence no \emph{transition}. Approximation is applicable only for transitions between discrete states.
So,
\begin{equation}
p_{fi} (t) = \frac{\Big| \langle \psi_f \mid \hat V (t) \mid \psi_i \rangle \Big|^2}{\hbar^2 \omega_{(fi)^2}}
\end{equation}


\chapter{Related Theorems}
\paragraph{Blochs' Theorem}\index{Blochs' Theorem}
\noindent

For periodic potentials, V(x + a) = V(x); So, $\Psi(x + a)$ = $e^{i K a} \Psi(x)$, for some constant K. \\

K = $\frac{2 \pi n}{a}$, n = $0, \pm 1, \pm 2, ....$  \\

$K^2$ = $\frac{2 m E}{\hbar^2}$  \\

$\therefore$ $\Psi(x)$ = $e^{i K a}$ [$A \sin(K(x + a))$ + $B(\cos(K(x + a))$]

\paragraph{Wicks' Theorem}
\paragraph{Wignar-Eckart Theorem}
\paragraph{virial Theorem} See Eq. \eqref{eq:virial Theorem}
\paragraph{Gram-Schmidt Theorem}
\paragraph{Optical Theorem}
\paragraph{Parsevals' Theorem}
\paragraph{Ritz Theorem}
\paragraph{Steiners' Parallel-Axis Theorem}
\paragraph{Larmors' Theorem}

\chapter{Time-Dependent Problems}

There's no possible state with E = 0, as illustrated by Figure \ref{fig:Ground state search.png}; Thus, it is convenient
to adopt periodic boundary conditions, which require that the wave-function is unaffected by a translation of any of the three Cartesian coordinates.

\chapter{Allowed Energies In Materials, Considering Temperature}
\section{Properties of Helium}

In \url{http:\\pages.cs.wisc.edu/kovar/hall.html},

\chapter{Quanglement}
\section{Entanglement}
There is a need to describe systems that extend over macroscopic distances in ways that are inconsistent with Classical ideas. In Quantum Mechanics, two observable particles ($x_1$), and ($x_2$) commute by Eq. \eqref{eq:Abelian/Commutatitive}. How can measurement of particle $(x_2$)s' Momentum, interfere with ($x_1$)s', if they're far apart? The way out of this dilemma within Quantum Mechanics is to suppose that the second measurement, which gives particle 1 a definite Position, does indeed prevent particle 2 from having a definite Momentum, even though the two particles are far apart. The states of the two particles are said to be $\mathbf{entangled}$. When two systems have wave-functions which differ by at most a constant phase factor, they are said to be in the same Quantum state. It is essential that in every measurement, there shall actually be a disturbance arising from the apparatus that destroys all elements of corresponding to observables that do not commute with the measured variable.


\chapter{Index Theorem}

Clifford Algebra, and (K-Theory).








\chapter{Problems}
\paragraph{Canonical Quantization}
\noindent

(1) What is the form of Lagrange Density?  \\

Answer:  \\

$\mathscr{L}$ = i$\psi^{*}$ $\partial \psi$ + b$\nabla \psi^{*}$ $\times$ $\nabla \psi$  \\

(2) Taking $\psi$, and $\psi^{*}$ as independent fields, what is the form for the Euler-Lagrange equations?  \\

Answer:  \\

For $\psi$, \quad $\frac{\partial \mathscr{L}}{\partial \psi^{*}}$ - $\partial_\mu \frac{\partial \mathscr{L}}{\partial(\partial_\mu \psi^{*})}$ = $i \partial_0 \psi$ - $0$ - $\partial_{i} (b \partial_{i} \psi)$ = $i \dot{\psi}$ - b$\nabla^{2} \psi$ = $0$ \\

For $\psi^{*}$, \quad $\frac{\partial \mathscr{L}}{\partial \psi}$ - $\partial_\mu \frac{\partial \mathscr{L}}{\partial(\partial_\mu \psi)}$ = $0$ - $\partial_0(i \psi^{*})$ - $\partial_{i} (b \partial_{i} \psi^{*})$ = $i \dot{\psi^{*}}$ - b$\nabla^{2} \psi^{*}$ = $0$ \\

\chapter{Appendix}


\begin{remark}


This appendix covers some of the  Mathematics used in  exposition of operators, and linear equations. Here you will find the special functions that are relevant to our current application. These require algebraic techniques.  \\

Calculus is also necessary to understand much of this material, for it is always used in advanced topics.  \\

A table on the history of Quantum Field Theory is listed.  \\

\end{remark}

\section{Inverse Matrix}

This ($2 \times 2$) matrix $\implies$ $\underbrace{\begin{pmatrix}a \quad b \cr c \quad d \end{pmatrix}}_{A}$ $\times$ $\underbrace{\begin{pmatrix}x_{1} \quad x_{2} \cr y_{1} \quad y_{2} \end{pmatrix}}_{A^{-1}}$ = $\underbrace{\begin{pmatrix}1 \quad 0 \cr 0 \quad 1 \end{pmatrix}}_{I}$, where we need to find the matrix (${A^{-1}}$).  \\

(${AB}^{-1}$) = (${B}^{-1}{A}^{-1}$)  \\

($({A}^{T})^{-1}$) = (${A)^{-1})^T}$, so, let's create a system of 4 equations, with 4 unknowns:  \\

(1) $ax_1$ + $by_1$ = $1$  \\

$
\left\{
\begin{array}{ll}
  ax_2 + by_2 = 0 \quad  \text{(2) (homogeneous equation)} \\
  cx_1 + dy_1 = 0 \quad  \text{(3) (homogeneous equation)} \\
\end{array}
\right.$  \\

(4) $cx_2$ + $dy_2$ = $1$  \\

We can solve $y_1$, and $y_2$ in Equations (2), and (3). $y_1$ = $\frac{-c}{d} {x_1}$; $y_2$ = $\frac{-a}{b} {x_2}$.  \\

Substituting for $y_1$ in Equation (1), we have: $ax_1$ - $\frac{-bc}{d} {x_1}$ = $1$. Following Algebra, the equation results to:  \\

$x_1$ = $\frac{d}{ad - bc}$.  \\

$y_1$ = $\frac{-c}{ad - bc}$  \\

Using the $4^{th}$ Equation, we have:  \\

$cx_2$ - $\frac{ad}{b} {x_2}$ = $1$ = $\frac{b}{bc - ad}$ = $\frac{-b}{ad - bc}$, and \ldots  \\

$y_2$ = $\frac{a}{ad - bc}$

So, $\boxed{({A}^{-1}) = \frac{1}{ad - bc} \begin{pmatrix}d \quad -b \cr -c \quad a \end{pmatrix}}$, hence the formula for the \textbf{Inverse Matrix}.  \\

The $\det{A}$ = $ad - bc$. \\

If the $\det{A}$ = $0$, then $A^{-1}$ DOES NOT EXIST! The matrix \emph{isn't} invertible.

\subsection{Finding Eigenvalues, and Eigenvectors in a Matrix}

This is essential for the commutation relations of Poincar$\acute{e}$ Generators, Lie Group Algebras, in Quantum Field Theory.  \\

Let the matrix (B) =  $\begin{pmatrix}0 \quad 1 \cr 0 \quad 0 \end{pmatrix}$, and [B(x) = $\lambda$(x)], is an Eigenvalue problem.  \\

Let (B - $\lambda$I = 0); So, the det(B -$\lambda$I) = $\begin{pmatrix} - \lambda \quad 0 \cr 0 \quad - \lambda \end{pmatrix}$ = $\lambda^2$ = $0$; Making it the characteristic polynomial equation. This results to, $\lambda$ = $0$.  \\

For the Eigenvector, [B(x) = 0]. Its' matrix becomes \ldots. $\begin{pmatrix}0 \quad 1 \cr 0 \quad 0 \end{pmatrix}$ $\times$ $\begin{pmatrix}v_1 \cr v_2 \end{pmatrix}$ = $\begin{pmatrix}0 \cr 0 \end{pmatrix}$; Here $v_2$ = $0$. For $v_1$ = $1$, makes the Eigenvector (x) - $\begin{pmatrix}1 \cr 0 \end{pmatrix}$, making it a 2 by 2 matrix, with single Eigenvalue, and Eigenvector. \\

Consider a skew-symmetric matrix C = $\begin{pmatrix}0 \quad -1 \cr 1 \quad 0 \end{pmatrix}$; ($C^T$ = $-C$)  \\

Its' det(C - $\lambda$I) = $\begin{pmatrix} - \lambda \quad -1 \cr 1 \quad - \lambda \end{pmatrix}$ = $\lambda^2 + 1$ = $0$; So, $\boxed{\lambda^2 = -1}$, which is a \textbf{lateral unit, or imaginary number}, since (i) = $\sqrt{-1}$.  \\

$\therefore$ $\lambda$ has two \emph{roots}; $\lambda_1$ = $(i)$, and $\lambda_2$ = $ - (i)$. Thus, the matrix has \emph{complex Eigenvalues}.  \\

"The shortest path between two truths in $\mathbb{R}$ domain, passes through the $\mathbb{C}$ domain. --- Jacques Hadamard"

\subsection{Matrix Diagonalization}

An [m by n] matrix, with $\mathbb{R}$ Eigenvalues, and \textbf{linear independent} Eigenvectors, is \index{symmetric}.  \\

If [A S = S $\Lambda$], [A = S $\Lambda$ $S^{-1}$], and hence [$\Lambda$ = $S^{-1}$ A S], then $\Lambda$, is the diagonal matrix.  \\

Let $\lambda_1$ = $a + b$. Then the Eigenvector is normalized as: $x_1$ = $\frac{1}{\sqrt(2)}$ $\begin{pmatrix}1 \cr 1 \end{pmatrix}$; \\

$\lambda_2$ = $a - b$, $x_2$ = $\frac{1}{srt{2}}$ $\begin{pmatrix}1 \cr -1 \end{pmatrix}$. So, the matrix S =$\frac{1}{\sqrt{2}}$ $\begin{pmatrix}1\quad 1 \cr 1 \quad -1 \end{pmatrix}$, and $\Lambda$ = $S^{-1}$ A S, hence \textbf{orthonormal}. \\

$S^{-1}$ = $S^{T}$ = S, making this matrix \emph{unique}. \\

$\therefore$ $\boxed{\begin{pmatrix}a + b \quad 0 \cr0 \quad a - b \end{pmatrix} = \frac{1}{2} \begin{pmatrix}1 \quad \quad 1 \cr 1 \quad -1 \end{pmatrix} \begin{pmatrix}a \quad b \cr b \quad a \end{pmatrix} \begin{pmatrix}1 \quad \quad 1 \cr 1 \quad -1 \end{pmatrix}}$   \\

Only the diagonal matrix elements, in this case, are different from zero. In particular, in the usual representation, with the wave-functions of the stationary states as the functions $\Psi_n$, the Energy matrix is diagonal; So are the matrices of all other Physical quantities having definite values, in the stationary states.   \\

\section{Rules for Continuous Functions[Apply Limit Laws]}

Product Rule $\rightarrow$ Let f(x) and g(x) be differential functions. Then,  \\
\begin{equation}
\boxed{\frac{d}{dx}f(x)g(x) = \frac{d}{dx} [ f(x) \cdot g(x) ] + \frac{d}{dx} [ g(x) \cdot f(x) ]}
\end{equation}

That is, if [ f(x) = f(x) g(x) ], then [ f(x) = $f^{'}(x) g(x) + g^{'}(x) f(x)$ ].  \\

This means that the derivative of a product of two functions is the derivative of the first function times the second function, plus, the derivative of the second function times the first function.


\begin{enumerate}
\item \textbf{Quotient Rule}
\item \textbf{Power Rule}
\end{enumerate}

Chain Rule $\rightarrow$ Let f and g be functions. $\forall$ x in the domain of g, for which g is differentiable at x, and f is differentiable at g(x), the derivative of the composite function h(x) = f(g(x) is given by
\begin{equation}
h'(x) = f'(g(x) g'(x)
\end{equation}

Alternatively, if (y) is a function of (u), and (u) is a function of (x), then $\frac{dy}{dx}$ = $\frac{dy}{du}$ $\frac{du}{dx}$ e.g $\frac{d}{dx}$ $(2x^2 - 5)^5$ = $20(x) (2x^2 - 5)^4$  \\

Contour Integrals $\rightarrow$ If the function f(z) has isolated singularities (poles) inside a closed contour $(\gamma)$, the integral around that contour is equal to
2$\pi$i times the sum of the residues at those singularities; $\oint _{\gamma} f(z) dz$ = 2 $\pi$i $\sum \text{Res}$ $(f,z_{k})$  \\

\section{Hermite Polynomials}

A polynomial multiplying $e^{\frac{(-u)^2}{2}}$ could be a solution of a differential equation. The solution could take the form $\psi$(u) = h(u)$e^{\frac{(-u)^2}{{2}}}$, where h(u) is a polynomial in (u). \\

An Energy parameter $\epsilon$, takes special values. Any function $\phi$(u) can be written as some other function, times $e^{\frac{(-u)^(2)}{2}}$; if h(u), then we've found $\phi$(u), leading to a second-order linear differential equation for h(u):The polynomial solutions $h_n(u)$ is normalized with a specific choice of the
leading coefficient, and satisfy a diffrential equation. The first few Hermite polynomials are:  \\

$\bullet$ $H_0(u)$ = $1$  \\

$\bullet$ $H_1(u)$ = $2u$  \\

$\bullet$ $H_2(u)$ = $4u^2 - 2$  \\

$\bullet$ $H_3(u)$ = $8u^3 - 12u$  \\

There exists a generating function for Hermite polynomials $H_n(u)$. It is a function of (u), and a formal expansion parameter (z):  \\

$e^{-z^2 + 2zu}$ = $\sum_{n = 0}^\infty$ $\frac{z^n}{n!}$ $H_n(u)$.  \\

Taking a normalization constant $N_n$, and using $H_0(u)$ = $1$, we get:  \\

$\phi_0$(x) = $N_0$ $e^- {\frac{(x)^2}{2L^{2}_0}}$, with $N^{2}_0$ = $\sqrt{\frac{(m) \omega}{\pi \hbar}}$ = $\frac{1}{\sqrt{\pi}L_0}$. \\

The ground-state is a pure Gaussian in (x), and the excited states are \emph{polynomials}, times that same Gaussian.  \\

\begin{tikzpicture}
\begin{axis}[
    title={First Four Hermite Polynomials},
    xlabel={$x$},
    ylabel={$H_n(x)$},
    xmin=-2, xmax=2,
    ymin=-15, ymax=15,
    samples=200,
    domain=-2:2,
    grid=major,
    legend pos=south east
]

% H_0(x) = 1
\addplot[thick, blue] {1};
\addlegendentry{$H_0$}

% H_1(x) = 2x
\addplot[thick, red] {2*x};
\addlegendentry{$H_1$}

% H_2(x) = 4x^2 - 2
\addplot[thick, green!70!black] {4*x^2 - 2};
\addlegendentry{$H_2$}

% H_3(x) = 8x^3 - 12x
\addplot[thick, orange] {8*x^3 - 12*x};
\addlegendentry{$H_3$}

\end{axis}
\end{tikzpicture}


\subsection{Laguerre Polynomials}

 A set of orthogonal polynomials defined over the interval $[0, \infty]$. In Quantum Mechanics, the Schrödinger equation for the Hydrogen-like atom, is solvable by separation of variables, in spherical coordinates. The radial part of the wave-function is a (generalized) Laguerre polynomial.  \\

 A potential (V) is the Coulomb potential $\frac{-Ze^2}{r}$, where (Z) is the atomic number. The reduced mass $\mu$ is nearly equal to the electron mass $m_e$, so $\mu$ $\approx$ $m_e$ can be made.  \\

The hydrogenic radial wave-function is:  \\

$\phi$(r) = $R_{nl}$(r) = $N_{nl}$ $(\frac{2Zr}{na_0})^l$ $e^{\frac{-Zr}{na_0}}$ $\mathbf{\textbf{L}}^{2l + 1}_{n + l}$ $(\frac{2Zr}{na_0})$, where   \\

$a_0$ = $\frac{\hbar^2}{m_e e^2}$, and $N_{nl}$ = $\left[\frac{4Z^3(n - l - 1)!}{n^4a^3_0[(n + l]!]^3} \right]^\frac{1}{2}$, and \\

$L^{2l + 1}_{n + !}$ is an associated \emph{Laguerre polynomial}. The Quantum Numbers in this case have the values l = 0, 1, \ldots, n - 1, and n = 1, 2, \ldots, $\infty$.  \\

Energy eigenvalues for the Hydrogen Atom are given by: $E_n$ = - $\frac{Z^2}{2n^2}$ $\frac{e^2}{a_0}$.  \\

\href{https://en.wikipedia.org/wiki/Laguerre_polynomials#cite_note-24} for in-depth exposition of Edmond Laguerre polynomials.






















\begin{tikzpicture}
\begin{axis}[
    title={Laguerre Polynomials $L_n(x)$},
    xlabel={$x$},
    ylabel={$L_n(x)$},
    xmin=0, xmax=5,
    ymin=-2, ymax=5,
    grid=major,
    axis lines=middle,
    legend pos=north east,
    domain=0:5,
    samples=100,
    thick
]

% L_0(x) = 1
\addplot[color=blue] {1};
\addlegendentry{$L_0(x)$}

% L_1(x) = -x + 1
\addplot[color=red] {-x + 1};
\addlegendentry{$L_1(x)$}

% L_2(x) = 0.5*(x^2 - 4*x + 2)
\addplot[color=green!60!black] {0.5*(x^2 - 4*x + 2)};
\addlegendentry{$L_2(x)$}

% L_3(x) = (1/6)*(-x^3 + 9*x^2 - 18*x + 6)
\addplot[color=orange] {(1/6)*(-x^3 + 9*x^2 - 18*x + 6)};
\addlegendentry{$L_3(x)$}

\end{axis}
\end{tikzpicture}


\begin{table}[htbp]
\centering
\caption{Table of First Few Laguerre and Associated Laguerre Polynomials}
\label{tab:laguerre_polynomials}
\begin{tabular}{ccl}
\toprule
$n$ & Standard Laguerre Polynomial $L_n(x)$ & Associated Laguerre Polynomial $L_n^\alpha(x)$ \\
\midrule
$0$ & $1$ & $1$ \\
\addlinespace
$1$ & $-x + 1$ & $-x + \alpha + 1$ \\
\addlinespace
$2$ & $\frac{1}{2}(x^2 - 4x + 2)$ & $\frac{1}{2}\left[x^2 - 2(\alpha + 2)x + (\alpha + 1)(\alpha + 2)\right]$ \\
\addlinespace
$3$ & $\frac{1}{6}(-x^3 + 9x^2 - 18x + 6)$ & $\frac{1}{6}\left[-x^3 + 3(\alpha + 3)x^2 - 3(\alpha + 2)(\alpha + 3)x + (\alpha + 1)(\alpha + 2)(\alpha + 3)\right]$ \\
\bottomrule
\end{tabular}
\end{table}


\subsection{Legendre Polynomials}















\section{Proof of Fouriers' Theorem}

\section{The Canonical Commutation Relations}

Annihilation Operation: $[\psi (\mathbf{x}, t), \psi (\mathbf{y}, t)]$ = $[\psi^{*} (\mathbf{x}, t). \psi^{*} (\mathbf{y}, t)]$ = $0$  \\

Creation Operation: $[\psi (\mathbf{x}, t), i \psi^{*} (\mathbf{y}, t)]$  = i$\delta^{(3)} (\mathbf{x} - \mathbf{y})$   \\

\section{Charge, Parity, and Time-Reversal}

\paragraph{Charge}

%\begin{align*}
%Q &= \int d^{3}\mathbf{x} \, \psi^{*}(\mathbf{x}, t) \psi(\mathbf{x}, t) \\
 % &= \int \frac{d^{3}\mathbf{p}}{(2 \pi)^3} d^{3}\mathbf{p}' \, a^{\dagger}_\mathbf{p} e^{-i(\mathbf{p} \cdot \mathbf{x} - \omega_\mathbf{p}t)} a_{\mathbf{p}'} %e^{i(\mathbf{p}' \cdot \mathbf{x} - \omega_{\mathbf{p}'}t)} \\
  %&= \int d^{3}\mathbf{p} \, a^{\dagger}_\mathbf{p} a_\mathbf{p}
%\end{align*}


$Q$  \\
= $\int$ $d^{3}$ $\mathbf{\textbf{x}}$ $\psi^{*}$ $(\mathbf{\textbf{x}}, t)$ $\psi$ $(\mathbf{x}, t)$ \\
= $\int$ $\frac{d^{3}}{(2 \pi)^3}$ $d^{3}$ $\mathbf{\textbf{p}}$ $d^{3}$ $\mathbf{\textbf{p}}^{\prime}$ $a^{\dagger}_\mathbf{\textbf{p}}$ $e^{-i(\mathbf{\textbf{p}}
\cdot \mathbf{\textbf{x}} - \omega \mathbf{\textbf{p}}t}$ $a_\mathbf{\textbf{p}}^{\prime}$ $e^{i(\mathbf{\textbf{p}} \cdot \mathbf{\textbf{x}} - \omega \mathbf{\textbf{p}}^{\prime}t)}$ \\
= $\int$ $d^{3}$ $\mathbf{\textbf{p}}$ $a^{\dagger}_\mathbf{\textbf{p}}$ $a_\mathbf{\textbf{p}}$  \\

This is the theory of a set of free, non-relativistic, identical bosons, all with mass (m). Each boson has momentum (p), and Energy [E] = $\omega$ = $\frac{\mid p \mid^2}{2m}$. The conserved charge (Q), is the number of bosons. All of the operators; the Hamiltonian (H), Momentum(P), and Charge(Q), are time-independent.  \\





















\subsection{Proof of the Derivative of the Delta Function}

The proof uses integration by parts. Since the delta function $(\delta(x))$ vanishes at infinity, the \emph{boundary term is zero}.  \\

$\int_{-\infty}^{\infty} \delta'(x) f(x) dx$ = $\left[ \delta(x) f(x) \right]_{-\infty}^{\infty} - \int_{-\infty}^{\infty} \delta(x) f'(x) dx$ = $0 - f'(0)$ = $-f'(0)$  \\

\subsection{Proof of Mathieu Functions}
















\newpage

\section{History Of Quantum Field Theory}

\begin{longtable}{p{0.1\textwidth} p{0.3\textwidth} p{0.6\textwidth}}
\caption{Historical Milestones of Quantum Field Theory} \\
\toprule
\textbf{Year} & \textbf{Event / Discovery} & \textbf{Significance / Description} \\
\midrule
\endfirsthead

\multicolumn{3}{c}%
{\tablename\ \thetable\ -- \textit{Continued from previous page}} \\
\toprule
\textbf{Year} & \textbf{Event / Discovery} & \textbf{Significance / Description} \\
\midrule
\endhead

\midrule
\multicolumn{3}{r}{\textit{Continued on next page}} \\
\endfoot

\bottomrule
\endlastfoot

1900 & Planck's Radiation Law & Max Planck introduces energy quantization, planting the seeds for quantum theory. \\
\midrule
1905 & Einstein's Light Quanta & Albert Einstein proposes the photon hypothesis, establishing wave-particle duality. \\
\midrule
1926 & Heisenberg, Born, and Jordan & Formulation of matrix mechanics, treating fields as infinite sets of harmonic oscillators. \\
\midrule
1927 & Dirac's Quantization of Fields & Paul Dirac publishes a foundational paper on emission/absorption of radiation, formally initiating Quantum Electrodynamics (QED). \\
\midrule
1930s & Divergences and Infinities & Robert Oppenheimer identifies that higher-order perturbative calculations yield infinite quantities (e.g., self-energy of the electron). \\
\midrule
1947--1949 & Renormalization & Breakthroughs by Feynman, Schwinger, Tomonaga, and Dyson. They establish renormalization to systematically remove infinities from QED calculations. \\
\midrule
1950s & Feynman Diagrams & Richard Feynman introduces graphical representations of particle interactions, revolutionizing perturbative calculations. \\
\midrule
1954 & Yang-Mills Theory & Chen-Ning Yang and Robert Mills propose non-Abelian gauge theories, laying the mathematical foundation for the strong and weak forces. \\
\midrule
1964 & Quark Model & Murray Gell-Mann and George Zweig independently propose the quark model to explain the classification of hadrons. \\
\midrule
1967--1968 & Electroweak Unification & Glashow, Salam, and Weinberg formulate the electroweak theory, unifying the electromagnetic and weak forces. \\
\midrule
1970s & Asymptotic Freedom & Gross, Politzer, and Wilczek discover asymptotic freedom in Quantum Chromodynamics (QCD), explaining why quarks behave as free particles at high energies. \\
\midrule
1973 & Standard Model Established & The theoretical framework of the Standard Model of particle physics takes its modern shape, organizing all known fundamental particles. \\
\midrule
1980s & Lattice Gauge Theory & Kenneth Wilson establishes non-perturbative approaches using lattice gauge theory, allowing computer simulations of QCD. \\
\midrule
1998 & Neutrino Oscillations & Evidence from Super-Kamiokande shows that neutrinos have mass, requiring minor extensions to the Standard Model. \\
\midrule
2012 & Higgs Boson Discovery & CERN physicists announce the discovery of the Higgs boson, confirming the mechanism of electroweak symmetry breaking. \\
\bottomrule
\end{longtable}

\newpage

\section{Hamilton Equations Of Motion}





\section{Wheeler-De Witt Equation}






\section{List of experiments}

\begin{enumerate}

\item Franck-Hertz
\item Davisson-Germer  
\item Photo-electric effect  
\item Youngs double-slit 
\item Stern-Gerlach  
\item EPR   
\item Badurek-Rauch-Tuppinger  
\item SQUID(A Superconducting Quantum Interference Device (SQUID) experiment, measures extremely weak magnetic fields by exploiting Josephson junctions, and magnetic flux quantization within a superconducting loop, cooled by liquid nitrogen at $77^\circ$K or $-196.1^\circ$C. 
\item Delayed-choice

\end{enumerate} 

\begin{thebibliography}{99}

\bibitem{Mermin}
N. D. Mermin,  Phys. Rev. Lett. 65, 1838 (1990).

\bibitem{Zweibach}
B. Zweibach, Mastering Quantum Mechanics, (2022).

\bibitem{Bates}
D.R.Bates, Vol I-III(Elements(I), Aggregate of Particles(II), Radiation and High Energy Physics(III)), (1961)

\bibitem{Strang}
G.Strang, Introduction To Linear Algebra(5th Edition), (2016)

\bibitem{Landau}
L.D.Landau and E.M.Lifshitz, Quantum Mechanics; Non-relativistic Theory(Volume 3), (1977)

\end{thebibliography}


\chapter{Reference}


\url{https://pubchem.ncbi.nlm.nih.gov/rest/pug/periodictable/CSV?response_type=display} \\

\begin{enumerate}

\item Quantum Theory for Mathematicians - Brian C. Hall  
\item Notes from Sidney Coleman's Physics 253a - Sidney Coleman$^*$  
\item Courant and Hilbert: Methods of Mathematical Physics volumes I and II - R. Courant and D. Hilbert 
\item An introduction to Quantum Optics - Gilbert Grynberg, Alain Aspect, and Claude Fabre  
\item Methods of Modern Mathematical Physics volumes I-IV - Michael Reed, and Barry Simon 
\item Mathematics for Engineers:Numerics - Gerd Baumann 
\item Functional Analysis - K$\dot{o}$saku Yosida  
\item Calculus volumes I-III - Edwin "Jed" Herman and Gilbert Strang  
\item Mathematical foundations of Quantum Mechanics - George W. Mackey 
\item The Feynmann lectures on Physics - R.P.Feynman, R.B.Leighton, and M.Sands  
\item Essential Calculus skills practice workbook with full solutions - Chris McMullen  
\item Applied partial differential equations - Richard Haberman  
\item Physics for Scientists and Engineers - Paul A. Tipler, Gene Mosca  
\item Understanding Quantum Mechanics - Roland Omn$\acute{e}$s  
\item A course of modern analysis - E.T.Whittaker, and G.N.Watson  
\item Handbook of Mathematical Functions; With formulas, Graphs, and Mathematical Tables - Milton Abramowitz, and Irene A. Stegun
\item Zaarur E Schaum's outline of Quantum Mechanics including hundred of solved problems  
\item Pauli and the Spin-Statistics Theorem - Ian Duck, and E.C.G.Sudarshan 
\item Lectures on Quantum Mechanics - Dirac Paul A.M.  
\item Quantum Mechanics; Concepts and Applications - Nouredine Zettili  
\item Riemannian Geometry and Geometric Analsysis - J$\ddot{u}gen$ Jost  
\item Lectures on Quantum Mechanics - Steven Weinberg

\end{enumerate}

\printindex
\label{sec:index}
\end{document}
